\chapter{142只高影响候选全量审计}

下表展示全量候选的当前价格、H1扣非利润、概率值和动作。它不是把142只都升级为投资建议，而是把每一只的当前处置公开化。

长鑫科技（688825）处于IPO待上市状态，发行价8.66元；截至本轮数据截点尚无二级市场当前价，因此目标价和空间标记为不适用，不是数据漏填。

\scriptsize
\setlength{\tabcolsep}{1.5pt}
\begin{longtable}{L{1.05cm}L{1.55cm}L{1.55cm}R{0.85cm}R{0.85cm}R{0.85cm}R{0.85cm}R{0.65cm}L{2.65cm}}
\toprule
\textbf{代码} & \textbf{公司} & \textbf{行业} & \textbf{现价} & \textbf{H1扣非} & \textbf{目标} & \textbf{空间} & \textbf{证据} & \textbf{动作} \\
\midrule
\endfirsthead
\toprule
\textbf{代码} & \textbf{公司} & \textbf{行业} & \textbf{现价} & \textbf{H1扣非} & \textbf{目标} & \textbf{空间} & \textbf{证据} & \textbf{动作} \\
\midrule
\endhead
000042 & 中洲控股 & 房地产 & 8.02元 & 10.6 & 9元 & 6.5\% & 中 & 市场观察 \\
600150 & 中国船舶 & 国防军工 & 34.31元 & 99.0 & 83元 & 142.9\% & 中高 & 高空间验证 \\
002379 & 宏桥控股 & 有色金属 & 18.66元 & 154.0 & 26元 & 38.3\% & 高 & 高空间验证 \\
601600 & 中国铝业 & 有色金属 & 8.98元 & 115.0 & 15元 & 62.8\% & 中高 & 高空间验证 \\
600346 & 恒力石化 & 石油石化 & 15.24元 & 53.3 & 14元 & -6.8\% & 中高 & 估值观察 \\
000858 & 五粮液 & 食品饮料 & 76.40元 & 87.0 & 75元 & -1.5\% & 中高 & 估值观察 \\
002460 & 赣锋锂业 & 有色金属 & 51.50元 & 36.0 & 54元 & 4.2\% & 中高 & 市场观察 \\
002466 & 天齐锂业 & 有色金属 & 47.43元 & 35.0 & 44元 & -6.4\% & 中高 & 估值观察 \\
601688 & 华泰证券 & 非银金融 & 20.63元 & 115.2 & 23元 & 9.4\% & 中高 & 市场观察 \\
601336 & 新华保险 & 非银金融 & 65.55元 & 222.7 & 78元 & 19.7\% & 中高 & 回撤验证 \\
600037 & 歌华有线 & 传媒 & 6.80元 & 2.1 & 6元 & -17.4\% & 中 & 高估值 \\
002738 & 中矿资源 & 有色金属 & 49.33元 & 11.0 & 39元 & -20.4\% & 中高 & 高估值 \\
601069 & 西部黄金 & 有色金属 & 22.96元 & 5.3 & 13元 & -44.5\% & 中 & 高估值 \\
000567 & 海德股份 & 非银金融 & 6.17元 & 14.4 & 7元 & 20.8\% & 中 & 回撤验证 \\
002602 & 世纪华通 & 传媒 & 14.26元 & 43.1 & 17元 & 15.8\% & 中高 & 回撤验证 \\
600026 & 中远海能 & 交通运输 & 15.83元 & 44.0 & 21元 & 34.8\% & 中高 & 高空间验证 \\
002048 & 宁波华翔 & 汽车 & 21.43元 & 5.3 & 28元 & 33.0\% & 中 & 高空间验证 \\
600688 & 上海石化 & 石油石化 & 2.84元 & 3.1 & 3元 & -8.1\% & 中 & 估值观察 \\
600095 & 湘财股份 & 非银金融 & 7.81元 & 5.2 & 4元 & -46.2\% & 中 & 高估值 \\
000426 & 兴业银锡 & 有色金属 & 30.56元 & 18.1 & 24元 & -22.9\% & 中高 & 高估值 \\
600489 & 中金黄金 & 有色金属 & 20.12元 & 43.0 & 21元 & 5.2\% & 中高 & 市场观察 \\
000737 & 北方铜业 & 有色金属 & 12.75元 & 13.5 & 15元 & 20.8\% & 中低 & 回撤验证 \\
002756 & 永兴材料 & 有色金属 & 44.90元 & 10.0 & 42元 & -6.2\% & 中高 & 估值观察 \\
300014 & 亿纬锂能 & 电力设备 & 52.28元 & 25.2 & 59元 & 12.1\% & 中高 & 市场观察 \\
601061 & 中信金属 & 商贸零售 & 10.82元 & 26.4 & 17元 & 57.7\% & 中 & 高空间验证 \\
600739 & 辽宁成大 & 医药生物 & 10.95元 & 15.8 & 43元 & 290.4\% & 中 & 高空间验证 \\
601360 & 三六零 & 计算机 & 9.00元 & 1.6 & 13元 & 45.6\% & 中 & 高空间验证 \\
600918 & 中泰证券 & 非银金融 & 5.71元 & 17.0 & 5元 & -6.3\% & 中 & 估值观察 \\
000975 & 山金国际 & 有色金属 & 19.34元 & 23.4 & 20元 & 5.0\% & 中高 & 市场观察 \\
000623 & 吉林敖东 & 医药生物 & 18.66元 & 22.2 & 78元 & 317.0\% & 中 & 高空间验证 \\
600292 & 电投水电 & 公用事业 & 12.50元 & 11.8 & 7元 & -47.5\% & 中 & 高估值 \\
000987 & 越秀资本 & 非银金融 & 7.77元 & 23.3 & 8元 & 2.5\% & 中 & 市场观察 \\
000155 & 川能动力 & 公用事业 & 11.09元 & 6.8 & 10元 & -9.5\% & 高 & 估值观察 \\
601377 & 兴业证券 & 非银金融 & 6.21元 & 20.9 & 6元 & -0.8\% & 中高 & 估值观察 \\
600120 & 浙江东方 & 非银金融 & 4.91元 & 8.4 & 5元 & 2.2\% & 中低 & 市场观察 \\
000703 & 恒逸石化 & 石油石化 & 14.48元 & 57.2 & 23元 & 61.7\% & 高 & 高空间验证 \\
000301 & 东方盛虹 & 石油石化 & 11.61元 & 44.2 & 11元 & -3.7\% & 中高 & 估值观察 \\
001339 & 智微智能 & 计算机 & 89.21元 & 3.7 & 65元 & -26.7\% & 中高 & 高估值 \\
002558 & 巨人网络 & 传媒 & 29.56元 & 23.0 & 34元 & 15.0\% & 中高 & 市场观察 \\
301308 & 江波龙 & 电子 & 488.00元 & 97.5 & 1069元 & 119.2\% & 中高 & 高空间验证 \\
601628 & 中国人寿 & 非银金融 & 40.42元 & 1332.8 & 49元 & 20.2\% & 中高 & 回撤验证 \\
002493 & 荣盛石化 & 石油石化 & 11.10元 & 53.0 & 9元 & -20.0\% & 中高 & 高估值 \\
001309 & 德明利 & 电子 & 662.40元 & 60.5 & 482元 & -27.3\% & 中高 & 高估值 \\
601233 & 桐昆股份 & 石油石化 & 20.72元 & 40.8 & 29元 & 38.0\% & 中高 & 高空间验证 \\
603986 & 兆易创新 & 电子 & 571.79元 & 48.5 & 351元 & -38.6\% & 中高 & 高估值 \\
601225 & 陕西煤业 & 煤炭 & 24.61元 & 97.5 & 18元 & -28.4\% & 中高 & 高估值 \\
601138 & 工业富联 & 电子 & 63.62元 & 232.0 & 69元 & 8.1\% & 中高 & 市场观察 \\
000807 & 云铝股份 & 有色金属 & 26.10元 & 75.5 & 45元 & 72.5\% & 中高 & 高空间验证 \\
603993 & 洛阳钼业 & 有色金属 & 17.91元 & 155.0 & 16元 & -8.2\% & 中高 & 估值观察 \\
601872 & 招商轮船 & 交通运输 & 15.26元 & 68.9 & 22元 & 41.4\% & 中高 & 高空间验证 \\
600999 & 招商证券 & 非银金融 & 20.90元 & 105.0 & 19元 & -11.4\% & 中高 & 估值观察 \\
601899 & 紫金矿业 & 有色金属 & 28.83元 & 379.0 & 35元 & 21.9\% & 中高 & 回撤验证 \\
600309 & 万华化学 & 基础化工 & 70.09元 & 99.0 & 80元 & 14.3\% & 中高 & 市场观察 \\
600989 & 宝丰能源 & 基础化工 & 22.07元 & 91.5 & 32元 & 43.5\% & 高 & 高空间验证 \\
300475 & 香农芯创 & 电子 & 221.20元 & 36.4 & 373元 & 68.4\% & 中 & 高空间验证 \\
000776 & 广发证券 & 非银金融 & 22.78元 & 116.9 & 23元 & 0.3\% & 中高 & 市场观察 \\
600030 & 中信证券 & 非银金融 & 28.60元 & 236.1 & 25元 & -12.7\% & 中高 & 估值观察 \\
601088 & 中国神华 & 煤炭 & 43.69元 & 271.0 & 22元 & -50.4\% & 中 & 高估值 \\
600362 & 江西铜业 & 有色金属 & 42.28元 & 81.0 & 51元 & 20.1\% & 中 & 回撤验证 \\
002709 & 天赐材料 & 电力设备 & 37.55元 & 28.0 & 60元 & 61.0\% & 高 & 高空间验证 \\
002414 & 高德红外 & 国防军工 & 13.32元 & 13.2 & 21元 & 56.1\% & 中 & 高空间验证 \\
300390 & 天华新能 & 电力设备 & 60.01元 & 22.6 & 77元 & 28.8\% & 中高 & 回撤验证 \\
002648 & 卫星化学 & 基础化工 & 23.43元 & 61.9 & 44元 & 85.7\% & 中高 & 高空间验证 \\
000792 & 盐湖股份 & 基础化工 & 25.27元 & 61.5 & 29元 & 16.5\% & 中 & 回撤验证 \\
601869 & 长飞光纤 & 通信 & 425.45元 & 23.0 & 148元 & -65.3\% & 中 & 高估值 \\
002432 & 九安医疗 & 医药生物 & 66.99元 & 31.5 & 284元 & 323.4\% & 中 & 高空间验证 \\
000657 & 中钨高新 & 有色金属 & 69.80元 & 20.6 & 23元 & -67.0\% & 中高 & 高估值 \\
601995 & 中金公司 & 非银金融 & 38.34元 & 78.1 & 26元 & -31.7\% & 中高 & 高估值 \\
002497 & 雅化集团 & 基础化工 & 18.07元 & 12.2 & 32元 & 76.0\% & 高 & 高空间验证 \\
000933 & 神火股份 & 有色金属 & 25.51元 & 48.1 & 44元 & 73.3\% & 中 & 高空间验证 \\
300223 & 北京君正 & 电子 & 190.18元 & 11.5 & 121元 & -36.4\% & 中高 & 高估值 \\
002192 & 融捷股份 & 有色金属 & 66.42元 & 10.0 & 84元 & 26.2\% & 中 & 回撤验证 \\
601066 & 中信建投 & 非银金融 & 27.82元 & 77.9 & 15元 & -46.5\% & 中高 & 高估值 \\
002812 & 恩捷股份 & 电力设备 & 54.96元 & 8.5 & 85元 & 55.2\% & 高 & 高空间验证 \\
002636 & 金安国纪 & 电子 & 94.79元 & 7.8 & 54元 & -42.5\% & 中 & 高估值 \\
002378 & 章源钨业 & 有色金属 & 25.84元 & 6.8 & 12元 & -52.4\% & 中 & 高估值 \\
601211 & 国泰海通 & 非银金融 & 18.58元 & 195.0 & 20元 & 10.1\% & 中高 & 市场观察 \\
600884 & 杉杉股份 & 电力设备 & 12.14元 & 7.8 & 15元 & 21.5\% & 中 & 回撤验证 \\
002842 & 翔鹭钨业 & 有色金属 & 35.96元 & 5.3 & 35元 & -1.3\% & 中 & 估值观察 \\
002407 & 多氟多 & 基础化工 & 33.88元 & 4.9 & 10元 & -69.3\% & 中 & 高估值 \\
002653 & 海思科 & 医药生物 & 74.13元 & 7.3 & 32元 & -57.0\% & 中高 & 高估值 \\
603629 & 利通电子 & 电子 & 112.09元 & 6.1 & 85元 & -24.4\% & 中 & 高估值 \\
600961 & 株冶集团 & 有色金属 & 24.76元 & 19.1 & 38元 & 51.8\% & 中高 & 高空间验证 \\
002128 & 电投能源 & 煤炭 & 26.51元 & 40.5 & 24元 & -9.5\% & 中高 & 估值观察 \\
601168 & 西部矿业 & 有色金属 & 33.18元 & 41.9 & 38元 & 15.1\% & 高 & 回撤验证 \\
600188 & 兖矿能源 & 煤炭 & 20.23元 & 45.0 & 10元 & -50.8\% & 中高 & 高估值 \\
603162 & 海通发展 & 交通运输 & 11.24元 & 5.2 & 9元 & -24.1\% & 中高 & 高估值 \\
002240 & 盛新锂能 & 有色金属 & 31.11元 & 14.0 & 38元 & 21.8\% & 中高 & 回撤验证 \\
002440 & 闰土股份 & 基础化工 & 11.29元 & 4.4 & 10元 & -12.1\% & 中 & 估值观察 \\
002532 & 天山铝业 & 有色金属 & 12.75元 & 40.9 & 20元 & 60.5\% & 中高 & 高空间验证 \\
301377 & 鼎泰高科 & 机械设备 & 467.88元 & 6.6 & 60元 & -87.3\% & 中高 & 高估值 \\
688257 & 新锐股份 & 机械设备 & 80.06元 & 5.9 & 64元 & -20.7\% & 中高 & 高估值 \\
000408 & 藏格矿业 & 基础化工 & 71.50元 & 37.1 & 65元 & -8.9\% & 中高 & 估值观察 \\
600601 & 方正科技 & 电子 & 11.71元 & 5.5 & 7元 & -43.6\% & 中 & 高估值 \\
002475 & 立讯精密 & 电子 & 60.47元 & 64.9 & 43元 & -29.4\% & 中 & 高估值 \\
600595 & 中孚实业 & 有色金属 & 6.46元 & 19.2 & 11元 & 66.6\% & 中高 & 高空间验证 \\
600487 & 亨通光电 & 通信 & 71.05元 & 33.3 & 69元 & -3.5\% & 中高 & 估值观察 \\
600183 & 生益科技 & 电子 & 152.43元 & 28.2 & 59元 & -61.3\% & 中高 & 高估值 \\
000630 & 铜陵有色 & 有色金属 & 6.06元 & 28.5 & 6元 & -7.4\% & 中高 & 估值观察 \\
600549 & 厦门钨业 & 有色金属 & 58.08元 & 21.8 & 32元 & -44.7\% & 中高 & 高估值 \\
000100 & TCL科技 & 电子 & 5.02元 & 30.3 & 8元 & 61.2\% & 中高 & 高空间验证 \\
001287 & 中电港 & 电子 & 26.65元 & 5.2 & 41元 & 52.0\% & 中高 & 高空间验证 \\
600111 & 北方稀土 & 有色金属 & 39.56元 & 20.3 & 13元 & -68.0\% & 中高 & 高估值 \\
000725 & 京东方A & 电子 & 6.38元 & 31.9 & 6元 & -13.6\% & 中高 & 估值观察 \\
600176 & 中国巨石 & 建筑材料 & 53.00元 & 29.8 & 27元 & -48.5\% & 中高 & 高估值 \\
002266 & 浙富控股 & 环保 & 5.09元 & 12.5 & 8元 & 52.6\% & 中高 & 高空间验证 \\
000612 & 焦作万方 & 有色金属 & 10.84元 & 12.1 & 21元 & 90.6\% & 中 & 高空间验证 \\
605117 & 德业股份 & 电力设备 & 85.40元 & 25.2 & 98元 & 14.2\% & 中高 & 市场观察 \\
000783 & 长江证券 & 非银金融 & 9.06元 & 31.8 & 9元 & -3.8\% & 中 & 估值观察 \\
600233 & 圆通速递 & 交通运输 & 17.50元 & 31.9 & 21元 & 19.2\% & 中高 & 回撤验证 \\
600392 & 盛和资源 & 有色金属 & 23.86元 & 8.6 & 11元 & -55.6\% & 中 & 高估值 \\
603225 & 新凤鸣 & 基础化工 & 17.97元 & 15.2 & 23元 & 26.5\% & 中高 & 回撤验证 \\
002008 & 大族激光 & 机械设备 & 114.36元 & 14.0 & 52元 & -54.2\% & 中高 & 高估值 \\
000685 & 中山公用 & 环保 & 11.19元 & 14.3 & 29元 & 162.8\% & 中高 & 高空间验证 \\
600909 & 华安证券 & 非银金融 & 9.50元 & 21.6 & 7元 & -29.5\% & 中高 & 高估值 \\
000833 & 粤桂股份 & 综合 & 20.07元 & 6.5 & 27元 & 34.7\% & 中 & 高空间验证 \\
300037 & 新宙邦 & 电力设备 & 67.20元 & 9.8 & 69元 & 2.5\% & 高 & 市场观察 \\
000060 & 中金岭南 & 有色金属 & 6.26元 & 11.1 & 6元 & -8.9\% & 中高 & 估值观察 \\
600522 & 中天科技 & 通信 & 40.70元 & 22.3 & 33元 & -18.3\% & 中 & 高估值 \\
002468 & 申通快递 & 交通运输 & 15.22元 & 10.1 & 14元 & -5.4\% & 中高 & 估值观察 \\
600267 & 海正药业 & 医药生物 & 11.98元 & 5.1 & 20元 & 65.7\% & 中 & 高空间验证 \\
600906 & 财达证券 & 非银金融 & 7.25元 & 7.6 & 4元 & -40.6\% & 中 & 高估值 \\
603268 & 松发股份 & 国防军工 & 160.46元 & 35.0 & 243元 & 51.7\% & 中高 & 高空间验证 \\
002458 & 益生股份 & 农林牧渔 & 9.02元 & 2.8 & 8元 & -15.0\% & 中 & 估值观察 \\
002384 & 东山精密 & 电子 & 262.49元 & 24.5 & 82元 & -68.7\% & 中高 & 高估值 \\
000977 & 浪潮信息 & 计算机 & 84.70元 & 23.1 & 88元 & 4.0\% & 中高 & 市场观察 \\
603127 & 昭衍新药 & 医药生物 & 50.62元 & 7.0 & 46元 & -9.6\% & 中高 & 估值观察 \\
301200 & 大族数控 & 机械设备 & 312.00元 & 9.5 & 80元 & -74.4\% & 中高 & 高估值 \\
000938 & 紫光股份 & 计算机 & 35.53元 & 18.1 & 39元 & 8.5\% & 中高 & 市场观察 \\
002463 & 沪电股份 & 电子 & 137.23元 & 28.1 & 75元 & -45.3\% & 中高 & 高估值 \\
688183 & 生益电子 & 电子 & 128.87元 & 11.1 & 67元 & -47.8\% & 中 & 高估值 \\
300604 & 长川科技 & 电子 & 304.44元 & 9.1 & 109元 & -64.3\% & 高 & 高估值 \\
002916 & 深南电路 & 电子 & 377.23元 & 21.8 & 185元 & -50.9\% & 中 & 高估值 \\
688825 & 长鑫科技 & 电子 & 未上市；发行价8.66元，暂无二级市场当前价 & 未披露 & 不适用；待上市交易 & 不适用 & 低 & 不可定价 \\
001248 & 华润新能源 & 公用事业 & 13.41元 & 35.4 & 7元 & -50.7\% & 中高 & 质量排除 \\
601633 & 长城汽车 & 汽车 & 15.58元 & 16.2 & 9元 & -44.1\% & 中高 & 质量排除 \\
600736 & 苏州高新 & 房地产 & 6.87元 & 3.0 & 7元 & 4.7\% & 中高 & 质量排除 \\
000415 & 渤海租赁 & 非银金融 & 4.45元 & 17.0 & 4元 & -1.1\% & 中高 & 质量排除 \\
002221 & 东华能源 & 石油石化 & 5.44元 & 0.9 & 1元 & -81.4\% & 中 & 质量排除 \\
002611 & 东方精工 & 机械设备 & 15.02元 & 0.9 & 3元 & -81.2\% & 中 & 质量排除 \\
002174 & 游族网络 & 传媒 & 12.54元 & 0.2 & 1元 & -95.4\% & 中 & 质量排除 \\
603296 & 华勤技术 & 电子 & 78.20元 & 16.8 & 78元 & -0.1\% & 中高 & 质量排除 \\
\bottomrule
\end{longtable}
\normalsize

\begin{riskbox}[候选层使用规则]
H1扣非利润、H2转换比例和行业PE只产生筛选区间。非经常性占比超过40%、预告预减、正分母缺失或研报/现金流质量不足的股票，不得从候选表直接升级为正式配置。
\end{riskbox}
\section{142只逐票估值算法与证据核算}

本节对142只候选逐票执行同一审计接口，但不强迫所有公司使用同一种估值方法。金融股使用PB/ROE与盈利混合，周期行业使用周期调整扣非EPS/PE，成长行业使用成长盈利情景PE，其他有稳定正分母的公司使用扣非EPS情景PE；低基数、结算型和近零EPS标的使用恢复模型，尚未上市的长鑫科技保留上市边界。每行均披露分母、公式、情景值、概率核算、外部锚权重和证据索引。

\scriptsize
\setlength{\tabcolsep}{2pt}
\begin{longtable}{L{0.85cm}L{1.55cm}L{2.25cm}L{6.0cm}L{4.0cm}}
\toprule
\textbf{代码} & \textbf{标的} & \textbf{方法/分母} & \textbf{公式与情景核算} & \textbf{证据/结论} \\
\midrule
\endfirsthead
\toprule
\textbf{代码} & \textbf{标的} & \textbf{方法/分母} & \textbf{公式与情景核算} & \textbf{证据/结论} \\
\midrule
\endhead
000042 & 中洲控股 & normalized EPS / PE with project settlement scenarios；recurring/normalized EPS after project settlement and tax/interest normalization & fair value = 4.80×0.30 + 8.80×0.50 + 13.50×0.20 熊/基准/牛=4.80/8.80/13.50元；概率值=8.54元；目标=8.54元；空间=6.5\%。 & official H1 preview; Q1 financial packet; quote/adjusted K-line; 1 archived source references；证据质量=中；状态=PASS。 \\
600150 & 中国船舶 & growth-earnings scenario PE；H1 deducted EPS bridged with H2/H1 conversion & 成长盈利情景EPS = H1扣非利润中值 / 股本；情景EPS = H1扣非EPS×(1+H2/H1转换比例)；情景价值 = 情景EPS×业务匹配PE；熊值 = min(原始熊值, 当前价×75\%) 熊/基准/牛=25.73/94.39/142.07元；概率值=83.33元；目标=83.33元；空间=142.9\%。 & official H1 preview; Q1 financial packet; quote/adjusted K-line; 华源证券 report 2026-07-14; 2 archived source references；证据质量=中高；状态=PASS。 \\
002379 & 宏桥控股 & cycle-adjusted deducted-EPS PE；H1 deducted EPS bridged with H2/H1 conversion & 周期调整EPS = H1扣非利润中值 / 股本；情景EPS = H1扣非EPS×(1+H2/H1转换比例)；情景价值 = 情景EPS×业务匹配PE；熊值 = min(原始熊值, 当前价×75\%) 熊/基准/牛=14.00/26.24/40.65元；概率值=25.45元；目标=25.80元；空间=38.3\%。 & official H1 preview; Q1 financial packet; quote/adjusted K-line; 招银国际 report 2026-05-13; 2 archived source references；证据质量=高；状态=PASS。 \\
601600 & 中国铝业 & cycle-adjusted deducted-EPS PE；H1 deducted EPS bridged with H2/H1 conversion & 周期调整EPS = H1扣非利润中值 / 股本；情景EPS = H1扣非EPS×(1+H2/H1转换比例)；情景价值 = 情景EPS×业务匹配PE；熊值 = min(原始熊值, 当前价×75\%) 熊/基准/牛=6.74/15.98/23.06元；概率值=14.62元；目标=14.62元；空间=62.8\%。 & official H1 preview; Q1 financial packet; quote/adjusted K-line; 国金证券 report 2026-04-24; 2 archived source references；证据质量=中高；状态=PASS。 \\
600346 & 恒力石化 & cycle-adjusted deducted-EPS PE；H1 deducted EPS bridged with H2/H1 conversion & 周期调整EPS = H1扣非利润中值 / 股本；情景EPS = H1扣非EPS×(1+H2/H1转换比例)；情景价值 = 情景EPS×业务匹配PE；熊值 = min(原始熊值, 当前价×75\%) 熊/基准/牛=7.04/16.02/20.44元；概率值=14.21元；目标=14.21元；空间=-6.8\%。 & official H1 preview; Q1 financial packet; quote/adjusted K-line; 中邮证券 report 2026-07-10; 2 archived source references；证据质量=中高；状态=PASS。 \\
000858 & 五粮液 & deducted-EPS scenario PE；H1 deducted EPS bridged with H2/H1 conversion & 扣非EPS = H1扣非利润中值 / 股本；情景EPS = H1扣非EPS×(1+H2/H1转换比例)；情景价值 = 情景EPS×业务匹配PE；熊值 = min(原始熊值, 当前价×75\%) 熊/基准/牛=44.35/76.61/118.27元；概率值=75.26元；目标=75.26元；空间=-1.5\%。 & official H1 preview; Q1 financial packet; quote/adjusted K-line; 中邮证券 report 2026-05-06; 2 archived source references；证据质量=中高；状态=PASS。 \\
002460 & 赣锋锂业 & cycle-adjusted EPS / PE；2026E cycle-adjusted EPS, not H1 turnaround annualization & fair value = 40.00×0.30 + 53.40×0.50 + 74.90×0.20 熊/基准/牛=40.00/53.40/74.90元；概率值=53.68元；目标=53.68元；空间=4.2\%。 & official H1 preview; Q1 financial packet; quote/adjusted K-line; 东吴证券 report 2026-07-15; 1 archived source references；证据质量=中高；状态=PASS。 \\
002466 & 天齐锂业 & cycle-adjusted deducted-EPS PE；H1 deducted EPS bridged with H2/H1 conversion & 周期调整EPS = H1扣非利润中值 / 股本；情景EPS = H1扣非EPS×(1+H2/H1转换比例)；情景价值 = 情景EPS×业务匹配PE；熊值 = min(原始熊值, 当前价×75\%) 熊/基准/牛=25.37/45.42/70.38元；概率值=44.40元；目标=44.40元；空间=-6.4\%。 & official H1 preview; Q1 financial packet; quote/adjusted K-line; 东莞证券 report 2026-05-22; 2 archived source references；证据质量=中高；状态=PASS。 \\
601688 & 华泰证券 & PB-ROE and earnings blended range；Q1 BPS plus H1 deducted EPS bridged to H2 & scenario value = 65\%×(Q1 BPS×PB multiple) + 35\%×(scenario EPS×PE multiple) 熊/基准/牛=15.47/23.29/31.38元；概率值=22.56元；目标=22.56元；空间=9.4\%。 & official H1 preview; Q1 financial packet; quote/adjusted K-line; 华源证券 report 2026-05-09; 2 archived source references；证据质量=中高；状态=PASS。 \\
601336 & 新华保险 & PB-ROE and earnings blended range；Q1 BPS plus H1 deducted EPS bridged to H2 & scenario value = 65\%×(Q1 BPS×PB multiple) + 35\%×(scenario EPS×PE multiple) 熊/基准/牛=49.16/82.01/113.47元；概率值=78.45元；目标=78.45元；空间=19.7\%。 & official H1 preview; Q1 financial packet; quote/adjusted K-line; 东吴证券 report 2026-07-14; 2 archived source references；证据质量=中高；状态=PASS。 \\
600037 & 歌华有线 & normalized EPS / PE with turnaround discount；2026E recurring EPS after turnaround, not H1 profit annualization & fair value = 2.40×0.30 + 5.40×0.50 + 11.00×0.20 熊/基准/牛=2.40/5.40/11.00元；概率值=5.62元；目标=5.62元；空间=-17.4\%。 & official H1 preview; Q1 financial packet; quote/adjusted K-line; 中邮证券 report 2026-03-23; 1 archived source references；证据质量=中；状态=PASS。 \\
002738 & 中矿资源 & cycle-adjusted deducted-EPS PE；H1 deducted EPS bridged with H2/H1 conversion & 周期调整EPS = H1扣非利润中值 / 股本；情景EPS = H1扣非EPS×(1+H2/H1转换比例)；情景价值 = 情景EPS×业务匹配PE；熊值 = min(原始熊值, 当前价×75\%) 熊/基准/牛=18.91/46.20/52.45元；概率值=39.26元；目标=39.26元；空间=-20.4\%。 & official H1 preview; Q1 financial packet; quote/adjusted K-line; 东吴证券 report 2026-07-14; 2 archived source references；证据质量=中高；状态=PASS。 \\
601069 & 西部黄金 & cycle-adjusted deducted-EPS PE；H1 deducted EPS bridged with H2/H1 conversion & 周期调整EPS = H1扣非利润中值 / 股本；情景EPS = H1扣非EPS×(1+H2/H1转换比例)；情景价值 = 情景EPS×业务匹配PE；熊值 = min(原始熊值, 当前价×75\%) 熊/基准/牛=7.28/13.04/20.20元；概率值=12.74元；目标=12.74元；空间=-44.5\%。 & official H1 preview; Q1 financial packet; quote/adjusted K-line; 民生证券 report 2025-10-31; 2 archived source references；证据质量=中；状态=PASS。 \\
000567 & 海德股份 & PB-ROE and earnings blended range；Q1 BPS plus H1 deducted EPS bridged to H2 & scenario value = 65\%×(Q1 BPS×PB multiple) + 35\%×(scenario EPS×PE multiple) 熊/基准/牛=4.63/7.78/10.83元；概率值=7.45元；目标=7.45元；空间=20.8\%。 & official H1 preview; Q1 financial packet; quote/adjusted K-line; 太平洋 report 2024-06-20; 2 archived source references；证据质量=中；状态=PASS。 \\
002602 & 世纪华通 & growth-earnings scenario PE；H1 deducted EPS bridged with H2/H1 conversion & 成长盈利情景EPS = H1扣非利润中值 / 股本；情景EPS = H1扣非EPS×(1+H2/H1转换比例)；情景价值 = 情景EPS×业务匹配PE；熊值 = min(原始熊值, 当前价×75\%) 熊/基准/牛=9.95/16.94/25.29元；概率值=16.51元；目标=16.51元；空间=15.8\%。 & official H1 preview; Q1 financial packet; quote/adjusted K-line; 开源证券 report 2026-06-30; 2 archived source references；证据质量=中高；状态=PASS。 \\
600026 & 中远海能 & cycle-adjusted deducted-EPS PE；H1 deducted EPS bridged with H2/H1 conversion & 周期调整EPS = H1扣非利润中值 / 股本；情景EPS = H1扣非EPS×(1+H2/H1转换比例)；情景价值 = 情景EPS×业务匹配PE；熊值 = min(原始熊值, 当前价×75\%) 熊/基准/牛=9.97/25.64/27.66元；概率值=21.34元；目标=21.34元；空间=34.8\%。 & official H1 preview; Q1 financial packet; quote/adjusted K-line; 国金证券 report 2026-07-12; 2 archived source references；证据质量=中高；状态=PASS。 \\
002048 & 宁波华翔 & normalized EPS / PE；2026E normalized EPS after auto-cycle and robot-order verification & fair value = 19.20×0.30 + 29.64×0.50 + 39.60×0.20 熊/基准/牛=19.20/29.64/39.60元；概率值=28.50元；目标=28.50元；空间=33.0\%。 & official H1 preview; Q1 financial packet; quote/adjusted K-line; 国金证券 report 2025-10-30; 1 archived source references；证据质量=中；状态=PASS。 \\
600688 & 上海石化 & PB / cycle-adjusted asset value；Q1 BPS and cycle-adjusted asset value & fair value = 2.01×0.30 + 2.68×0.50 + 3.35×0.20 熊/基准/牛=2.01/2.68/3.35元；概率值=2.61元；目标=2.61元；空间=-8.1\%。 & official H1 preview; Q1 financial packet; quote/adjusted K-line; 西南证券 report 2023-03-30; 1 archived source references；证据质量=中；状态=PASS。 \\
600095 & 湘财股份 & PB-ROE and earnings blended range；Q1 BPS plus H1 deducted EPS bridged to H2 & scenario value = 65\%×(Q1 BPS×PB multiple) + 35\%×(scenario EPS×PE multiple) 熊/基准/牛=2.95/4.32/5.76元；概率值=4.20元；目标=4.20元；空间=-46.2\%。 & official H1 preview; Q1 financial packet; quote/adjusted K-line; 国金证券 report 2026-03-24; 2 archived source references；证据质量=中；状态=PASS。 \\
000426 & 兴业银锡 & cycle-adjusted deducted-EPS PE；H1 deducted EPS bridged with H2/H1 conversion & 周期调整EPS = H1扣非利润中值 / 股本；情景EPS = H1扣非EPS×(1+H2/H1转换比例)；情景价值 = 情景EPS×业务匹配PE；熊值 = min(原始熊值, 当前价×75\%) 熊/基准/牛=12.61/25.56/34.97元；概率值=23.56元；目标=23.56元；空间=-22.9\%。 & official H1 preview; Q1 financial packet; quote/adjusted K-line; 国信证券 report 2026-04-24; 2 archived source references；证据质量=中高；状态=PASS。 \\
600489 & 中金黄金 & cycle-adjusted deducted-EPS PE；H1 deducted EPS bridged with H2/H1 conversion & 周期调整EPS = H1扣非利润中值 / 股本；情景EPS = H1扣非EPS×(1+H2/H1转换比例)；情景价值 = 情景EPS×业务匹配PE；熊值 = min(原始熊值, 当前价×75\%) 熊/基准/牛=11.00/23.52/30.52元；概率值=21.16元；目标=21.16元；空间=5.2\%。 & official H1 preview; Q1 financial packet; quote/adjusted K-line; 中邮证券 report 2026-05-08; 2 archived source references；证据质量=中高；状态=PASS。 \\
000737 & 北方铜业 & cycle-adjusted deducted-EPS PE；H1 deducted EPS bridged with H2/H1 conversion & 周期调整EPS = H1扣非利润中值 / 股本；情景EPS = H1扣非EPS×(1+H2/H1转换比例)；情景价值 = 情景EPS×业务匹配PE；熊值 = min(原始熊值, 当前价×75\%) 熊/基准/牛=8.80/15.76/24.42元；概率值=15.40元；目标=15.40元；空间=20.8\%。 & official H1 preview; Q1 financial packet; quote/adjusted K-line; 1 archived source references；证据质量=中低；状态=PASS。 \\
002756 & 永兴材料 & cycle-adjusted deducted-EPS PE；H1 deducted EPS bridged with H2/H1 conversion & 周期调整EPS = H1扣非利润中值 / 股本；情景EPS = H1扣非EPS×(1+H2/H1转换比例)；情景价值 = 情景EPS×业务匹配PE；熊值 = min(原始熊值, 当前价×75\%) 熊/基准/牛=23.00/44.88/63.81元；概率值=42.10元；目标=42.10元；空间=-6.2\%。 & official H1 preview; Q1 financial packet; quote/adjusted K-line; 东吴证券 report 2026-07-14; 2 archived source references；证据质量=中高；状态=PASS。 \\
300014 & 亿纬锂能 & growth-earnings scenario PE；H1 deducted EPS bridged with H2/H1 conversion & 成长盈利情景EPS = H1扣非利润中值 / 股本；情景EPS = H1扣非EPS×(1+H2/H1转换比例)；情景价值 = 情景EPS×业务匹配PE；熊值 = min(原始熊值, 当前价×75\%) 熊/基准/牛=31.49/68.46/74.57元；概率值=58.59元；目标=58.59元；空间=12.1\%。 & official H1 preview; Q1 financial packet; quote/adjusted K-line; 群益证券 report 2026-06-26; 2 archived source references；证据质量=中高；状态=PASS。 \\
601061 & 中信金属 & deducted-EPS scenario PE；H1 deducted EPS bridged with H2/H1 conversion & 扣非EPS = H1扣非利润中值 / 股本；情景EPS = H1扣非EPS×(1+H2/H1转换比例)；情景价值 = 情景EPS×业务匹配PE；熊值 = min(原始熊值, 当前价×75\%) 熊/基准/牛=8.12/18.46/28.50元；概率值=17.37元；目标=17.06元；空间=57.7\%。 & official H1 preview; Q1 financial packet; quote/adjusted K-line; 西南证券 report 2025-09-05; 2 archived source references；证据质量=中；状态=PASS。 \\
600739 & 辽宁成大 & growth-earnings scenario PE；H1 deducted EPS bridged with H2/H1 conversion & 成长盈利情景EPS = H1扣非利润中值 / 股本；情景EPS = H1扣非EPS×(1+H2/H1转换比例)；情景价值 = 情景EPS×业务匹配PE；熊值 = min(原始熊值, 当前价×75\%) 熊/基准/牛=8.21/52.03/76.60元；概率值=43.80元；目标=42.75元；空间=290.4\%。 & official H1 preview; Q1 financial packet; quote/adjusted K-line; 太平洋 report 2017-10-30; 2 archived source references；证据质量=中；状态=PASS。 \\
601360 & 三六零 & PS on 2026E revenue；2026E revenue consensus and recurring AI/security monetization, not H1 turnaround EPS & fair value = 10.28×0.30 + 13.70×0.50 + 15.07×0.20 熊/基准/牛=10.28/13.70/15.07元；概率值=12.95元；目标=13.10元；空间=45.6\%。 & official H1 preview; Q1 financial packet; quote/adjusted K-line; 信达证券 report 2023-04-25; 1 archived source references；证据质量=中；状态=PASS。 \\
600918 & 中泰证券 & PB-ROE and earnings blended range；Q1 BPS plus H1 deducted EPS bridged to H2 & scenario value = 65\%×(Q1 BPS×PB multiple) + 35\%×(scenario EPS×PE multiple) 熊/基准/牛=3.77/5.50/7.32元；概率值=5.35元；目标=5.35元；空间=-6.3\%。 & official H1 preview; Q1 financial packet; quote/adjusted K-line; 太平洋 report 2025-10-31; 2 archived source references；证据质量=中；状态=PASS。 \\
000975 & 山金国际 & cycle-adjusted deducted-EPS PE；H1 deducted EPS bridged with H2/H1 conversion & 周期调整EPS = H1扣非利润中值 / 股本；情景EPS = H1扣非EPS×(1+H2/H1转换比例)；情景价值 = 情景EPS×业务匹配PE；熊值 = min(原始熊值, 当前价×75\%) 熊/基准/牛=10.49/22.68/29.10元；概率值=20.31元；目标=20.31元；空间=5.0\%。 & official H1 preview; Q1 financial packet; quote/adjusted K-line; 太平洋 report 2026-05-13; 2 archived source references；证据质量=中高；状态=PASS。 \\
000623 & 吉林敖东 & growth-earnings scenario PE；H1 deducted EPS bridged with H2/H1 conversion & 成长盈利情景EPS = H1扣非利润中值 / 股本；情景EPS = H1扣非EPS×(1+H2/H1转换比例)；情景价值 = 情景EPS×业务匹配PE；熊值 = min(原始熊值, 当前价×75\%) 熊/基准/牛=14.00/92.67/136.41元；概率值=77.82元；目标=77.82元；空间=317.0\%。 & official H1 preview; Q1 financial packet; quote/adjusted K-line; 国元证券 report 2024-09-03; 2 archived source references；证据质量=中；状态=PASS。 \\
600292 & 电投水电 & deducted-EPS scenario PE；H1 deducted EPS bridged with H2/H1 conversion & 扣非EPS = H1扣非利润中值 / 股本；情景EPS = H1扣非EPS×(1+H2/H1转换比例)；情景价值 = 情景EPS×业务匹配PE；熊值 = min(原始熊值, 当前价×75\%) 熊/基准/牛=4.01/6.67/10.10元；概率值=6.56元；目标=6.56元；空间=-47.5\%。 & official H1 preview; Q1 financial packet; quote/adjusted K-line; 西南证券 report 2025-09-12; 2 archived source references；证据质量=中；状态=PASS。 \\
000987 & 越秀资本 & PB-ROE and earnings blended range；Q1 BPS plus H1 deducted EPS bridged to H2 & scenario value = 65\%×(Q1 BPS×PB multiple) + 35\%×(scenario EPS×PE multiple) 熊/基准/牛=5.45/8.11/10.95元；概率值=7.88元；目标=7.96元；空间=2.5\%。 & official H1 preview; Q1 financial packet; quote/adjusted K-line; 太平洋 report 2019-08-19; 2 archived source references；证据质量=中；状态=PASS。 \\
000155 & 川能动力 & deducted-EPS scenario PE；H1 deducted EPS bridged with H2/H1 conversion & 扣非EPS = H1扣非利润中值 / 股本；情景EPS = H1扣非EPS×(1+H2/H1转换比例)；情景价值 = 情景EPS×业务匹配PE；熊值 = min(原始熊值, 当前价×75\%) 熊/基准/牛=5.43/9.03/13.67元；概率值=8.88元；目标=10.04元；空间=-9.5\%。 & official H1 preview; Q1 financial packet; quote/adjusted K-line; 国信证券 report 2026-04-20; 2 archived source references；证据质量=高；状态=PASS。 \\
601377 & 兴业证券 & PB-ROE and earnings blended range；Q1 BPS plus H1 deducted EPS bridged to H2 & scenario value = 65\%×(Q1 BPS×PB multiple) + 35\%×(scenario EPS×PE multiple) 熊/基准/牛=4.35/6.34/8.43元；概率值=6.16元；目标=6.16元；空间=-0.8\%。 & official H1 preview; Q1 financial packet; quote/adjusted K-line; 开源证券 report 2026-05-05; 2 archived source references；证据质量=中高；状态=PASS。 \\
600120 & 浙江东方 & PB-ROE and earnings blended range；Q1 BPS plus H1 deducted EPS bridged to H2 & scenario value = 65\%×(Q1 BPS×PB multiple) + 35\%×(scenario EPS×PE multiple) 熊/基准/牛=3.51/5.17/6.93元；概率值=5.02元；目标=5.02元；空间=2.2\%。 & official H1 preview; Q1 financial packet; quote/adjusted K-line; 1 archived source references；证据质量=中低；状态=PASS。 \\
000703 & 恒逸石化 & cycle-adjusted deducted-EPS PE；H1 deducted EPS bridged with H2/H1 conversion & 周期调整EPS = H1扣非利润中值 / 股本；情景EPS = H1扣非EPS×(1+H2/H1转换比例)；情景价值 = 情景EPS×业务匹配PE；熊值 = min(原始熊值, 当前价×75\%) 熊/基准/牛=10.86/25.59/40.41元；概率值=24.13元；目标=23.42元；空间=61.7\%。 & official H1 preview; Q1 financial packet; quote/adjusted K-line; 国信证券 report 2026-07-05; 2 archived source references；证据质量=高；状态=PASS。 \\
000301 & 东方盛虹 & cycle-adjusted deducted-EPS PE；H1 deducted EPS bridged with H2/H1 conversion & 周期调整EPS = H1扣非利润中值 / 股本；情景EPS = H1扣非EPS×(1+H2/H1转换比例)；情景价值 = 情景EPS×业务匹配PE；熊值 = min(原始熊值, 当前价×75\%) 熊/基准/牛=6.21/11.42/18.03元；概率值=11.18元；目标=11.18元；空间=-3.7\%。 & official H1 preview; Q1 financial packet; quote/adjusted K-line; 中邮证券 report 2026-05-06; 2 archived source references；证据质量=中高；状态=PASS。 \\
001339 & 智微智能 & growth-earnings scenario PE；H1 deducted EPS bridged with H2/H1 conversion & 成长盈利情景EPS = H1扣非利润中值 / 股本；情景EPS = H1扣非EPS×(1+H2/H1转换比例)；情景价值 = 情景EPS×业务匹配PE；熊值 = min(原始熊值, 当前价×75\%) 熊/基准/牛=34.27/69.44/102.13元；概率值=65.43元；目标=65.43元；空间=-26.7\%。 & official H1 preview; Q1 financial packet; quote/adjusted K-line; 中邮证券 report 2026-06-18; 2 archived source references；证据质量=中高；状态=PASS。 \\
002558 & 巨人网络 & growth-earnings scenario PE；H1 deducted EPS bridged with H2/H1 conversion & 成长盈利情景EPS = H1扣非利润中值 / 股本；情景EPS = H1扣非EPS×(1+H2/H1转换比例)；情景价值 = 情景EPS×业务匹配PE；熊值 = min(原始熊值, 当前价×75\%) 熊/基准/牛=20.57/34.73/52.28元；概率值=33.99元；目标=33.99元；空间=15.0\%。 & official H1 preview; Q1 financial packet; quote/adjusted K-line; 西南证券 report 2026-07-15; 2 archived source references；证据质量=中高；状态=PASS。 \\
301308 & 江波龙 & growth-earnings scenario PE；H1 deducted EPS bridged with H2/H1 conversion & 成长盈利情景EPS = H1扣非利润中值 / 股本；情景EPS = H1扣非EPS×(1+H2/H1转换比例)；情景价值 = 情景EPS×业务匹配PE；熊值 = min(原始熊值, 当前价×75\%) 熊/基准/牛=366.00/1198.41/1802.22元；概率值=1069.45元；目标=1069.45元；空间=119.2\%。 & official H1 preview; Q1 financial packet; quote/adjusted K-line; 国信证券 report 2026-05-07; 2 archived source references；证据质量=中高；状态=PASS。 \\
601628 & 中国人寿 & PB-ROE and earnings blended range；Q1 BPS plus H1 deducted EPS bridged to H2 & scenario value = 65\%×(Q1 BPS×PB multiple) + 35\%×(scenario EPS×PE multiple) 熊/基准/牛=30.32/50.75/70.56元；概率值=48.58元；目标=48.58元；空间=20.2\%。 & official H1 preview; Q1 financial packet; quote/adjusted K-line; 东吴证券 report 2026-07-15; 2 archived source references；证据质量=中高；状态=PASS。 \\
002493 & 荣盛石化 & cycle-adjusted deducted-EPS PE；H1 deducted EPS bridged with H2/H1 conversion & 周期调整EPS = H1扣非利润中值 / 股本；情景EPS = H1扣非EPS×(1+H2/H1转换比例)；情景价值 = 情景EPS×业务匹配PE；熊值 = min(原始熊值, 当前价×75\%) 熊/基准/牛=4.93/9.07/14.33元；概率值=8.88元；目标=8.88元；空间=-20.0\%。 & official H1 preview; Q1 financial packet; quote/adjusted K-line; 国信证券 report 2026-07-15; 2 archived source references；证据质量=中高；状态=PASS。 \\
001309 & 德明利 & forward EPS / PE with memory-cycle discount；2026E memory-cycle-adjusted EPS, not H1 annualized EPS & fair value = 330.00×0.30 + 513.00×0.50 + 630.00×0.20 熊/基准/牛=330.00/513.00/630.00元；概率值=481.50元；目标=481.50元；空间=-27.3\%。 & official H1 preview; Q1 financial packet; quote/adjusted K-line; 国信证券 report 2026-04-21; 1 archived source references；证据质量=中高；状态=PASS。 \\
601233 & 桐昆股份 & cycle-adjusted deducted-EPS PE；H1 deducted EPS bridged with H2/H1 conversion & 周期调整EPS = H1扣非利润中值 / 股本；情景EPS = H1扣非EPS×(1+H2/H1转换比例)；情景价值 = 情景EPS×业务匹配PE；熊值 = min(原始熊值, 当前价×75\%) 熊/基准/牛=15.54/29.33/46.31元；概率值=28.59元；目标=28.59元；空间=38.0\%。 & official H1 preview; Q1 financial packet; quote/adjusted K-line; 东海证券 report 2026-05-08; 2 archived source references；证据质量=中高；状态=PASS。 \\
603986 & 兆易创新 & growth-earnings scenario PE；H1 deducted EPS bridged with H2/H1 conversion & 成长盈利情景EPS = H1扣非利润中值 / 股本；情景EPS = H1扣非EPS×(1+H2/H1转换比例)；情景价值 = 情景EPS×业务匹配PE；熊值 = min(原始熊值, 当前价×75\%) 熊/基准/牛=211.49/359.39/540.47元；概率值=351.24元；目标=351.24元；空间=-38.6\%。 & official H1 preview; Q1 financial packet; quote/adjusted K-line; 华鑫证券 report 2026-06-29; 2 archived source references；证据质量=中高；状态=PASS。 \\
601225 & 陕西煤业 & cycle-adjusted deducted-EPS PE；H1 deducted EPS bridged with H2/H1 conversion & 周期调整EPS = H1扣非利润中值 / 股本；情景EPS = H1扣非EPS×(1+H2/H1转换比例)；情景价值 = 情景EPS×业务匹配PE；熊值 = min(原始熊值, 当前价×75\%) 熊/基准/牛=10.91/18.72/24.88元；概率值=17.61元；目标=17.61元；空间=-28.4\%。 & official H1 preview; Q1 financial packet; quote/adjusted K-line; 国信证券 report 2026-04-28; 2 archived source references；证据质量=中高；状态=PASS。 \\
601138 & 工业富联 & growth-earnings scenario PE；H1 deducted EPS bridged with H2/H1 conversion & 成长盈利情景EPS = H1扣非利润中值 / 股本；情景EPS = H1扣非EPS×(1+H2/H1转换比例)；情景价值 = 情景EPS×业务匹配PE；熊值 = min(原始熊值, 当前价×75\%) 熊/基准/牛=35.77/79.56/91.42元；概率值=68.80元；目标=68.80元；空间=8.1\%。 & official H1 preview; Q1 financial packet; quote/adjusted K-line; 金元证券 report 2026-06-18; 2 archived source references；证据质量=中高；状态=PASS。 \\
000807 & 云铝股份 & cycle-adjusted deducted-EPS PE；H1 deducted EPS bridged with H2/H1 conversion & 周期调整EPS = H1扣非利润中值 / 股本；情景EPS = H1扣非EPS×(1+H2/H1转换比例)；情景价值 = 情景EPS×业务匹配PE；熊值 = min(原始熊值, 当前价×75\%) 熊/基准/牛=19.58/48.33/74.89元；概率值=45.02元；目标=45.02元；空间=72.5\%。 & official H1 preview; Q1 financial packet; quote/adjusted K-line; 东吴证券 report 2026-04-27; 2 archived source references；证据质量=中高；状态=PASS。 \\
603993 & 洛阳钼业 & cycle-adjusted deducted-EPS PE；H1 deducted EPS bridged with H2/H1 conversion & 周期调整EPS = H1扣非利润中值 / 股本；情景EPS = H1扣非EPS×(1+H2/H1转换比例)；情景价值 = 情景EPS×业务匹配PE；熊值 = min(原始熊值, 当前价×75\%) 熊/基准/牛=8.98/17.52/24.92元；概率值=16.44元；目标=16.44元；空间=-8.2\%。 & official H1 preview; Q1 financial packet; quote/adjusted K-line; 太平洋 report 2026-05-14; 2 archived source references；证据质量=中高；状态=PASS。 \\
601872 & 招商轮船 & cycle-adjusted deducted-EPS PE；H1 deducted EPS bridged with H2/H1 conversion & 周期调整EPS = H1扣非利润中值 / 股本；情景EPS = H1扣非EPS×(1+H2/H1转换比例)；情景价值 = 情景EPS×业务匹配PE；熊值 = min(原始熊值, 当前价×75\%) 熊/基准/牛=10.58/25.08/29.35元；概率值=21.58元；目标=21.58元；空间=41.4\%。 & official H1 preview; Q1 financial packet; quote/adjusted K-line; 华源证券 report 2026-07-06; 2 archived source references；证据质量=中高；状态=PASS。 \\
600999 & 招商证券 & PB-ROE and earnings blended range；Q1 BPS plus H1 deducted EPS bridged to H2 & scenario value = 65\%×(Q1 BPS×PB multiple) + 35\%×(scenario EPS×PE multiple) 熊/基准/牛=12.72/19.05/25.86元；概率值=18.51元；目标=18.51元；空间=-11.4\%。 & official H1 preview; Q1 financial packet; quote/adjusted K-line; 国信证券 report 2026-07-10; 2 archived source references；证据质量=中高；状态=PASS。 \\
601899 & 紫金矿业 & cycle-adjusted deducted-EPS PE；H1 deducted EPS bridged with H2/H1 conversion & 周期调整EPS = H1扣非利润中值 / 股本；情景EPS = H1扣非EPS×(1+H2/H1转换比例)；情景价值 = 情景EPS×业务匹配PE；熊值 = min(原始熊值, 当前价×75\%) 熊/基准/牛=17.67/40.08/49.03元；概率值=35.15元；目标=35.15元；空间=21.9\%。 & official H1 preview; Q1 financial packet; quote/adjusted K-line; 中邮证券 report 2026-04-28; 2 archived source references；证据质量=中高；状态=PASS。 \\
600309 & 万华化学 & cycle-adjusted deducted-EPS PE；H1 deducted EPS bridged with H2/H1 conversion & 周期调整EPS = H1扣非利润中值 / 股本；情景EPS = H1扣非EPS×(1+H2/H1转换比例)；情景价值 = 情景EPS×业务匹配PE；熊值 = min(原始熊值, 当前价×75\%) 熊/基准/牛=49.02/81.91/122.39元；概率值=80.14元；目标=80.14元；空间=14.3\%。 & official H1 preview; Q1 financial packet; quote/adjusted K-line; 国信证券 report 2026-07-07; 2 archived source references；证据质量=中高；状态=PASS。 \\
600989 & 宝丰能源 & cycle-adjusted deducted-EPS PE；H1 deducted EPS bridged with H2/H1 conversion & 周期调整EPS = H1扣非利润中值 / 股本；情景EPS = H1扣非EPS×(1+H2/H1转换比例)；情景价值 = 情景EPS×业务匹配PE；熊值 = min(原始熊值, 当前价×75\%) 熊/基准/牛=16.55/32.32/48.29元；概率值=30.78元；目标=31.66元；空间=43.5\%。 & official H1 preview; Q1 financial packet; quote/adjusted K-line; 环球富盛理财 report 2026-05-05; 2 archived source references；证据质量=高；状态=PASS。 \\
300475 & 香农芯创 & growth-earnings scenario PE；H1 deducted EPS bridged with H2/H1 conversion & 成长盈利情景EPS = H1扣非利润中值 / 股本；情景EPS = H1扣非EPS×(1+H2/H1转换比例)；情景价值 = 情景EPS×业务匹配PE；熊值 = min(原始熊值, 当前价×75\%) 熊/基准/牛=165.90/403.12/606.22元；概率值=372.57元；目标=372.57元；空间=68.4\%。 & official H1 preview; Q1 financial packet; quote/adjusted K-line; 华鑫证券 report 2026-03-22; 2 archived source references；证据质量=中；状态=PASS。 \\
000776 & 广发证券 & PB-ROE and earnings blended range；Q1 BPS plus H1 deducted EPS bridged to H2 & scenario value = 65\%×(Q1 BPS×PB multiple) + 35\%×(scenario EPS×PE multiple) 熊/基准/牛=15.69/23.52/31.92元；概率值=22.85元；目标=22.85元；空间=0.3\%。 & official H1 preview; Q1 financial packet; quote/adjusted K-line; 华源证券 report 2026-04-30; 2 archived source references；证据质量=中高；状态=PASS。 \\
600030 & 中信证券 & PB-ROE and earnings blended range；Q1 BPS plus H1 deducted EPS bridged to H2 & scenario value = 65\%×(Q1 BPS×PB multiple) + 35\%×(scenario EPS×PE multiple) 熊/基准/牛=17.17/25.69/34.82元；概率值=24.96元；目标=24.96元；空间=-12.7\%。 & official H1 preview; Q1 financial packet; quote/adjusted K-line; 东吴证券 report 2026-07-10; 2 archived source references；证据质量=中高；状态=PASS。 \\
601088 & 中国神华 & cycle-adjusted deducted-EPS PE；H1 deducted EPS bridged with H2/H1 conversion & 周期调整EPS = H1扣非利润中值 / 股本；情景EPS = H1扣非EPS×(1+H2/H1转换比例)；情景价值 = 情景EPS×业务匹配PE；熊值 = min(原始熊值, 当前价×75\%) 熊/基准/牛=13.56/22.86/30.92元；概率值=21.68元；目标=21.68元；空间=-50.4\%。 & official H1 preview; Q1 financial packet; quote/adjusted K-line; 华源证券 report 2026-04-11; 2 archived source references；证据质量=中；状态=PASS。 \\
600362 & 江西铜业 & cycle-adjusted deducted-EPS PE；H1 deducted EPS bridged with H2/H1 conversion & 周期调整EPS = H1扣非利润中值 / 股本；情景EPS = H1扣非EPS×(1+H2/H1转换比例)；情景价值 = 情景EPS×业务匹配PE；熊值 = min(原始熊值, 当前价×75\%) 熊/基准/牛=29.02/51.96/80.52元；概率值=50.79元；目标=50.79元；空间=20.1\%。 & official H1 preview; Q1 financial packet; quote/adjusted K-line; 国信证券 report 2026-04-03; 2 archived source references；证据质量=中；状态=PASS。 \\
002709 & 天赐材料 & growth-earnings scenario PE；H1 deducted EPS bridged with H2/H1 conversion & 成长盈利情景EPS = H1扣非利润中值 / 股本；情景EPS = H1扣非EPS×(1+H2/H1转换比例)；情景价值 = 情景EPS×业务匹配PE；熊值 = min(原始熊值, 当前价×75\%) 熊/基准/牛=28.16/67.32/88.45元；概率值=59.80元；目标=60.46元；空间=61.0\%。 & official H1 preview; Q1 financial packet; quote/adjusted K-line; 东吴证券 report 2026-07-10; 2 archived source references；证据质量=高；状态=PASS。 \\
002414 & 高德红外 & growth-earnings scenario PE；H1 deducted EPS bridged with H2/H1 conversion & 成长盈利情景EPS = H1扣非利润中值 / 股本；情景EPS = H1扣非EPS×(1+H2/H1转换比例)；情景价值 = 情景EPS×业务匹配PE；熊值 = min(原始熊值, 当前价×75\%) 熊/基准/牛=9.99/22.26/33.51元；概率值=20.83元；目标=20.79元；空间=56.1\%。 & official H1 preview; Q1 financial packet; quote/adjusted K-line; 太平洋 report 2025-11-06; 2 archived source references；证据质量=中；状态=PASS。 \\
300390 & 天华新能 & cycle-adjusted EPS / PE；2026E cycle-adjusted EPS and lithium-price sensitivity & fair value = 56.40×0.30 + 79.80×0.50 + 102.40×0.20 熊/基准/牛=56.40/79.80/102.40元；概率值=77.30元；目标=77.30元；空间=28.8\%。 & official H1 preview; Q1 financial packet; quote/adjusted K-line; 华源证券 report 2026-06-06; 1 archived source references；证据质量=中高；状态=PASS。 \\
002648 & 卫星化学 & cycle-adjusted deducted-EPS PE；H1 deducted EPS bridged with H2/H1 conversion & 周期调整EPS = H1扣非利润中值 / 股本；情景EPS = H1扣非EPS×(1+H2/H1转换比例)；情景价值 = 情景EPS×业务匹配PE；熊值 = min(原始熊值, 当前价×75\%) 熊/基准/牛=17.57/48.02/71.08元；概率值=43.50元；目标=43.50元；空间=85.7\%。 & official H1 preview; Q1 financial packet; quote/adjusted K-line; 华源证券 report 2026-07-09; 2 archived source references；证据质量=中高；状态=PASS。 \\
000792 & 盐湖股份 & cycle-adjusted deducted-EPS PE；H1 deducted EPS bridged with H2/H1 conversion & 周期调整EPS = H1扣非利润中值 / 股本；情景EPS = H1扣非EPS×(1+H2/H1转换比例)；情景价值 = 情景EPS×业务匹配PE；熊值 = min(原始熊值, 当前价×75\%) 熊/基准/牛=18.01/30.10/44.98元；概率值=29.45元；目标=29.45元；空间=16.5\%。 & official H1 preview; Q1 financial packet; quote/adjusted K-line; 华鑫证券 report 2026-04-09; 2 archived source references；证据质量=中；状态=PASS。 \\
601869 & 长飞光纤 & growth-earnings scenario PE；H1 deducted EPS bridged with H2/H1 conversion & 成长盈利情景EPS = H1扣非利润中值 / 股本；情景EPS = H1扣非EPS×(1+H2/H1转换比例)；情景价值 = 情景EPS×业务匹配PE；熊值 = min(原始熊值, 当前价×75\%) 熊/基准/牛=85.01/157.30/217.25元；概率值=147.60元；目标=147.60元；空间=-65.3\%。 & official H1 preview; Q1 financial packet; quote/adjusted K-line; 华鑫证券 report 2026-04-12; 2 archived source references；证据质量=中；状态=PASS。 \\
002432 & 九安医疗 & growth-earnings scenario PE；H1 deducted EPS bridged with H2/H1 conversion & 成长盈利情景EPS = H1扣非利润中值 / 股本；情景EPS = H1扣非EPS×(1+H2/H1转换比例)；情景价值 = 情景EPS×业务匹配PE；熊值 = min(原始熊值, 当前价×75\%) 熊/基准/牛=50.24/338.06/497.63元；概率值=283.63元；目标=283.63元；空间=323.4\%。 & official H1 preview; Q1 financial packet; quote/adjusted K-line; 东吴证券 report 2025-05-06; 2 archived source references；证据质量=中；状态=PASS。 \\
000657 & 中钨高新 & cycle-adjusted deducted-EPS PE；H1 deducted EPS bridged with H2/H1 conversion & 周期调整EPS = H1扣非利润中值 / 股本；情景EPS = H1扣非EPS×(1+H2/H1转换比例)；情景价值 = 情景EPS×业务匹配PE；熊值 = min(原始熊值, 当前价×75\%) 熊/基准/牛=11.18/27.00/31.02元；概率值=23.06元；目标=23.06元；空间=-67.0\%。 & official H1 preview; Q1 financial packet; quote/adjusted K-line; 东吴证券 report 2026-07-14; 2 archived source references；证据质量=中高；状态=PASS。 \\
601995 & 中金公司 & PB-ROE and earnings blended range；Q1 BPS plus H1 deducted EPS bridged to H2 & scenario value = 65\%×(Q1 BPS×PB multiple) + 35\%×(scenario EPS×PE multiple) 熊/基准/牛=18.07/26.97/36.49元；概率值=26.20元；目标=26.20元；空间=-31.7\%。 & official H1 preview; Q1 financial packet; quote/adjusted K-line; 国信证券 report 2026-05-05; 2 archived source references；证据质量=中高；状态=PASS。 \\
002497 & 雅化集团 & cycle-adjusted deducted-EPS PE；H1 deducted EPS bridged with H2/H1 conversion & 周期调整EPS = H1扣非利润中值 / 股本；情景EPS = H1扣非EPS×(1+H2/H1转换比例)；情景价值 = 情景EPS×业务匹配PE；熊值 = min(原始熊值, 当前价×75\%) 熊/基准/牛=13.55/36.82/40.96元；概率值=30.67元；目标=31.80元；空间=76.0\%。 & official H1 preview; Q1 financial packet; quote/adjusted K-line; 东吴证券 report 2026-07-08; 2 archived source references；证据质量=高；状态=PASS。 \\
000933 & 神火股份 & cycle-adjusted deducted-EPS PE；H1 deducted EPS bridged with H2/H1 conversion & 周期调整EPS = H1扣非利润中值 / 股本；情景EPS = H1扣非EPS×(1+H2/H1转换比例)；情景价值 = 情景EPS×业务匹配PE；熊值 = min(原始熊值, 当前价×75\%) 熊/基准/牛=19.13/47.54/73.51元；概率值=44.21元；目标=44.21元；空间=73.3\%。 & official H1 preview; Q1 financial packet; quote/adjusted K-line; 国金证券 report 2026-04-14; 2 archived source references；证据质量=中；状态=PASS。 \\
300223 & 北京君正 & growth-earnings scenario PE；H1 deducted EPS bridged with H2/H1 conversion & 成长盈利情景EPS = H1扣非利润中值 / 股本；情景EPS = H1扣非EPS×(1+H2/H1转换比例)；情景价值 = 情景EPS×业务匹配PE；熊值 = min(原始熊值, 当前价×75\%) 熊/基准/牛=72.82/123.75/186.09元；概率值=120.94元；目标=120.94元；空间=-36.4\%。 & official H1 preview; Q1 financial packet; quote/adjusted K-line; 国信证券 report 2026-05-06; 2 archived source references；证据质量=中高；状态=PASS。 \\
002192 & 融捷股份 & cycle-adjusted deducted-EPS PE；H1 deducted EPS bridged with H2/H1 conversion & 周期调整EPS = H1扣非利润中值 / 股本；情景EPS = H1扣非EPS×(1+H2/H1转换比例)；情景价值 = 情景EPS×业务匹配PE；熊值 = min(原始熊值, 当前价×75\%) 熊/基准/牛=47.90/85.76/132.88元；概率值=83.83元；目标=83.83元；空间=26.2\%。 & official H1 preview; Q1 financial packet; quote/adjusted K-line; 东吴证券 report 2022-10-26; 2 archived source references；证据质量=中；状态=PASS。 \\
601066 & 中信建投 & PB-ROE and earnings blended range；Q1 BPS plus H1 deducted EPS bridged to H2 & scenario value = 65\%×(Q1 BPS×PB multiple) + 35\%×(scenario EPS×PE multiple) 熊/基准/牛=10.20/15.32/20.83元；概率值=14.89元；目标=14.89元；空间=-46.5\%。 & official H1 preview; Q1 financial packet; quote/adjusted K-line; 东吴证券 report 2026-04-29; 2 archived source references；证据质量=中高；状态=PASS。 \\
002812 & 恩捷股份 & forward EPS / PE with capacity-cycle normalization；2027E normalized EPS after separator price and utilization recovery & fair value = 21.60×0.30 + 103.20×0.50 + 126.20×0.20 熊/基准/牛=21.60/103.20/126.20元；概率值=83.32元；目标=85.29元；空间=55.2\%。 & official H1 preview; Q1 financial packet; quote/adjusted K-line; 东吴证券 report 2026-07-10; 1 archived source references；证据质量=高；状态=PASS。 \\
002636 & 金安国纪 & growth-earnings scenario PE；H1 deducted EPS bridged with H2/H1 conversion & 成长盈利情景EPS = H1扣非利润中值 / 股本；情景EPS = H1扣非EPS×(1+H2/H1转换比例)；情景价值 = 情景EPS×业务匹配PE；熊值 = min(原始熊值, 当前价×75\%) 熊/基准/牛=32.81/55.75/83.84元；概率值=54.49元；目标=54.49元；空间=-42.5\%。 & official H1 preview; Q1 financial packet; quote/adjusted K-line; 开源证券 report 2021-09-02; 2 archived source references；证据质量=中；状态=PASS。 \\
002378 & 章源钨业 & cycle-adjusted deducted-EPS PE；H1 deducted EPS bridged with H2/H1 conversion & 周期调整EPS = H1扣非利润中值 / 股本；情景EPS = H1扣非EPS×(1+H2/H1转换比例)；情景价值 = 情景EPS×业务匹配PE；熊值 = min(原始熊值, 当前价×75\%) 熊/基准/牛=7.02/12.57/19.47元；概率值=12.29元；目标=12.29元；空间=-52.4\%。 & official H1 preview; Q1 financial packet; quote/adjusted K-line; 民生证券 report 2025-08-27; 2 archived source references；证据质量=中；状态=PASS。 \\
601211 & 国泰海通 & PB-ROE and earnings blended range；Q1 BPS plus H1 deducted EPS bridged to H2 & scenario value = 65\%×(Q1 BPS×PB multiple) + 35\%×(scenario EPS×PE multiple) 熊/基准/牛=13.93/21.16/28.45元；概率值=20.45元；目标=20.45元；空间=10.1\%。 & official H1 preview; Q1 financial packet; quote/adjusted K-line; 东吴证券 report 2026-07-05; 2 archived source references；证据质量=中高；状态=PASS。 \\
600884 & 杉杉股份 & growth-earnings scenario PE；H1 deducted EPS bridged with H2/H1 conversion & 成长盈利情景EPS = H1扣非利润中值 / 股本；情景EPS = H1扣非EPS×(1+H2/H1转换比例)；情景价值 = 情景EPS×业务匹配PE；熊值 = min(原始熊值, 当前价×75\%) 熊/基准/牛=9.11/15.16/22.19元；概率值=14.75元；目标=14.75元；空间=21.5\%。 & official H1 preview; Q1 financial packet; quote/adjusted K-line; 民生证券 report 2024-04-28; 2 archived source references；证据质量=中；状态=PASS。 \\
002842 & 翔鹭钨业 & cycle-adjusted deducted-EPS PE；H1 deducted EPS bridged with H2/H1 conversion & 周期调整EPS = H1扣非利润中值 / 股本；情景EPS = H1扣非EPS×(1+H2/H1转换比例)；情景价值 = 情景EPS×业务匹配PE；熊值 = min(原始熊值, 当前价×75\%) 熊/基准/牛=20.28/36.30/56.25元；概率值=35.48元；目标=35.48元；空间=-1.3\%。 & official H1 preview; Q1 financial packet; quote/adjusted K-line; 中邮证券 report 2018-12-04; 2 archived source references；证据质量=中；状态=PASS。 \\
002407 & 多氟多 & cycle-adjusted deducted-EPS PE；H1 deducted EPS bridged with H2/H1 conversion & 周期调整EPS = H1扣非利润中值 / 股本；情景EPS = H1扣非EPS×(1+H2/H1转换比例)；情景价值 = 情景EPS×业务匹配PE；熊值 = min(原始熊值, 当前价×75\%) 熊/基准/牛=6.37/10.64/15.89元；概率值=10.41元；目标=10.41元；空间=-69.3\%。 & official H1 preview; Q1 financial packet; quote/adjusted K-line; 民生证券 report 2025-04-24; 2 archived source references；证据质量=中；状态=PASS。 \\
002653 & 海思科 & growth-earnings scenario PE；H1 deducted EPS bridged with H2/H1 conversion & 成长盈利情景EPS = H1扣非利润中值 / 股本；情景EPS = H1扣非EPS×(1+H2/H1转换比例)；情景价值 = 情景EPS×业务匹配PE；熊值 = min(原始熊值, 当前价×75\%) 熊/基准/牛=19.95/32.59/47.97元；概率值=31.87元；目标=31.87元；空间=-57.0\%。 & official H1 preview; Q1 financial packet; quote/adjusted K-line; 太平洋 report 2026-06-12; 2 archived source references；证据质量=中高；状态=PASS。 \\
603629 & 利通电子 & growth-earnings scenario PE；H1 deducted EPS bridged with H2/H1 conversion & 成长盈利情景EPS = H1扣非利润中值 / 股本；情景EPS = H1扣非EPS×(1+H2/H1转换比例)；情景价值 = 情景EPS×业务匹配PE；熊值 = min(原始熊值, 当前价×75\%) 熊/基准/牛=51.02/86.71/130.39元；概率值=84.74元；目标=84.74元；空间=-24.4\%。 & official H1 preview; Q1 financial packet; quote/adjusted K-line; 东莞证券 report 2018-12-04; 2 archived source references；证据质量=中；状态=PASS。 \\
600961 & 株冶集团 & cycle-adjusted deducted-EPS PE；H1 deducted EPS bridged with H2/H1 conversion & 周期调整EPS = H1扣非利润中值 / 股本；情景EPS = H1扣非EPS×(1+H2/H1转换比例)；情景价值 = 情景EPS×业务匹配PE；熊值 = min(原始熊值, 当前价×75\%) 熊/基准/牛=18.57/39.52/61.24元；概率值=37.58元；目标=37.58元；空间=51.8\%。 & official H1 preview; Q1 financial packet; quote/adjusted K-line; 国信证券 report 2026-04-21; 2 archived source references；证据质量=中高；状态=PASS。 \\
002128 & 电投能源 & cycle-adjusted deducted-EPS PE；H1 deducted EPS bridged with H2/H1 conversion & 周期调整EPS = H1扣非利润中值 / 股本；情景EPS = H1扣非EPS×(1+H2/H1转换比例)；情景价值 = 情景EPS×业务匹配PE；熊值 = min(原始熊值, 当前价×75\%) 熊/基准/牛=14.87/25.47/33.91元；概率值=23.98元；目标=23.98元；空间=-9.5\%。 & official H1 preview; Q1 financial packet; quote/adjusted K-line; 开源证券 report 2026-04-17; 2 archived source references；证据质量=中高；状态=PASS。 \\
601168 & 西部矿业 & cycle-adjusted deducted-EPS PE；H1 deducted EPS bridged with H2/H1 conversion & 周期调整EPS = H1扣非利润中值 / 股本；情景EPS = H1扣非EPS×(1+H2/H1转换比例)；情景价值 = 情景EPS×业务匹配PE；熊值 = min(原始熊值, 当前价×75\%) 熊/基准/牛=21.80/39.03/60.49元；概率值=38.15元；目标=38.19元；空间=15.1\%。 & official H1 preview; Q1 financial packet; quote/adjusted K-line; 国信证券 report 2026-07-06; 2 archived source references；证据质量=高；状态=PASS。 \\
600188 & 兖矿能源 & cycle-adjusted deducted-EPS PE；H1 deducted EPS bridged with H2/H1 conversion & 周期调整EPS = H1扣非利润中值 / 股本；情景EPS = H1扣非EPS×(1+H2/H1转换比例)；情景价值 = 情景EPS×业务匹配PE；熊值 = min(原始熊值, 当前价×75\%) 熊/基准/牛=4.86/11.10/14.76元；概率值=9.96元；目标=9.96元；空间=-50.8\%。 & official H1 preview; Q1 financial packet; quote/adjusted K-line; 开源证券 report 2026-05-05; 2 archived source references；证据质量=中高；状态=PASS。 \\
603162 & 海通发展 & cycle-adjusted deducted-EPS PE；H1 deducted EPS bridged with H2/H1 conversion & 周期调整EPS = H1扣非利润中值 / 股本；情景EPS = H1扣非EPS×(1+H2/H1转换比例)；情景价值 = 情景EPS×业务匹配PE；熊值 = min(原始熊值, 当前价×75\%) 熊/基准/牛=4.71/9.00/13.08元；概率值=8.53元；目标=8.53元；空间=-24.1\%。 & official H1 preview; Q1 financial packet; quote/adjusted K-line; 西南证券 report 2026-07-09; 2 archived source references；证据质量=中高；状态=PASS。 \\
002240 & 盛新锂能 & cycle-adjusted EPS / PE；2026E cycle-adjusted EPS with lithium-price and cash-flow bridge & fair value = 21.60×0.30 + 39.20×0.50 + 59.04×0.20 熊/基准/牛=21.60/39.20/59.04元；概率值=37.89元；目标=37.89元；空间=21.8\%。 & official H1 preview; Q1 financial packet; quote/adjusted K-line; 东吴证券 report 2026-07-09; 1 archived source references；证据质量=中高；状态=PASS。 \\
002440 & 闰土股份 & cycle-adjusted deducted-EPS PE；H1 deducted EPS bridged with H2/H1 conversion & 周期调整EPS = H1扣非利润中值 / 股本；情景EPS = H1扣非EPS×(1+H2/H1转换比例)；情景价值 = 情景EPS×业务匹配PE；熊值 = min(原始熊值, 当前价×75\%) 熊/基准/牛=6.07/10.14/15.15元；概率值=9.92元；目标=9.92元；空间=-12.1\%。 & official H1 preview; Q1 financial packet; quote/adjusted K-line; 华鑫证券 report 2025-05-02; 2 archived source references；证据质量=中；状态=PASS。 \\
002532 & 天山铝业 & cycle-adjusted deducted-EPS PE；H1 deducted EPS bridged with H2/H1 conversion & 周期调整EPS = H1扣非利润中值 / 股本；情景EPS = H1扣非EPS×(1+H2/H1转换比例)；情景价值 = 情景EPS×业务匹配PE；熊值 = min(原始熊值, 当前价×75\%) 熊/基准/牛=9.56/23.03/30.40元；概率值=20.46元；目标=20.46元；空间=60.5\%。 & official H1 preview; Q1 financial packet; quote/adjusted K-line; 国金证券 report 2026-04-29; 2 archived source references；证据质量=中高；状态=PASS。 \\
301377 & 鼎泰高科 & deducted-EPS scenario PE；H1 deducted EPS bridged with H2/H1 conversion & 扣非EPS = H1扣非利润中值 / 股本；情景EPS = H1扣非EPS×(1+H2/H1转换比例)；情景价值 = 情景EPS×业务匹配PE；熊值 = min(原始熊值, 当前价×75\%) 熊/基准/牛=35.95/59.14/95.88元；概率值=59.53元；目标=59.53元；空间=-87.3\%。 & official H1 preview; Q1 financial packet; quote/adjusted K-line; 东吴证券 report 2026-04-16; 2 archived source references；证据质量=中高；状态=PASS。 \\
688257 & 新锐股份 & deducted-EPS scenario PE；H1 deducted EPS bridged with H2/H1 conversion & 扣非EPS = H1扣非利润中值 / 股本；情景EPS = H1扣非EPS×(1+H2/H1转换比例)；情景价值 = 情景EPS×业务匹配PE；熊值 = min(原始熊值, 当前价×75\%) 熊/基准/牛=38.36/63.11/102.30元；概率值=63.52元；目标=63.52元；空间=-20.7\%。 & official H1 preview; Q1 financial packet; quote/adjusted K-line; 东吴证券 report 2026-07-15; 2 archived source references；证据质量=中高；状态=PASS。 \\
000408 & 藏格矿业 & cycle-adjusted deducted-EPS PE；H1 deducted EPS bridged with H2/H1 conversion & 周期调整EPS = H1扣非利润中值 / 股本；情景EPS = H1扣非EPS×(1+H2/H1转换比例)；情景价值 = 情景EPS×业务匹配PE；熊值 = min(原始熊值, 当前价×75\%) 熊/基准/牛=36.68/71.68/91.59元；概率值=65.16元；目标=65.16元；空间=-8.9\%。 & official H1 preview; Q1 financial packet; quote/adjusted K-line; 东吴证券 report 2026-07-08; 2 archived source references；证据质量=中高；状态=PASS。 \\
600601 & 方正科技 & growth-earnings scenario PE；H1 deducted EPS bridged with H2/H1 conversion & 成长盈利情景EPS = H1扣非利润中值 / 股本；情景EPS = H1扣非EPS×(1+H2/H1转换比例)；情景价值 = 情景EPS×业务匹配PE；熊值 = min(原始熊值, 当前价×75\%) 熊/基准/牛=3.97/6.75/10.16元；概率值=6.60元；目标=6.60元；空间=-43.6\%。 & official H1 preview; Q1 financial packet; quote/adjusted K-line; 中航证券 report 2025-07-24; 2 archived source references；证据质量=中；状态=PASS。 \\
002475 & 立讯精密 & growth-earnings scenario PE；H1 deducted EPS bridged with H2/H1 conversion & 成长盈利情景EPS = H1扣非利润中值 / 股本；情景EPS = H1扣非EPS×(1+H2/H1转换比例)；情景价值 = 情景EPS×业务匹配PE；熊值 = min(原始熊值, 当前价×75\%) 熊/基准/牛=25.71/43.68/65.69元；概率值=42.69元；目标=42.69元；空间=-29.4\%。 & official H1 preview; Q1 financial packet; quote/adjusted K-line; 招银国际 report 2025-11-05; 2 archived source references；证据质量=中；状态=PASS。 \\
600595 & 中孚实业 & cycle-adjusted deducted-EPS PE；H1 deducted EPS bridged with H2/H1 conversion & 周期调整EPS = H1扣非利润中值 / 股本；情景EPS = H1扣非EPS×(1+H2/H1转换比例)；情景价值 = 情景EPS×业务匹配PE；熊值 = min(原始熊值, 当前价×75\%) 熊/基准/牛=4.84/12.00/16.52元；概率值=10.76元；目标=10.76元；空间=66.6\%。 & official H1 preview; Q1 financial packet; quote/adjusted K-line; 国信证券 report 2026-04-28; 2 archived source references；证据质量=中高；状态=PASS。 \\
600487 & 亨通光电 & growth-earnings scenario PE；H1 deducted EPS bridged with H2/H1 conversion & 成长盈利情景EPS = H1扣非利润中值 / 股本；情景EPS = H1扣非EPS×(1+H2/H1转换比例)；情景价值 = 情景EPS×业务匹配PE；熊值 = min(原始熊值, 当前价×75\%) 熊/基准/牛=41.26/70.12/105.45元；概率值=68.53元；目标=68.53元；空间=-3.5\%。 & official H1 preview; Q1 financial packet; quote/adjusted K-line; 群益证券 report 2026-07-14; 2 archived source references；证据质量=中高；状态=PASS。 \\
600183 & 生益科技 & growth-earnings scenario PE；H1 deducted EPS bridged with H2/H1 conversion & 成长盈利情景EPS = H1扣非利润中值 / 股本；情景EPS = H1扣非EPS×(1+H2/H1转换比例)；情景价值 = 情景EPS×业务匹配PE；熊值 = min(原始熊值, 当前价×75\%) 熊/基准/牛=35.51/60.34/90.74元；概率值=58.97元；目标=58.97元；空间=-61.3\%。 & official H1 preview; Q1 financial packet; quote/adjusted K-line; 太平洋 report 2026-05-27; 2 archived source references；证据质量=中高；状态=PASS。 \\
000630 & 铜陵有色 & cycle-adjusted deducted-EPS PE；H1 deducted EPS bridged with H2/H1 conversion & 周期调整EPS = H1扣非利润中值 / 股本；情景EPS = H1扣非EPS×(1+H2/H1转换比例)；情景价值 = 情景EPS×业务匹配PE；熊值 = min(原始熊值, 当前价×75\%) 熊/基准/牛=2.64/6.72/7.31元；概率值=5.61元；目标=5.61元；空间=-7.4\%。 & official H1 preview; Q1 financial packet; quote/adjusted K-line; 国信证券 report 2026-04-22; 2 archived source references；证据质量=中高；状态=PASS。 \\
600549 & 厦门钨业 & cycle-adjusted deducted-EPS PE；H1 deducted EPS bridged with H2/H1 conversion & 周期调整EPS = H1扣非利润中值 / 股本；情景EPS = H1扣非EPS×(1+H2/H1转换比例)；情景价值 = 情景EPS×业务匹配PE；熊值 = min(原始熊值, 当前价×75\%) 熊/基准/牛=17.00/35.16/47.15元；概率值=32.11元；目标=32.11元；空间=-44.7\%。 & official H1 preview; Q1 financial packet; quote/adjusted K-line; 太平洋 report 2026-06-21; 2 archived source references；证据质量=中高；状态=PASS。 \\
000100 & TCL科技 & growth-earnings scenario PE；H1 deducted EPS bridged with H2/H1 conversion & 成长盈利情景EPS = H1扣非利润中值 / 股本；情景EPS = H1扣非EPS×(1+H2/H1转换比例)；情景价值 = 情景EPS×业务匹配PE；熊值 = min(原始熊值, 当前价×75\%) 熊/基准/牛=3.76/9.36/11.39元；概率值=8.09元；目标=8.09元；空间=61.2\%。 & official H1 preview; Q1 financial packet; quote/adjusted K-line; 中邮证券 report 2026-04-27; 2 archived source references；证据质量=中高；状态=PASS。 \\
001287 & 中电港 & growth-earnings scenario PE；H1 deducted EPS bridged with H2/H1 conversion & 成长盈利情景EPS = H1扣非利润中值 / 股本；情景EPS = H1扣非EPS×(1+H2/H1转换比例)；情景价值 = 情景EPS×业务匹配PE；熊值 = min(原始熊值, 当前价×75\%) 熊/基准/牛=19.99/47.84/53.00元；概率值=40.52元；目标=40.52元；空间=52.0\%。 & official H1 preview; Q1 financial packet; quote/adjusted K-line; 爱建证券 report 2026-07-13; 2 archived source references；证据质量=中高；状态=PASS。 \\
600111 & 北方稀土 & cycle-adjusted deducted-EPS PE；H1 deducted EPS bridged with H2/H1 conversion & 周期调整EPS = H1扣非利润中值 / 股本；情景EPS = H1扣非EPS×(1+H2/H1转换比例)；情景价值 = 情景EPS×业务匹配PE；熊值 = min(原始熊值, 当前价×75\%) 熊/基准/牛=6.96/13.44/19.32元；概率值=12.67元；目标=12.67元；空间=-68.0\%。 & official H1 preview; Q1 financial packet; quote/adjusted K-line; 东莞证券 report 2026-04-23; 2 archived source references；证据质量=中高；状态=PASS。 \\
000725 & 京东方A & growth-earnings scenario PE；H1 deducted EPS bridged with H2/H1 conversion & 成长盈利情景EPS = H1扣非利润中值 / 股本；情景EPS = H1扣非EPS×(1+H2/H1转换比例)；情景价值 = 情景EPS×业务匹配PE；熊值 = min(原始熊值, 当前价×75\%) 熊/基准/牛=2.64/6.74/6.76元；概率值=5.51元；目标=5.51元；空间=-13.6\%。 & official H1 preview; Q1 financial packet; quote/adjusted K-line; 金元证券 report 2026-07-10; 2 archived source references；证据质量=中高；状态=PASS。 \\
600176 & 中国巨石 & deducted-EPS scenario PE；H1 deducted EPS bridged with H2/H1 conversion & 扣非EPS = H1扣非利润中值 / 股本；情景EPS = H1扣非EPS×(1+H2/H1转换比例)；情景价值 = 情景EPS×业务匹配PE；熊值 = min(原始熊值, 当前价×75\%) 熊/基准/牛=14.73/30.06/39.27元；概率值=27.30元；目标=27.30元；空间=-48.5\%。 & official H1 preview; Q1 financial packet; quote/adjusted K-line; 华源证券 report 2026-05-27; 2 archived source references；证据质量=中高；状态=PASS。 \\
002266 & 浙富控股 & deducted-EPS scenario PE；H1 deducted EPS bridged with H2/H1 conversion & 扣非EPS = H1扣非利润中值 / 股本；情景EPS = H1扣非EPS×(1+H2/H1转换比例)；情景价值 = 情景EPS×业务匹配PE；熊值 = min(原始熊值, 当前价×75\%) 熊/基准/牛=3.82/8.19/12.65元；概率值=7.77元；目标=7.77元；空间=52.6\%。 & official H1 preview; Q1 financial packet; quote/adjusted K-line; 中邮证券 report 2026-06-29; 2 archived source references；证据质量=中高；状态=PASS。 \\
000612 & 焦作万方 & cycle-adjusted deducted-EPS PE；H1 deducted EPS bridged with H2/H1 conversion & 周期调整EPS = H1扣非利润中值 / 股本；情景EPS = H1扣非EPS×(1+H2/H1转换比例)；情景价值 = 情景EPS×业务匹配PE；熊值 = min(原始熊值, 当前价×75\%) 熊/基准/牛=8.13/22.50/34.86元；概率值=20.66元；目标=20.66元；空间=90.6\%。 & official H1 preview; Q1 financial packet; quote/adjusted K-line; 开源证券 report 2025-09-26; 2 archived source references；证据质量=中；状态=PASS。 \\
605117 & 德业股份 & growth-earnings scenario PE；H1 deducted EPS bridged with H2/H1 conversion & 成长盈利情景EPS = H1扣非利润中值 / 股本；情景EPS = H1扣非EPS×(1+H2/H1转换比例)；情景价值 = 情景EPS×业务匹配PE；熊值 = min(原始熊值, 当前价×75\%) 熊/基准/牛=53.90/111.76/127.61元；概率值=97.57元；目标=97.57元；空间=14.2\%。 & official H1 preview; Q1 financial packet; quote/adjusted K-line; 国信证券 report 2026-07-10; 2 archived source references；证据质量=中高；状态=PASS。 \\
000783 & 长江证券 & PB-ROE and earnings blended range；Q1 BPS plus H1 deducted EPS bridged to H2 & scenario value = 65\%×(Q1 BPS×PB multiple) + 35\%×(scenario EPS×PE multiple) 熊/基准/牛=5.99/8.98/12.19元；概率值=8.72元；目标=8.72元；空间=-3.8\%。 & official H1 preview; Q1 financial packet; quote/adjusted K-line; 中航证券 report 2021-05-06; 2 archived source references；证据质量=中；状态=PASS。 \\
600233 & 圆通速递 & cycle-adjusted deducted-EPS PE；H1 deducted EPS bridged with H2/H1 conversion & 周期调整EPS = H1扣非利润中值 / 股本；情景EPS = H1扣非EPS×(1+H2/H1转换比例)；情景价值 = 情景EPS×业务匹配PE；熊值 = min(原始熊值, 当前价×75\%) 熊/基准/牛=11.56/21.96/32.06元；概率值=20.86元；目标=20.86元；空间=19.2\%。 & official H1 preview; Q1 financial packet; quote/adjusted K-line; 群益证券 report 2026-07-01; 2 archived source references；证据质量=中高；状态=PASS。 \\
600392 & 盛和资源 & cycle-adjusted deducted-EPS PE；H1 deducted EPS bridged with H2/H1 conversion & 周期调整EPS = H1扣非利润中值 / 股本；情景EPS = H1扣非EPS×(1+H2/H1转换比例)；情景价值 = 情景EPS×业务匹配PE；熊值 = min(原始熊值, 当前价×75\%) 熊/基准/牛=6.05/10.83/16.78元；概率值=10.59元；目标=10.59元；空间=-55.6\%。 & official H1 preview; Q1 financial packet; quote/adjusted K-line; 国元证券 report 2025-11-12; 2 archived source references；证据质量=中；状态=PASS。 \\
603225 & 新凤鸣 & cycle-adjusted deducted-EPS PE；H1 deducted EPS bridged with H2/H1 conversion & 周期调整EPS = H1扣非利润中值 / 股本；情景EPS = H1扣非EPS×(1+H2/H1转换比例)；情景价值 = 情景EPS×业务匹配PE；熊值 = min(原始熊值, 当前价×75\%) 熊/基准/牛=13.48/23.39/34.95元；概率值=22.73元；目标=22.73元；空间=26.5\%。 & official H1 preview; Q1 financial packet; quote/adjusted K-line; 中邮证券 report 2026-07-10; 2 archived source references；证据质量=中高；状态=PASS。 \\
002008 & 大族激光 & deducted-EPS scenario PE；H1 deducted EPS bridged with H2/H1 conversion & 扣非EPS = H1扣非利润中值 / 股本；情景EPS = H1扣非EPS×(1+H2/H1转换比例)；情景价值 = 情景EPS×业务匹配PE；熊值 = min(原始熊值, 当前价×75\%) 熊/基准/牛=31.41/52.40/83.76元；概率值=52.38元；目标=52.38元；空间=-54.2\%。 & official H1 preview; Q1 financial packet; quote/adjusted K-line; 东吴证券 report 2026-07-10; 2 archived source references；证据质量=中高；状态=PASS。 \\
000685 & 中山公用 & deducted-EPS scenario PE；H1 deducted EPS bridged with H2/H1 conversion & 扣非EPS = H1扣非利润中值 / 股本；情景EPS = H1扣非EPS×(1+H2/H1转换比例)；情景价值 = 情景EPS×业务匹配PE；熊值 = min(原始熊值, 当前价×75\%) 熊/基准/牛=8.39/33.26/51.34元；概率值=29.41元；目标=29.41元；空间=162.8\%。 & official H1 preview; Q1 financial packet; quote/adjusted K-line; 国信证券 report 2026-05-11; 2 archived source references；证据质量=中高；状态=PASS。 \\
600909 & 华安证券 & PB-ROE and earnings blended range；Q1 BPS plus H1 deducted EPS bridged to H2 & scenario value = 65\%×(Q1 BPS×PB multiple) + 35\%×(scenario EPS×PE multiple) 熊/基准/牛=4.61/6.89/9.35元；概率值=6.70元；目标=6.70元；空间=-29.5\%。 & official H1 preview; Q1 financial packet; quote/adjusted K-line; 太平洋 report 2026-04-23; 2 archived source references；证据质量=中高；状态=PASS。 \\
000833 & 粤桂股份 & deducted-EPS scenario PE；H1 deducted EPS bridged with H2/H1 conversion & 扣非EPS = H1扣非利润中值 / 股本；情景EPS = H1扣非EPS×(1+H2/H1转换比例)；情景价值 = 情景EPS×业务匹配PE；熊值 = min(原始熊值, 当前价×75\%) 熊/基准/牛=15.05/27.84/42.99元；概率值=27.03元；目标=27.03元；空间=34.7\%。 & official H1 preview; Q1 financial packet; quote/adjusted K-line; 东兴证券 report 2018-08-31; 2 archived source references；证据质量=中；状态=PASS。 \\
300037 & 新宙邦 & growth-earnings scenario PE；H1 deducted EPS bridged with H2/H1 conversion & 成长盈利情景EPS = H1扣非利润中值 / 股本；情景EPS = H1扣非EPS×(1+H2/H1转换比例)；情景价值 = 情景EPS×业务匹配PE；熊值 = min(原始熊值, 当前价×75\%) 熊/基准/牛=35.54/70.62/84.14元；概率值=62.80元；目标=68.85元；空间=2.5\%。 & official H1 preview; Q1 financial packet; quote/adjusted K-line; 东吴证券 report 2026-07-10; 2 archived source references；证据质量=高；状态=PASS。 \\
000060 & 中金岭南 & cycle-adjusted deducted-EPS PE；H1 deducted EPS bridged with H2/H1 conversion & 周期调整EPS = H1扣非利润中值 / 股本；情景EPS = H1扣非EPS×(1+H2/H1转换比例)；情景价值 = 情景EPS×业务匹配PE；熊值 = min(原始熊值, 当前价×75\%) 熊/基准/牛=3.09/6.12/8.56元；概率值=5.70元；目标=5.70元；空间=-8.9\%。 & official H1 preview; Q1 financial packet; quote/adjusted K-line; 中邮证券 report 2026-05-13; 2 archived source references；证据质量=中高；状态=PASS。 \\
600522 & 中天科技 & growth-earnings scenario PE；H1 deducted EPS bridged with H2/H1 conversion & 成长盈利情景EPS = H1扣非利润中值 / 股本；情景EPS = H1扣非EPS×(1+H2/H1转换比例)；情景价值 = 情景EPS×业务匹配PE；熊值 = min(原始熊值, 当前价×75\%) 熊/基准/牛=20.02/34.02/51.16元；概率值=33.25元；目标=33.25元；空间=-18.3\%。 & official H1 preview; Q1 financial packet; quote/adjusted K-line; 中邮证券 report 2026-02-25; 2 archived source references；证据质量=中；状态=PASS。 \\
002468 & 申通快递 & cycle-adjusted deducted-EPS PE；H1 deducted EPS bridged with H2/H1 conversion & 周期调整EPS = H1扣非利润中值 / 股本；情景EPS = H1扣非EPS×(1+H2/H1转换比例)；情景价值 = 情景EPS×业务匹配PE；熊值 = min(原始熊值, 当前价×75\%) 熊/基准/牛=8.14/14.88/22.58元；概率值=14.40元；目标=14.40元；空间=-5.4\%。 & official H1 preview; Q1 financial packet; quote/adjusted K-line; 国信证券 report 2026-04-16; 2 archived source references；证据质量=中高；状态=PASS。 \\
600267 & 海正药业 & growth-earnings scenario PE；H1 deducted EPS bridged with H2/H1 conversion & 成长盈利情景EPS = H1扣非利润中值 / 股本；情景EPS = H1扣非EPS×(1+H2/H1转换比例)；情景价值 = 情景EPS×业务匹配PE；熊值 = min(原始熊值, 当前价×75\%) 熊/基准/牛=8.98/21.60/31.80元；概率值=19.85元；目标=19.85元；空间=65.7\%。 & official H1 preview; Q1 financial packet; quote/adjusted K-line; 中银证券 report 2022-08-24; 2 archived source references；证据质量=中；状态=PASS。 \\
600906 & 财达证券 & PB-ROE and earnings blended range；Q1 BPS plus H1 deducted EPS bridged to H2 & scenario value = 65\%×(Q1 BPS×PB multiple) + 35\%×(scenario EPS×PE multiple) 熊/基准/牛=3.00/4.44/5.97元；概率值=4.31元；目标=4.31元；空间=-40.6\%。 & official H1 preview; Q1 financial packet; quote/adjusted K-line; 东兴证券 report 2022-02-24; 2 archived source references；证据质量=中；状态=PASS。 \\
603268 & 松发股份 & growth-earnings scenario PE；H1 deducted EPS bridged with H2/H1 conversion & 成长盈利情景EPS = H1扣非利润中值 / 股本；情景EPS = H1扣非EPS×(1+H2/H1转换比例)；情景价值 = 情景EPS×业务匹配PE；熊值 = min(原始熊值, 当前价×75\%) 熊/基准/牛=120.34/259.00/389.38元；概率值=243.48元；目标=243.48元；空间=51.7\%。 & official H1 preview; Q1 financial packet; quote/adjusted K-line; 东吴证券 report 2026-06-02; 2 archived source references；证据质量=中高；状态=PASS。 \\
002458 & 益生股份 & deducted-EPS scenario PE；H1 deducted EPS bridged with H2/H1 conversion & 扣非EPS = H1扣非利润中值 / 股本；情景EPS = H1扣非EPS×(1+H2/H1转换比例)；情景价值 = 情景EPS×业务匹配PE；熊值 = min(原始熊值, 当前价×75\%) 熊/基准/牛=3.92/8.82/10.44元；概率值=7.67元；目标=7.67元；空间=-15.0\%。 & official H1 preview; Q1 financial packet; quote/adjusted K-line; 太平洋 report 2026-03-30; 2 archived source references；证据质量=中；状态=PASS。 \\
002384 & 东山精密 & growth-earnings scenario PE；H1 deducted EPS bridged with H2/H1 conversion & 成长盈利情景EPS = H1扣非利润中值 / 股本；情景EPS = H1扣非EPS×(1+H2/H1转换比例)；情景价值 = 情景EPS×业务匹配PE；熊值 = min(原始熊值, 当前价×75\%) 熊/基准/牛=40.93/97.76/104.60元；概率值=82.08元；目标=82.08元；空间=-68.7\%。 & official H1 preview; Q1 financial packet; quote/adjusted K-line; 开源证券 report 2026-05-05; 2 archived source references；证据质量=中高；状态=PASS。 \\
000977 & 浪潮信息 & growth-earnings scenario PE；H1 deducted EPS bridged with H2/H1 conversion & 成长盈利情景EPS = H1扣非利润中值 / 股本；情景EPS = H1扣非EPS×(1+H2/H1转换比例)；情景价值 = 情景EPS×业务匹配PE；熊值 = min(原始熊值, 当前价×75\%) 熊/基准/牛=48.03/90.10/143.15元；概率值=88.09元；目标=88.09元；空间=4.0\%。 & official H1 preview; Q1 financial packet; quote/adjusted K-line; 中邮证券 report 2026-07-09; 2 archived source references；证据质量=中高；状态=PASS。 \\
603127 & 昭衍新药 & growth-earnings scenario PE；H1 deducted EPS bridged with H2/H1 conversion & 成长盈利情景EPS = H1扣非利润中值 / 股本；情景EPS = H1扣非EPS×(1+H2/H1转换比例)；情景价值 = 情景EPS×业务匹配PE；熊值 = min(原始熊值, 当前价×75\%) 熊/基准/牛=28.64/46.80/68.89元；概率值=45.77元；目标=45.77元；空间=-9.6\%。 & official H1 preview; Q1 financial packet; quote/adjusted K-line; 国信证券 report 2026-04-20; 2 archived source references；证据质量=中高；状态=PASS。 \\
301200 & 大族数控 & deducted-EPS scenario PE；H1 deducted EPS bridged with H2/H1 conversion & 扣非EPS = H1扣非利润中值 / 股本；情景EPS = H1扣非EPS×(1+H2/H1转换比例)；情景价值 = 情景EPS×业务匹配PE；熊值 = min(原始熊值, 当前价×75\%) 熊/基准/牛=44.88/85.00/119.67元；概率值=79.90元；目标=79.90元；空间=-74.4\%。 & official H1 preview; Q1 financial packet; quote/adjusted K-line; 东吴证券 report 2026-07-10; 2 archived source references；证据质量=中高；状态=PASS。 \\
000938 & 紫光股份 & growth-earnings scenario PE；H1 deducted EPS bridged with H2/H1 conversion & 成长盈利情景EPS = H1扣非利润中值 / 股本；情景EPS = H1扣非EPS×(1+H2/H1转换比例)；情景价值 = 情景EPS×业务匹配PE；熊值 = min(原始熊值, 当前价×75\%) 熊/基准/牛=19.42/42.28/57.88元；概率值=38.54元；目标=38.54元；空间=8.5\%。 & official H1 preview; Q1 financial packet; quote/adjusted K-line; 开源证券 report 2026-07-13; 2 archived source references；证据质量=中高；状态=PASS。 \\
002463 & 沪电股份 & growth-earnings scenario PE；H1 deducted EPS bridged with H2/H1 conversion & 成长盈利情景EPS = H1扣非利润中值 / 股本；情景EPS = H1扣非EPS×(1+H2/H1转换比例)；情景价值 = 情景EPS×业务匹配PE；熊值 = min(原始熊值, 当前价×75\%) 熊/基准/牛=44.60/77.74/113.99元；概率值=75.05元；目标=75.05元；空间=-45.3\%。 & official H1 preview; Q1 financial packet; quote/adjusted K-line; 中邮证券 report 2026-04-24; 2 archived source references；证据质量=中高；状态=PASS。 \\
688183 & 生益电子 & growth-earnings scenario PE；H1 deducted EPS bridged with H2/H1 conversion & 成长盈利情景EPS = H1扣非利润中值 / 股本；情景EPS = H1扣非EPS×(1+H2/H1转换比例)；情景价值 = 情景EPS×业务匹配PE；熊值 = min(原始熊值, 当前价×75\%) 熊/基准/牛=40.52/68.86/103.56元；概率值=67.30元；目标=67.30元；空间=-47.8\%。 & official H1 preview; Q1 financial packet; quote/adjusted K-line; 华安证券 report 2025-12-31; 2 archived source references；证据质量=中；状态=PASS。 \\
300604 & 长川科技 & growth-earnings scenario PE；H1 deducted EPS bridged with H2/H1 conversion & 成长盈利情景EPS = H1扣非利润中值 / 股本；情景EPS = H1扣非EPS×(1+H2/H1转换比例)；情景价值 = 情景EPS×业务匹配PE；熊值 = min(原始熊值, 当前价×75\%) 熊/基准/牛=43.65/86.06/111.55元；概率值=78.44元；目标=108.60元；空间=-64.3\%。 & official H1 preview; Q1 financial packet; quote/adjusted K-line; 群益证券 report 2026-06-26; 2 archived source references；证据质量=高；状态=PASS。 \\
002916 & 深南电路 & growth-earnings scenario PE；H1 deducted EPS bridged with H2/H1 conversion & 成长盈利情景EPS = H1扣非利润中值 / 股本；情景EPS = H1扣非EPS×(1+H2/H1转换比例)；情景价值 = 情景EPS×业务匹配PE；熊值 = min(原始熊值, 当前价×75\%) 熊/基准/牛=97.93/211.64/250.27元；概率值=185.25元；目标=185.25元；空间=-50.9\%。 & official H1 preview; Q1 financial packet; quote/adjusted K-line; 太平洋 report 2026-03-26; 2 archived source references；证据质量=中；状态=PASS。 \\
688825 & 长鑫科技 & 上市边界，不作二级市场估值；无二级市场当前价 & 目标价/空间 = 不适用；待上市交易后重新建立当前价模型 目标价/空间不适用；上市交易后重建模型。 & IPO发行价公告；截至数据截点无二级市场交易价格；证据质量=低；状态=PASS。 \\
001248 & 华润新能源 & deducted-EPS scenario PE；H1 deducted EPS bridged with H2/H1 conversion & 扣非EPS = H1扣非利润中值 / 股本；情景EPS = H1扣非EPS×(1+H2/H1转换比例)；情景价值 = 情景EPS×业务匹配PE；熊值 = min(原始熊值, 当前价×75\%) 熊/基准/牛=4.04/6.72/10.18元；概率值=6.61元；目标=6.61元；空间=-50.7\%。 & official H1 preview; Q1 financial packet; quote/adjusted K-line; 华金证券 report 2026-06-21; 2 archived source references；证据质量=中高；状态=PASS。 \\
601633 & 长城汽车 & deducted-EPS scenario PE；H1 deducted EPS bridged with H2/H1 conversion & 扣非EPS = H1扣非利润中值 / 股本；情景EPS = H1扣非EPS×(1+H2/H1转换比例)；情景价值 = 情景EPS×业务匹配PE；熊值 = min(原始熊值, 当前价×75\%) 熊/基准/牛=3.76/8.36/16.99元；概率值=8.71元；目标=8.71元；空间=-44.1\%。 & official H1 preview; Q1 financial packet; quote/adjusted K-line; 国金证券 report 2026-07-15; 2 archived source references；证据质量=中高；状态=PASS。 \\
600736 & 苏州高新 & normalized EPS / PE with PB cross-check；2026E normalized EPS after property settlement and investment-income normalization & fair value = 4.20×0.30 + 7.70×0.50 + 10.40×0.20 熊/基准/牛=4.20/7.70/10.40元；概率值=7.19元；目标=7.19元；空间=4.7\%。 & official H1 preview; Q1 financial packet; quote/adjusted K-line; 东吴证券 report 2026-04-28; 1 archived source references；证据质量=中高；状态=PASS。 \\
000415 & 渤海租赁 & normalized EPS / PE with PB cross-check；2026E normalized EPS after lease income, asset quality and FX normalization & fair value = 3.00×0.30 + 4.62×0.50 + 5.95×0.20 熊/基准/牛=3.00/4.62/5.95元；概率值=4.40元；目标=4.40元；空间=-1.1\%。 & official H1 preview; Q1 financial packet; quote/adjusted K-line; 西南证券 report 2026-07-09; 1 archived source references；证据质量=中高；状态=PASS。 \\
002221 & 东华能源 & cycle-adjusted deducted-EPS PE；H1 deducted EPS bridged with H2/H1 conversion & 周期调整EPS = H1扣非利润中值 / 股本；情景EPS = H1扣非EPS×(1+H2/H1转换比例)；情景价值 = 情景EPS×业务匹配PE；熊值 = min(原始熊值, 当前价×75\%) 熊/基准/牛=0.56/1.03/1.63元；概率值=1.01元；目标=1.01元；空间=-81.4\%。 & official H1 preview; Q1 financial packet; quote/adjusted K-line; 东吴证券 report 2024-12-30; 2 archived source references；证据质量=中；状态=PASS。 \\
002611 & 东方精工 & deducted-EPS scenario PE；H1 deducted EPS bridged with H2/H1 conversion & 扣非EPS = H1扣非利润中值 / 股本；情景EPS = H1扣非EPS×(1+H2/H1转换比例)；情景价值 = 情景EPS×业务匹配PE；熊值 = min(原始熊值, 当前价×75\%) 熊/基准/牛=1.71/2.81/4.55元；概率值=2.83元；目标=2.83元；空间=-81.2\%。 & official H1 preview; Q1 financial packet; quote/adjusted K-line; 爱建证券 report 2025-12-25; 2 archived source references；证据质量=中；状态=PASS。 \\
002174 & 游族网络 & growth-earnings scenario PE；H1 deducted EPS bridged with H2/H1 conversion & 成长盈利情景EPS = H1扣非利润中值 / 股本；情景EPS = H1扣非EPS×(1+H2/H1转换比例)；情景价值 = 情景EPS×业务匹配PE；熊值 = min(原始熊值, 当前价×75\%) 熊/基准/牛=0.35/0.59/0.88元；概率值=0.58元；目标=0.58元；空间=-95.4\%。 & official H1 preview; Q1 financial packet; quote/adjusted K-line; 东吴证券 report 2020-08-28; 2 archived source references；证据质量=中；状态=PASS。 \\
603296 & 华勤技术 & growth-earnings scenario PE；H1 deducted EPS bridged with H2/H1 conversion & 成长盈利情景EPS = H1扣非利润中值 / 股本；情景EPS = H1扣非EPS×(1+H2/H1转换比例)；情景价值 = 情景EPS×业务匹配PE；熊值 = min(原始熊值, 当前价×75\%) 熊/基准/牛=33.96/86.78/122.72元；概率值=78.12元；目标=78.12元；空间=-0.1\%。 & official H1 preview; Q1 financial packet; quote/adjusted K-line; 太平洋 report 2026-05-22; 2 archived source references；证据质量=中高；状态=PASS。 \\
\bottomrule
\end{longtable}
\normalsize

逐票审计包的完整证据路径、原始研报文件、输入字段和重算结果保存在\texttt{refresh-20260715/data/full\_market\_candidate\_valuation\_audit\_20260715.json}；正文表格使用短证据索引，避免把原始路径压缩成不可读的横向长串。
\section{估值推演与情景分析}

下列每个标的均使用真实LaTeX数学环境展示代入步骤。表格是阅读索引，以下公式块是可复核的计算正文；结构化审计包保存未四舍五入的输入字段和证据路径。

\scriptsize
\paragraph{000042 中洲控股：项目结算正常化 EPS / PE}
分母：recurring/normalized EPS after project settlement and tax/interest normalization；证据质量：中；审计状态：PASS。
为什么选这个方法：最新H1利润存在低基数、扭亏或周期扭曲，因此先正常化EPS，再套业务匹配PE区间；业务匹配可比包括保利发展、招商蛇口、滨江集团、越秀地产。
为什么选这个分母：正常化分母为：recurring/normalized EPS after project settlement and tax/interest normalization；H1机械年化指标只作为市场隐含压力测试，不作为最终目标价分母。
为什么用这组倍数：熊/基准/牛倍数分别表达下行、正常兑现和上行重估；区间通过可比方法和当前价格隐含指标交叉检查（业绩预告刺激有限，7/15主力净流出约1473万元、连续2日减仓；地产高负债和项目结算风险压制估值），因此只能作为条件性House价值。
为什么用这组概率：30\%/50\%/20\%是统一筛选先验：基准情景权重最高，同时保留熊市和牛市的实质性下行/上行贡献；它是House情景权重，不是经过统计估计的客观事件概率。
证据如何参与：官方半年报预告、Q1财务包和行情/K线提供事实输入；研报/来源用于分母或可比背景。外部目标价11.84元，仅作背景，权重为0。
\begin{align*}
V_s=\mathrm{EPS}_{normalized,s}\times \mathrm{PE}_s\\
V_{b/m/b} &= 4.80/8.80/13.50\ \mathrm{CNY/share}\\
V_{prob} &= 0.30\times 4.80+0.50\times 8.80+0.20\times 13.50=8.54\\
P_{target} &= V_{prob}=8.54\quad(\text{外部锚权重}=0)\\
\mathrm{Upside} &= \frac{8.54}{8.02}-1=6.5\%\\
\end{align*}

\paragraph{600150 中国船舶：成长盈利情景 PE}
分母：H1 deducted EPS bridged with H2/H1 conversion；证据质量：中高；审计状态：PASS。
为什么选这个方法：国防军工具有成长或盈利持续期特征，因此模型给盈利桥接配置较高PE区间，而不是直接给主题收入估值；价格和现金流检查防止把泛化成长叙事直接变成EPS。
为什么选这个分母：使用H1扣非利润以减少非经常性收益污染；股本由总市值/当前价反推；H2/H1转换比例是明确的季节性和交付假设，不是机械年化。
为什么用这组倍数：PE区间是House对成长持续期的情景判断；牛市情景要求业绩桥接和经营证据持续，不是市场一致目标价。
为什么用这组概率：30\%/50\%/20\%是统一筛选先验：基准情景权重最高，同时保留熊市和牛市的实质性下行/上行贡献；它是House情景权重，不是经过统计估计的客观事件概率。
证据如何参与：官方半年报预告、Q1财务包和行情/K线提供事实输入；研报/来源用于分母或可比背景。没有正权重外部目标锚，House价值不伪装为Street共识。
\begin{align*}
\mathrm{Shares} &= \frac{2582.04}{34.31}=75.2562\ \mathrm{100m\ shares}\\
\mathrm{EPS}_{H1}^d &= \frac{99.00}{75.2562}=1.3155\\
\mathrm{EPS}_s &= \mathrm{EPS}_{H1}^d(1+r_s)=2.24/2.70/3.16\\
V_s &= \mathrm{EPS}_s\times \mathrm{PE}_s=25.73/94.39/142.07\\
V_{prob} &= 0.30\times 25.73+0.50\times 94.39+0.20\times 142.07=83.33\\
P_{target} &= V_{prob}=83.33\quad(\text{外部锚权重}=0)\\
\mathrm{Upside} &= \frac{83.33}{34.31}-1=142.9\%\\
\end{align*}

\paragraph{002379 宏桥控股：周期调整扣非 EPS / PE}
分母：H1 deducted EPS bridged with H2/H1 conversion；证据质量：高；审计状态：PASS。
为什么选这个方法：有色金属盈利受商品价格、运价、价差或利用率周期影响，因此使用周期调整EPS/PE区间，而不是把某一个季度的PE当作全年估值。H2转换比例用于检验当前周期能否延续。
为什么选这个分母：使用H1扣非利润以减少非经常性收益污染；股本由总市值/当前价反推；H2/H1转换比例是明确的季节性和交付假设，不是机械年化。
为什么用这组倍数：熊/基准/牛PE区间是针对该周期业务类型的House假设，不是外部券商目标价；熊值较低反映周期反转，牛值较高必须由价格、销量和现金流共同确认。
为什么用这组概率：30\%/50\%/20\%是统一筛选先验：基准情景权重最高，同时保留熊市和牛市的实质性下行/上行贡献；它是House情景权重，不是经过统计估计的客观事件概率。
证据如何参与：官方半年报预告、Q1财务包和行情/K线提供事实输入；研报/来源用于分母或可比背景。外部目标价29.00元，赋予正权重10\%。
\begin{align*}
\mathrm{Shares} &= \frac{2431.61}{18.66}=130.3114\ \mathrm{100m\ shares}\\
\mathrm{EPS}_{H1}^d &= \frac{154.00}{130.3114}=1.1818\\
\mathrm{EPS}_s &= \mathrm{EPS}_{H1}^d(1+r_s)=1.83/2.19/2.54\\
V_s &= \mathrm{EPS}_s\times \mathrm{PE}_s=14.00/26.24/40.65\\
V_{prob} &= 0.30\times 14.00+0.50\times 26.24+0.20\times 40.65=25.45\\
P_{target} &= 25.45\times 0.90+29.00\times 0.10=25.80\\
\mathrm{Upside} &= \frac{25.80}{18.66}-1=38.3\%\\
\end{align*}

\paragraph{601600 中国铝业：周期调整扣非 EPS / PE}
分母：H1 deducted EPS bridged with H2/H1 conversion；证据质量：中高；审计状态：PASS。
为什么选这个方法：有色金属盈利受商品价格、运价、价差或利用率周期影响，因此使用周期调整EPS/PE区间，而不是把某一个季度的PE当作全年估值。H2转换比例用于检验当前周期能否延续。
为什么选这个分母：使用H1扣非利润以减少非经常性收益污染；股本由总市值/当前价反推；H2/H1转换比例是明确的季节性和交付假设，不是机械年化。
为什么用这组倍数：熊/基准/牛PE区间是针对该周期业务类型的House假设，不是外部券商目标价；熊值较低反映周期反转，牛值较高必须由价格、销量和现金流共同确认。
为什么用这组概率：30\%/50\%/20\%是统一筛选先验：基准情景权重最高，同时保留熊市和牛市的实质性下行/上行贡献；它是House情景权重，不是经过统计估计的客观事件概率。
证据如何参与：官方半年报预告、Q1财务包和行情/K线提供事实输入；研报/来源用于分母或可比背景。没有正权重外部目标锚，House价值不伪装为Street共识。
\begin{align*}
\mathrm{Shares} &= \frac{1540.52}{8.98}=171.5501\ \mathrm{100m\ shares}\\
\mathrm{EPS}_{H1}^d &= \frac{115.00}{171.5501}=0.6704\\
\mathrm{EPS}_s &= \mathrm{EPS}_{H1}^d(1+r_s)=1.04/1.33/1.44\\
V_s &= \mathrm{EPS}_s\times \mathrm{PE}_s=6.74/15.98/23.06\\
V_{prob} &= 0.30\times 6.74+0.50\times 15.98+0.20\times 23.06=14.62\\
P_{target} &= V_{prob}=14.62\quad(\text{外部锚权重}=0)\\
\mathrm{Upside} &= \frac{14.62}{8.98}-1=62.8\%\\
\end{align*}

\paragraph{600346 恒力石化：周期调整扣非 EPS / PE}
分母：H1 deducted EPS bridged with H2/H1 conversion；证据质量：中高；审计状态：PASS。
为什么选这个方法：石油石化盈利受商品价格、运价、价差或利用率周期影响，因此使用周期调整EPS/PE区间，而不是把某一个季度的PE当作全年估值。H2转换比例用于检验当前周期能否延续。
为什么选这个分母：使用H1扣非利润以减少非经常性收益污染；股本由总市值/当前价反推；H2/H1转换比例是明确的季节性和交付假设，不是机械年化。
为什么用这组倍数：熊/基准/牛PE区间是针对该周期业务类型的House假设，不是外部券商目标价；熊值较低反映周期反转，牛值较高必须由价格、销量和现金流共同确认。
为什么用这组概率：30\%/50\%/20\%是统一筛选先验：基准情景权重最高，同时保留熊市和牛市的实质性下行/上行贡献；它是House情景权重，不是经过统计估计的客观事件概率。
证据如何参与：官方半年报预告、Q1财务包和行情/K线提供事实输入；研报/来源用于分母或可比背景。没有正权重外部目标锚，House价值不伪装为Street共识。
\begin{align*}
\mathrm{Shares} &= \frac{1072.76}{15.24}=70.3911\ \mathrm{100m\ shares}\\
\mathrm{EPS}_{H1}^d &= \frac{53.30}{70.3911}=0.7572\\
\mathrm{EPS}_s &= \mathrm{EPS}_{H1}^d(1+r_s)=1.17/1.78/1.70\\
V_s &= \mathrm{EPS}_s\times \mathrm{PE}_s=7.04/16.02/20.44\\
V_{prob} &= 0.30\times 7.04+0.50\times 16.02+0.20\times 20.44=14.21\\
P_{target} &= V_{prob}=14.21\quad(\text{外部锚权重}=0)\\
\mathrm{Upside} &= \frac{14.21}{15.24}-1=-6.8\%\\
\end{align*}

\paragraph{000858 五粮液：常规扣非 EPS / PE}
分母：H1 deducted EPS bridged with H2/H1 conversion；证据质量：中高；审计状态：PASS。
为什么选这个方法：该标的H1扣非利润为正且未触发恢复估值阻断，因此扣非EPS是最透明的筛选分母；情景桥接避免把一个H1数字直接当作全年点预测。
为什么选这个分母：使用H1扣非利润以减少非经常性收益污染；股本由总市值/当前价反推；H2/H1转换比例是明确的季节性和交付假设，不是机械年化。
为什么用这组倍数：PE区间按公司所属行业匹配，表达估值压缩/正常/重估三种情景；这是House筛选区间，正式升级前仍需当前可审计的Street区间。
为什么用这组概率：30\%/50\%/20\%是统一筛选先验：基准情景权重最高，同时保留熊市和牛市的实质性下行/上行贡献；它是House情景权重，不是经过统计估计的客观事件概率。
证据如何参与：官方半年报预告、Q1财务包和行情/K线提供事实输入；研报/来源用于分母或可比背景。没有正权重外部目标锚，House价值不伪装为Street共识。
\begin{align*}
\mathrm{Shares} &= \frac{2965.55}{76.40}=38.8161\ \mathrm{100m\ shares}\\
\mathrm{EPS}_{H1}^d &= \frac{86.95}{38.8161}=2.2400\\
\mathrm{EPS}_s &= \mathrm{EPS}_{H1}^d(1+r_s)=3.70/4.26/4.93\\
V_s &= \mathrm{EPS}_s\times \mathrm{PE}_s=44.35/76.61/118.27\\
V_{prob} &= 0.30\times 44.35+0.50\times 76.61+0.20\times 118.27=75.26\\
P_{target} &= V_{prob}=75.26\quad(\text{外部锚权重}=0)\\
\mathrm{Upside} &= \frac{75.26}{76.40}-1=-1.5\%\\
\end{align*}

\paragraph{002460 赣锋锂业：周期调整 EPS / PE}
分母：2026E cycle-adjusted EPS, not H1 turnaround annualization；证据质量：中高；审计状态：PASS。
为什么选这个方法：最新H1利润存在低基数、扭亏或周期扭曲，因此先正常化EPS，再套业务匹配PE区间；业务匹配可比包括天齐锂业、中矿资源、盛新锂能。
为什么选这个分母：正常化分母为：2026E cycle-adjusted EPS, not H1 turnaround annualization；H1机械年化指标只作为市场隐含压力测试，不作为最终目标价分母。
为什么用这组倍数：熊/基准/牛倍数分别表达下行、正常兑现和上行重估；区间通过可比方法和当前价格隐含指标交叉检查（锂价上行带来盈利弹性，但周期持续性和海外资源兑现仍需验证），因此只能作为条件性House价值。
为什么用这组概率：30\%/50\%/20\%是统一筛选先验：基准情景权重最高，同时保留熊市和牛市的实质性下行/上行贡献；它是House情景权重，不是经过统计估计的客观事件概率。
证据如何参与：官方半年报预告、Q1财务包和行情/K线提供事实输入；研报/来源用于分母或可比背景。没有正权重外部目标锚，House价值不伪装为Street共识。
\begin{align*}
V_s=\mathrm{EPS}_{normalized,s}\times \mathrm{PE}_s\\
V_{b/m/b} &= 40.00/53.40/74.90\ \mathrm{CNY/share}\\
V_{prob} &= 0.30\times 40.00+0.50\times 53.40+0.20\times 74.90=53.68\\
P_{target} &= V_{prob}=53.68\quad(\text{外部锚权重}=0)\\
\mathrm{Upside} &= \frac{53.68}{51.50}-1=4.2\%\\
\end{align*}

\paragraph{002466 天齐锂业：周期调整扣非 EPS / PE}
分母：H1 deducted EPS bridged with H2/H1 conversion；证据质量：中高；审计状态：PASS。
为什么选这个方法：有色金属盈利受商品价格、运价、价差或利用率周期影响，因此使用周期调整EPS/PE区间，而不是把某一个季度的PE当作全年估值。H2转换比例用于检验当前周期能否延续。
为什么选这个分母：使用H1扣非利润以减少非经常性收益污染；股本由总市值/当前价反推；H2/H1转换比例是明确的季节性和交付假设，不是机械年化。
为什么用这组倍数：熊/基准/牛PE区间是针对该周期业务类型的House假设，不是外部券商目标价；熊值较低反映周期反转，牛值较高必须由价格、销量和现金流共同确认。
为什么用这组概率：30\%/50\%/20\%是统一筛选先验：基准情景权重最高，同时保留熊市和牛市的实质性下行/上行贡献；它是House情景权重，不是经过统计估计的客观事件概率。
证据如何参与：官方半年报预告、Q1财务包和行情/K线提供事实输入；研报/来源用于分母或可比背景。没有正权重外部目标锚，House价值不伪装为Street共识。
\begin{align*}
\mathrm{Shares} &= \frac{812.58}{47.43}=17.1322\ \mathrm{100m\ shares}\\
\mathrm{EPS}_{H1}^d &= \frac{35.05}{17.1322}=2.0459\\
\mathrm{EPS}_s &= \mathrm{EPS}_{H1}^d(1+r_s)=3.17/3.78/4.40\\
V_s &= \mathrm{EPS}_s\times \mathrm{PE}_s=25.37/45.42/70.38\\
V_{prob} &= 0.30\times 25.37+0.50\times 45.42+0.20\times 70.38=44.40\\
P_{target} &= V_{prob}=44.40\quad(\text{外部锚权重}=0)\\
\mathrm{Upside} &= \frac{44.40}{47.43}-1=-6.4\%\\
\end{align*}

\paragraph{601688 华泰证券：金融 PB/ROE + 盈利混合}
分母：Q1 BPS plus H1 deducted EPS bridged to H2；证据质量：中高；审计状态：PASS。
为什么选这个方法：金融公司主要通过净资产、ROE和盈利质量估值；PB保留资产负债表/ROE锚，PE提供独立盈利交叉检查，混合使用可以避免只依赖投资收益驱动的单季EPS。
为什么选这个分母：Q1 BPS作为资产分母，H1扣非EPS通过H2情景转换为熊/基准/牛盈利；这样将资产价值和盈利能力分开处理。
为什么用这组倍数：PB 0.75/1.05/1.35倍分别代表资产质量压力/正常/重估；行业PE区间提供盈利交叉检查；这些是House区间，不是外部目标价。
为什么用这组概率：30\%/50\%/20\%是统一筛选先验：基准情景权重最高，同时保留熊市和牛市的实质性下行/上行贡献；它是House情景权重，不是经过统计估计的客观事件概率。
证据如何参与：官方半年报预告、Q1财务包和行情/K线提供事实输入；研报/来源用于分母或可比背景。没有正权重外部目标锚，House价值不伪装为Street共识。
\begin{align*}
\mathrm{EPS}_{H1}^d &= \frac{115.19}{90.2685} = 1.2761\ \mathrm{CNY/share}\\
V_s &= 0.65\left(\mathrm{BPS}_1\times \mathrm{PB}_s\right)+0.35\left(\mathrm{EPS}_s\times \mathrm{PE}_s\right)\\
V_{b/m/b} &= 15.69/23.29/31.38\ \mathrm{CNY/share}\quad(\mathrm{PB}=14.80/20.72/26.64,\ \mathrm{PE}=17.36/28.08/40.20)\\
V_{prob} &= 0.30\times 15.47+0.50\times 23.29+0.20\times 31.38=22.56\\
P_{target} &= V_{prob}=22.56\quad(\text{外部锚权重}=0)\\
\mathrm{Upside} &= \frac{22.56}{20.63}-1=9.4\%\\
\end{align*}

\paragraph{601336 新华保险：金融 PB/ROE + 盈利混合}
分母：Q1 BPS plus H1 deducted EPS bridged to H2；证据质量：中高；审计状态：PASS。
为什么选这个方法：金融公司主要通过净资产、ROE和盈利质量估值；PB保留资产负债表/ROE锚，PE提供独立盈利交叉检查，混合使用可以避免只依赖投资收益驱动的单季EPS。
为什么选这个分母：Q1 BPS作为资产分母，H1扣非EPS通过H2情景转换为熊/基准/牛盈利；这样将资产价值和盈利能力分开处理。
为什么用这组倍数：PB 0.75/1.05/1.35倍分别代表资产质量压力/正常/重估；行业PE区间提供盈利交叉检查；这些是House区间，不是外部目标价。
为什么用这组概率：30\%/50\%/20\%是统一筛选先验：基准情景权重最高，同时保留熊市和牛市的实质性下行/上行贡献；它是House情景权重，不是经过统计估计的客观事件概率。
证据如何参与：官方半年报预告、Q1财务包和行情/K线提供事实输入；研报/来源用于分母或可比背景。没有正权重外部目标锚，House价值不伪装为Street共识。
\begin{align*}
\mathrm{EPS}_{H1}^d &= \frac{222.74}{31.1954} = 7.1400\ \mathrm{CNY/share}\\
V_s &= 0.65\left(\mathrm{BPS}_1\times \mathrm{PB}_s\right)+0.35\left(\mathrm{EPS}_s\times \mathrm{PE}_s\right)\\
V_{b/m/b} &= 53.29/82.01/113.47\ \mathrm{CNY/share}\quad(\mathrm{PB}=29.70/41.58/53.46,\ \mathrm{PE}=97.10/157.08/224.91)\\
V_{prob} &= 0.30\times 49.16+0.50\times 82.01+0.20\times 113.47=78.45\\
P_{target} &= V_{prob}=78.45\quad(\text{外部锚权重}=0)\\
\mathrm{Upside} &= \frac{78.45}{65.55}-1=19.7\%\\
\end{align*}

\paragraph{600037 歌华有线：扭亏折价正常化 EPS / PE}
分母：2026E recurring EPS after turnaround, not H1 profit annualization；证据质量：中；审计状态：PASS。
为什么选这个方法：最新H1利润存在低基数、扭亏或周期扭曲，因此先正常化EPS，再套业务匹配PE区间；业务匹配可比包括东方明珠、电广传媒、新媒股份。
为什么选这个分母：正常化分母为：2026E recurring EPS after turnaround, not H1 profit annualization；H1机械年化指标只作为市场隐含压力测试，不作为最终目标价分母。
为什么用这组倍数：熊/基准/牛倍数分别表达下行、正常兑现和上行重估；区间通过可比方法和当前价格隐含指标交叉检查（扭亏预期已被部分交易，通信业务转型和盈利持续性仍弱），因此只能作为条件性House价值。
为什么用这组概率：30\%/50\%/20\%是统一筛选先验：基准情景权重最高，同时保留熊市和牛市的实质性下行/上行贡献；它是House情景权重，不是经过统计估计的客观事件概率。
证据如何参与：官方半年报预告、Q1财务包和行情/K线提供事实输入；研报/来源用于分母或可比背景。没有正权重外部目标锚，House价值不伪装为Street共识。
\begin{align*}
V_s=\mathrm{EPS}_{normalized,s}\times \mathrm{PE}_s\\
V_{b/m/b} &= 2.40/5.40/11.00\ \mathrm{CNY/share}\\
V_{prob} &= 0.30\times 2.40+0.50\times 5.40+0.20\times 11.00=5.62\\
P_{target} &= V_{prob}=5.62\quad(\text{外部锚权重}=0)\\
\mathrm{Upside} &= \frac{5.62}{6.80}-1=-17.3\%\\
\end{align*}

\paragraph{002738 中矿资源：周期调整扣非 EPS / PE}
分母：H1 deducted EPS bridged with H2/H1 conversion；证据质量：中高；审计状态：PASS。
为什么选这个方法：有色金属盈利受商品价格、运价、价差或利用率周期影响，因此使用周期调整EPS/PE区间，而不是把某一个季度的PE当作全年估值。H2转换比例用于检验当前周期能否延续。
为什么选这个分母：使用H1扣非利润以减少非经常性收益污染；股本由总市值/当前价反推；H2/H1转换比例是明确的季节性和交付假设，不是机械年化。
为什么用这组倍数：熊/基准/牛PE区间是针对该周期业务类型的House假设，不是外部券商目标价；熊值较低反映周期反转，牛值较高必须由价格、销量和现金流共同确认。
为什么用这组概率：30\%/50\%/20\%是统一筛选先验：基准情景权重最高，同时保留熊市和牛市的实质性下行/上行贡献；它是House情景权重，不是经过统计估计的客观事件概率。
证据如何参与：官方半年报预告、Q1财务包和行情/K线提供事实输入；研报/来源用于分母或可比背景。没有正权重外部目标锚，House价值不伪装为Street共识。
\begin{align*}
\mathrm{Shares} &= \frac{355.91}{49.33}=7.2149\ \mathrm{100m\ shares}\\
\mathrm{EPS}_{H1}^d &= \frac{11.00}{7.2149}=1.5246\\
\mathrm{EPS}_s &= \mathrm{EPS}_{H1}^d(1+r_s)=2.36/3.85/3.28\\
V_s &= \mathrm{EPS}_s\times \mathrm{PE}_s=18.91/46.20/52.45\\
V_{prob} &= 0.30\times 18.91+0.50\times 46.20+0.20\times 52.45=39.26\\
P_{target} &= V_{prob}=39.26\quad(\text{外部锚权重}=0)\\
\mathrm{Upside} &= \frac{39.26}{49.33}-1=-20.4\%\\
\end{align*}

\paragraph{601069 西部黄金：周期调整扣非 EPS / PE}
分母：H1 deducted EPS bridged with H2/H1 conversion；证据质量：中；审计状态：PASS。
为什么选这个方法：有色金属盈利受商品价格、运价、价差或利用率周期影响，因此使用周期调整EPS/PE区间，而不是把某一个季度的PE当作全年估值。H2转换比例用于检验当前周期能否延续。
为什么选这个分母：使用H1扣非利润以减少非经常性收益污染；股本由总市值/当前价反推；H2/H1转换比例是明确的季节性和交付假设，不是机械年化。
为什么用这组倍数：熊/基准/牛PE区间是针对该周期业务类型的House假设，不是外部券商目标价；熊值较低反映周期反转，牛值较高必须由价格、销量和现金流共同确认。
为什么用这组概率：30\%/50\%/20\%是统一筛选先验：基准情景权重最高，同时保留熊市和牛市的实质性下行/上行贡献；它是House情景权重，不是经过统计估计的客观事件概率。
证据如何参与：官方半年报预告、Q1财务包和行情/K线提供事实输入；研报/来源用于分母或可比背景。没有正权重外部目标锚，House价值不伪装为Street共识。
\begin{align*}
\mathrm{Shares} &= \frac{209.17}{22.96}=9.1102\ \mathrm{100m\ shares}\\
\mathrm{EPS}_{H1}^d &= \frac{5.35}{9.1102}=0.5873\\
\mathrm{EPS}_s &= \mathrm{EPS}_{H1}^d(1+r_s)=0.91/1.09/1.26\\
V_s &= \mathrm{EPS}_s\times \mathrm{PE}_s=7.28/13.04/20.20\\
V_{prob} &= 0.30\times 7.28+0.50\times 13.04+0.20\times 20.20=12.74\\
P_{target} &= V_{prob}=12.74\quad(\text{外部锚权重}=0)\\
\mathrm{Upside} &= \frac{12.74}{22.96}-1=-44.5\%\\
\end{align*}

\paragraph{000567 海德股份：金融 PB/ROE + 盈利混合}
分母：Q1 BPS plus H1 deducted EPS bridged to H2；证据质量：中；审计状态：PASS。
为什么选这个方法：金融公司主要通过净资产、ROE和盈利质量估值；PB保留资产负债表/ROE锚，PE提供独立盈利交叉检查，混合使用可以避免只依赖投资收益驱动的单季EPS。
为什么选这个分母：Q1 BPS作为资产分母，H1扣非EPS通过H2情景转换为熊/基准/牛盈利；这样将资产价值和盈利能力分开处理。
为什么用这组倍数：PB 0.75/1.05/1.35倍分别代表资产质量压力/正常/重估；行业PE区间提供盈利交叉检查；这些是House区间，不是外部目标价。
为什么用这组概率：30\%/50\%/20\%是统一筛选先验：基准情景权重最高，同时保留熊市和牛市的实质性下行/上行贡献；它是House情景权重，不是经过统计估计的客观事件概率。
证据如何参与：官方半年报预告、Q1财务包和行情/K线提供事实输入；研报/来源用于分母或可比背景。没有正权重外部目标锚，House价值不伪装为Street共识。
\begin{align*}
\mathrm{EPS}_{H1}^d &= \frac{14.40}{19.5462} = 0.7367\ \mathrm{CNY/share}\\
V_s &= 0.65\left(\mathrm{BPS}_1\times \mathrm{PB}_s\right)+0.35\left(\mathrm{EPS}_s\times \mathrm{PE}_s\right)\\
V_{b/m/b} &= 5.01/7.78/10.83\ \mathrm{CNY/share}\quad(\mathrm{PB}=2.31/3.24/4.16,\ \mathrm{PE}=10.02/16.21/23.21)\\
V_{prob} &= 0.30\times 4.63+0.50\times 7.78+0.20\times 10.83=7.45\\
P_{target} &= V_{prob}=7.45\quad(\text{外部锚权重}=0)\\
\mathrm{Upside} &= \frac{7.45}{6.17}-1=20.8\%\\
\end{align*}

\paragraph{002602 世纪华通：成长盈利情景 PE}
分母：H1 deducted EPS bridged with H2/H1 conversion；证据质量：中高；审计状态：PASS。
为什么选这个方法：传媒具有成长或盈利持续期特征，因此模型给盈利桥接配置较高PE区间，而不是直接给主题收入估值；价格和现金流检查防止把泛化成长叙事直接变成EPS。
为什么选这个分母：使用H1扣非利润以减少非经常性收益污染；股本由总市值/当前价反推；H2/H1转换比例是明确的季节性和交付假设，不是机械年化。
为什么用这组倍数：PE区间是House对成长持续期的情景判断；牛市情景要求业绩桥接和经营证据持续，不是市场一致目标价。
为什么用这组概率：30\%/50\%/20\%是统一筛选先验：基准情景权重最高，同时保留熊市和牛市的实质性下行/上行贡献；它是House情景权重，不是经过统计估计的客观事件概率。
证据如何参与：官方半年报预告、Q1财务包和行情/K线提供事实输入；研报/来源用于分母或可比背景。没有正权重外部目标锚，House价值不伪装为Street共识。
\begin{align*}
\mathrm{Shares} &= \frac{1051.06}{14.26}=73.7069\ \mathrm{100m\ shares}\\
\mathrm{EPS}_{H1}^d &= \frac{43.15}{73.7069}=0.5854\\
\mathrm{EPS}_s &= \mathrm{EPS}_{H1}^d(1+r_s)=1.00/1.21/1.41\\
V_s &= \mathrm{EPS}_s\times \mathrm{PE}_s=9.95/16.94/25.29\\
V_{prob} &= 0.30\times 9.95+0.50\times 16.94+0.20\times 25.29=16.51\\
P_{target} &= V_{prob}=16.51\quad(\text{外部锚权重}=0)\\
\mathrm{Upside} &= \frac{16.51}{14.26}-1=15.8\%\\
\end{align*}

\paragraph{600026 中远海能：周期调整扣非 EPS / PE}
分母：H1 deducted EPS bridged with H2/H1 conversion；证据质量：中高；审计状态：PASS。
为什么选这个方法：交通运输盈利受商品价格、运价、价差或利用率周期影响，因此使用周期调整EPS/PE区间，而不是把某一个季度的PE当作全年估值。H2转换比例用于检验当前周期能否延续。
为什么选这个分母：使用H1扣非利润以减少非经常性收益污染；股本由总市值/当前价反推；H2/H1转换比例是明确的季节性和交付假设，不是机械年化。
为什么用这组倍数：熊/基准/牛PE区间是针对该周期业务类型的House假设，不是外部券商目标价；熊值较低反映周期反转，牛值较高必须由价格、销量和现金流共同确认。
为什么用这组概率：30\%/50\%/20\%是统一筛选先验：基准情景权重最高，同时保留熊市和牛市的实质性下行/上行贡献；它是House情景权重，不是经过统计估计的客观事件概率。
证据如何参与：官方半年报预告、Q1财务包和行情/K线提供事实输入；研报/来源用于分母或可比背景。没有正权重外部目标锚，House价值不伪装为Street共识。
\begin{align*}
\mathrm{Shares} &= \frac{866.09}{15.83}=54.7119\ \mathrm{100m\ shares}\\
\mathrm{EPS}_{H1}^d &= \frac{44.00}{54.7119}=0.8042\\
\mathrm{EPS}_s &= \mathrm{EPS}_{H1}^d(1+r_s)=1.25/2.14/1.73\\
V_s &= \mathrm{EPS}_s\times \mathrm{PE}_s=9.97/25.64/27.66\\
V_{prob} &= 0.30\times 9.97+0.50\times 25.64+0.20\times 27.66=21.34\\
P_{target} &= V_{prob}=21.34\quad(\text{外部锚权重}=0)\\
\mathrm{Upside} &= \frac{21.34}{15.83}-1=34.8\%\\
\end{align*}

\paragraph{002048 宁波华翔：正常化 EPS / PE}
分母：2026E normalized EPS after auto-cycle and robot-order verification；证据质量：中；审计状态：PASS。
为什么选这个方法：最新H1利润存在低基数、扭亏或周期扭曲，因此先正常化EPS，再套业务匹配PE区间；业务匹配可比包括华域汽车、福耀玻璃、拓普集团。
为什么选这个分母：正常化分母为：2026E normalized EPS after auto-cycle and robot-order verification；H1机械年化指标只作为市场隐含压力测试，不作为最终目标价分母。
为什么用这组倍数：熊/基准/牛倍数分别表达下行、正常兑现和上行重估；区间通过可比方法和当前价格隐含指标交叉检查（汽车零部件和机器人业务提供弹性，但预测时点较旧），因此只能作为条件性House价值。
为什么用这组概率：30\%/50\%/20\%是统一筛选先验：基准情景权重最高，同时保留熊市和牛市的实质性下行/上行贡献；它是House情景权重，不是经过统计估计的客观事件概率。
证据如何参与：官方半年报预告、Q1财务包和行情/K线提供事实输入；研报/来源用于分母或可比背景。没有正权重外部目标锚，House价值不伪装为Street共识。
\begin{align*}
V_s=\mathrm{EPS}_{normalized,s}\times \mathrm{PE}_s\\
V_{b/m/b} &= 19.20/29.64/39.60\ \mathrm{CNY/share}\\
V_{prob} &= 0.30\times 19.20+0.50\times 29.64+0.20\times 39.60=28.50\\
P_{target} &= V_{prob}=28.50\quad(\text{外部锚权重}=0)\\
\mathrm{Upside} &= \frac{28.50}{21.43}-1=33.0\%\\
\end{align*}

\paragraph{600688 上海石化：PB + 周期调整资产价值}
分母：Q1 BPS and cycle-adjusted asset value；证据质量：中；审计状态：PASS。
为什么选这个方法：该业务资产较重或受项目周期影响，因此资产价值/BPS和周期调整盈利比直接年化结算利润更稳定；业务匹配可比包括中国石化、恒力石化、华锦股份。
为什么选这个分母：正常化分母为：Q1 BPS and cycle-adjusted asset value；H1机械年化指标只作为市场隐含压力测试，不作为最终目标价分母。
为什么用这组倍数：熊/基准/牛倍数分别表达下行、正常兑现和上行重估；区间通过可比方法和当前价格隐含指标交叉检查（炼化周期和现金流承压，历史目标价不能作为当前Street锚），因此只能作为条件性House价值。
为什么用这组概率：30\%/50\%/20\%是统一筛选先验：基准情景权重最高，同时保留熊市和牛市的实质性下行/上行贡献；它是House情景权重，不是经过统计估计的客观事件概率。
证据如何参与：官方半年报预告、Q1财务包和行情/K线提供事实输入；研报/来源用于分母或可比背景。外部目标价4.06元，仅作背景，权重为0。
\begin{align*}
V_s=\mathrm{BPS}_{normalized}\times \mathrm{PB}_s\\
V_{b/m/b} &= 2.01/2.68/3.35\ \mathrm{CNY/share}\\
V_{prob} &= 0.30\times 2.01+0.50\times 2.68+0.20\times 3.35=2.61\\
P_{target} &= V_{prob}=2.61\quad(\text{外部锚权重}=0)\\
\mathrm{Upside} &= \frac{2.61}{2.84}-1=-8.1\%\\
\end{align*}

\paragraph{600095 湘财股份：金融 PB/ROE + 盈利混合}
分母：Q1 BPS plus H1 deducted EPS bridged to H2；证据质量：中；审计状态：PASS。
为什么选这个方法：金融公司主要通过净资产、ROE和盈利质量估值；PB保留资产负债表/ROE锚，PE提供独立盈利交叉检查，混合使用可以避免只依赖投资收益驱动的单季EPS。
为什么选这个分母：Q1 BPS作为资产分母，H1扣非EPS通过H2情景转换为熊/基准/牛盈利；这样将资产价值和盈利能力分开处理。
为什么用这组倍数：PB 0.75/1.05/1.35倍分别代表资产质量压力/正常/重估；行业PE区间提供盈利交叉检查；这些是House区间，不是外部目标价。
为什么用这组概率：30\%/50\%/20\%是统一筛选先验：基准情景权重最高，同时保留熊市和牛市的实质性下行/上行贡献；它是House情景权重，不是经过统计估计的客观事件概率。
证据如何参与：官方半年报预告、Q1财务包和行情/K线提供事实输入；研报/来源用于分母或可比背景。没有正权重外部目标锚，House价值不伪装为Street共识。
\begin{align*}
\mathrm{EPS}_{H1}^d &= \frac{5.15}{28.5915} = 0.1801\ \mathrm{CNY/share}\\
V_s &= 0.65\left(\mathrm{BPS}_1\times \mathrm{PB}_s\right)+0.35\left(\mathrm{EPS}_s\times \mathrm{PE}_s\right)\\
V_{b/m/b} &= 2.96/4.32/5.76\ \mathrm{CNY/share}\quad(\mathrm{PB}=3.23/4.52/5.81,\ \mathrm{PE}=2.45/3.96/5.67)\\
V_{prob} &= 0.30\times 2.95+0.50\times 4.32+0.20\times 5.76=4.20\\
P_{target} &= V_{prob}=4.20\quad(\text{外部锚权重}=0)\\
\mathrm{Upside} &= \frac{4.20}{7.81}-1=-46.2\%\\
\end{align*}

\paragraph{000426 兴业银锡：周期调整扣非 EPS / PE}
分母：H1 deducted EPS bridged with H2/H1 conversion；证据质量：中高；审计状态：PASS。
为什么选这个方法：有色金属盈利受商品价格、运价、价差或利用率周期影响，因此使用周期调整EPS/PE区间，而不是把某一个季度的PE当作全年估值。H2转换比例用于检验当前周期能否延续。
为什么选这个分母：使用H1扣非利润以减少非经常性收益污染；股本由总市值/当前价反推；H2/H1转换比例是明确的季节性和交付假设，不是机械年化。
为什么用这组倍数：熊/基准/牛PE区间是针对该周期业务类型的House假设，不是外部券商目标价；熊值较低反映周期反转，牛值较高必须由价格、销量和现金流共同确认。
为什么用这组概率：30\%/50\%/20\%是统一筛选先验：基准情景权重最高，同时保留熊市和牛市的实质性下行/上行贡献；它是House情景权重，不是经过统计估计的客观事件概率。
证据如何参与：官方半年报预告、Q1财务包和行情/K线提供事实输入；研报/来源用于分母或可比背景。没有正权重外部目标锚，House价值不伪装为Street共识。
\begin{align*}
\mathrm{Shares} &= \frac{542.63}{30.56}=17.7562\ \mathrm{100m\ shares}\\
\mathrm{EPS}_{H1}^d &= \frac{18.05}{17.7562}=1.0165\\
\mathrm{EPS}_s &= \mathrm{EPS}_{H1}^d(1+r_s)=1.58/2.13/2.19\\
V_s &= \mathrm{EPS}_s\times \mathrm{PE}_s=12.61/25.56/34.97\\
V_{prob} &= 0.30\times 12.61+0.50\times 25.56+0.20\times 34.97=23.56\\
P_{target} &= V_{prob}=23.56\quad(\text{外部锚权重}=0)\\
\mathrm{Upside} &= \frac{23.56}{30.56}-1=-22.9\%\\
\end{align*}

\paragraph{600489 中金黄金：周期调整扣非 EPS / PE}
分母：H1 deducted EPS bridged with H2/H1 conversion；证据质量：中高；审计状态：PASS。
为什么选这个方法：有色金属盈利受商品价格、运价、价差或利用率周期影响，因此使用周期调整EPS/PE区间，而不是把某一个季度的PE当作全年估值。H2转换比例用于检验当前周期能否延续。
为什么选这个分母：使用H1扣非利润以减少非经常性收益污染；股本由总市值/当前价反推；H2/H1转换比例是明确的季节性和交付假设，不是机械年化。
为什么用这组倍数：熊/基准/牛PE区间是针对该周期业务类型的House假设，不是外部券商目标价；熊值较低反映周期反转，牛值较高必须由价格、销量和现金流共同确认。
为什么用这组概率：30\%/50\%/20\%是统一筛选先验：基准情景权重最高，同时保留熊市和牛市的实质性下行/上行贡献；它是House情景权重，不是经过统计估计的客观事件概率。
证据如何参与：官方半年报预告、Q1财务包和行情/K线提供事实输入；研报/来源用于分母或可比背景。没有正权重外部目标锚，House价值不伪装为Street共识。
\begin{align*}
\mathrm{Shares} &= \frac{975.28}{20.12}=48.4732\ \mathrm{100m\ shares}\\
\mathrm{EPS}_{H1}^d &= \frac{43.00}{48.4732}=0.8871\\
\mathrm{EPS}_s &= \mathrm{EPS}_{H1}^d(1+r_s)=1.38/1.96/1.91\\
V_s &= \mathrm{EPS}_s\times \mathrm{PE}_s=11.00/23.52/30.52\\
V_{prob} &= 0.30\times 11.00+0.50\times 23.52+0.20\times 30.52=21.16\\
P_{target} &= V_{prob}=21.16\quad(\text{外部锚权重}=0)\\
\mathrm{Upside} &= \frac{21.16}{20.12}-1=5.2\%\\
\end{align*}

\paragraph{000737 北方铜业：周期调整扣非 EPS / PE}
分母：H1 deducted EPS bridged with H2/H1 conversion；证据质量：中低；审计状态：PASS。
为什么选这个方法：有色金属盈利受商品价格、运价、价差或利用率周期影响，因此使用周期调整EPS/PE区间，而不是把某一个季度的PE当作全年估值。H2转换比例用于检验当前周期能否延续。
为什么选这个分母：使用H1扣非利润以减少非经常性收益污染；股本由总市值/当前价反推；H2/H1转换比例是明确的季节性和交付假设，不是机械年化。
为什么用这组倍数：熊/基准/牛PE区间是针对该周期业务类型的House假设，不是外部券商目标价；熊值较低反映周期反转，牛值较高必须由价格、销量和现金流共同确认。
为什么用这组概率：30\%/50\%/20\%是统一筛选先验：基准情景权重最高，同时保留熊市和牛市的实质性下行/上行贡献；它是House情景权重，不是经过统计估计的客观事件概率。
证据如何参与：官方半年报预告、Q1财务包和行情/K线提供事实输入；研报/来源用于分母或可比背景。没有正权重外部目标锚，House价值不伪装为Street共识。
\begin{align*}
\mathrm{Shares} &= \frac{242.85}{12.75}=19.0471\ \mathrm{100m\ shares}\\
\mathrm{EPS}_{H1}^d &= \frac{13.52}{19.0471}=0.7099\\
\mathrm{EPS}_s &= \mathrm{EPS}_{H1}^d(1+r_s)=1.10/1.31/1.53\\
V_s &= \mathrm{EPS}_s\times \mathrm{PE}_s=8.80/15.76/24.42\\
V_{prob} &= 0.30\times 8.80+0.50\times 15.76+0.20\times 24.42=15.40\\
P_{target} &= V_{prob}=15.40\quad(\text{外部锚权重}=0)\\
\mathrm{Upside} &= \frac{15.40}{12.75}-1=20.8\%\\
\end{align*}

\paragraph{002756 永兴材料：周期调整扣非 EPS / PE}
分母：H1 deducted EPS bridged with H2/H1 conversion；证据质量：中高；审计状态：PASS。
为什么选这个方法：有色金属盈利受商品价格、运价、价差或利用率周期影响，因此使用周期调整EPS/PE区间，而不是把某一个季度的PE当作全年估值。H2转换比例用于检验当前周期能否延续。
为什么选这个分母：使用H1扣非利润以减少非经常性收益污染；股本由总市值/当前价反推；H2/H1转换比例是明确的季节性和交付假设，不是机械年化。
为什么用这组倍数：熊/基准/牛PE区间是针对该周期业务类型的House假设，不是外部券商目标价；熊值较低反映周期反转，牛值较高必须由价格、销量和现金流共同确认。
为什么用这组概率：30\%/50\%/20\%是统一筛选先验：基准情景权重最高，同时保留熊市和牛市的实质性下行/上行贡献；它是House情景权重，不是经过统计估计的客观事件概率。
证据如何参与：官方半年报预告、Q1财务包和行情/K线提供事实输入；研报/来源用于分母或可比背景。没有正权重外部目标锚，House价值不伪装为Street共识。
\begin{align*}
\mathrm{Shares} &= \frac{242.06}{44.90}=5.3911\ \mathrm{100m\ shares}\\
\mathrm{EPS}_{H1}^d &= \frac{10.00}{5.3911}=1.8549\\
\mathrm{EPS}_s &= \mathrm{EPS}_{H1}^d(1+r_s)=2.88/3.74/3.99\\
V_s &= \mathrm{EPS}_s\times \mathrm{PE}_s=23.00/44.88/63.81\\
V_{prob} &= 0.30\times 23.00+0.50\times 44.88+0.20\times 63.81=42.10\\
P_{target} &= V_{prob}=42.10\quad(\text{外部锚权重}=0)\\
\mathrm{Upside} &= \frac{42.10}{44.90}-1=-6.2\%\\
\end{align*}

\paragraph{300014 亿纬锂能：成长盈利情景 PE}
分母：H1 deducted EPS bridged with H2/H1 conversion；证据质量：中高；审计状态：PASS。
为什么选这个方法：电力设备具有成长或盈利持续期特征，因此模型给盈利桥接配置较高PE区间，而不是直接给主题收入估值；价格和现金流检查防止把泛化成长叙事直接变成EPS。
为什么选这个分母：使用H1扣非利润以减少非经常性收益污染；股本由总市值/当前价反推；H2/H1转换比例是明确的季节性和交付假设，不是机械年化。
为什么用这组倍数：PE区间是House对成长持续期的情景判断；牛市情景要求业绩桥接和经营证据持续，不是市场一致目标价。
为什么用这组概率：30\%/50\%/20\%是统一筛选先验：基准情景权重最高，同时保留熊市和牛市的实质性下行/上行贡献；它是House情景权重，不是经过统计估计的客观事件概率。
证据如何参与：官方半年报预告、Q1财务包和行情/K线提供事实输入；研报/来源用于分母或可比背景。没有正权重外部目标锚，House价值不伪装为Street共识。
\begin{align*}
\mathrm{Shares} &= \frac{1136.20}{52.28}=21.7330\ \mathrm{100m\ shares}\\
\mathrm{EPS}_{H1}^d &= \frac{25.16}{21.7330}=1.1579\\
\mathrm{EPS}_s &= \mathrm{EPS}_{H1}^d(1+r_s)=1.97/3.11/2.66\\
V_s &= \mathrm{EPS}_s\times \mathrm{PE}_s=31.49/68.46/74.57\\
V_{prob} &= 0.30\times 31.49+0.50\times 68.46+0.20\times 74.57=58.59\\
P_{target} &= V_{prob}=58.59\quad(\text{外部锚权重}=0)\\
\mathrm{Upside} &= \frac{58.59}{52.28}-1=12.1\%\\
\end{align*}

\paragraph{601061 中信金属：常规扣非 EPS / PE}
分母：H1 deducted EPS bridged with H2/H1 conversion；证据质量：中；审计状态：PASS。
为什么选这个方法：该标的H1扣非利润为正且未触发恢复估值阻断，因此扣非EPS是最透明的筛选分母；情景桥接避免把一个H1数字直接当作全年点预测。
为什么选这个分母：使用H1扣非利润以减少非经常性收益污染；股本由总市值/当前价反推；H2/H1转换比例是明确的季节性和交付假设，不是机械年化。
为什么用这组倍数：PE区间按公司所属行业匹配，表达估值压缩/正常/重估三种情景；这是House筛选区间，正式升级前仍需当前可审计的Street区间。
为什么用这组概率：30\%/50\%/20\%是统一筛选先验：基准情景权重最高，同时保留熊市和牛市的实质性下行/上行贡献；它是House情景权重，不是经过统计估计的客观事件概率。
证据如何参与：官方半年报预告、Q1财务包和行情/K线提供事实输入；研报/来源用于分母或可比背景。外部目标价11.25元，赋予正权重5\%。
\begin{align*}
\mathrm{Shares} &= \frac{530.18}{10.82}=49.0000\ \mathrm{100m\ shares}\\
\mathrm{EPS}_{H1}^d &= \frac{26.45}{49.0000}=0.5398\\
\mathrm{EPS}_s &= \mathrm{EPS}_{H1}^d(1+r_s)=0.89/1.03/1.19\\
V_s &= \mathrm{EPS}_s\times \mathrm{PE}_s=8.12/18.46/28.50\\
V_{prob} &= 0.30\times 8.12+0.50\times 18.46+0.20\times 28.50=17.37\\
P_{target} &= 17.37\times 0.95+11.25\times 0.05=17.06\\
\mathrm{Upside} &= \frac{17.06}{10.82}-1=57.7\%\\
\end{align*}

\paragraph{600739 辽宁成大：成长盈利情景 PE}
分母：H1 deducted EPS bridged with H2/H1 conversion；证据质量：中；审计状态：PASS。
为什么选这个方法：医药生物具有成长或盈利持续期特征，因此模型给盈利桥接配置较高PE区间，而不是直接给主题收入估值；价格和现金流检查防止把泛化成长叙事直接变成EPS。
为什么选这个分母：使用H1扣非利润以减少非经常性收益污染；股本由总市值/当前价反推；H2/H1转换比例是明确的季节性和交付假设，不是机械年化。
为什么用这组倍数：PE区间是House对成长持续期的情景判断；牛市情景要求业绩桥接和经营证据持续，不是市场一致目标价。
为什么用这组概率：30\%/50\%/20\%是统一筛选先验：基准情景权重最高，同时保留熊市和牛市的实质性下行/上行贡献；它是House情景权重，不是经过统计估计的客观事件概率。
证据如何参与：官方半年报预告、Q1财务包和行情/K线提供事实输入；研报/来源用于分母或可比背景。外部目标价22.79元，赋予正权重5\%。
\begin{align*}
\mathrm{Shares} &= \frac{166.69}{10.95}=15.2228\ \mathrm{100m\ shares}\\
\mathrm{EPS}_{H1}^d &= \frac{15.84}{15.2228}=1.0407\\
\mathrm{EPS}_s &= \mathrm{EPS}_{H1}^d(1+r_s)=1.77/2.08/2.39\\
V_s &= \mathrm{EPS}_s\times \mathrm{PE}_s=8.21/52.03/76.60\\
V_{prob} &= 0.30\times 8.21+0.50\times 52.03+0.20\times 76.60=43.80\\
P_{target} &= 43.80\times 0.95+22.79\times 0.05=42.75\\
\mathrm{Upside} &= \frac{42.75}{10.95}-1=290.4\%\\
\end{align*}

\paragraph{601360 三六零：2026E 收入 PS}
分母：2026E revenue consensus and recurring AI/security monetization, not H1 turnaround EPS；证据质量：中；审计状态：PASS。
为什么选这个方法：EPS暂时过于波动或接近零，因此收入作为主分母，再用公开共识和可比公司进行PS三角验证；业务匹配可比包括昆仑万维、启明星辰、金山办公。
为什么选这个分母：正常化分母为：2026E revenue consensus and recurring AI/security monetization, not H1 turnaround EPS；H1机械年化指标只作为市场隐含压力测试，不作为最终目标价分母。
为什么用这组倍数：熊/基准/牛倍数分别表达下行、正常兑现和上行重估；区间通过可比方法和当前价格隐含指标交叉检查（AI+安全重估，但资金参与度偏弱、融资余额低于一年30\%分位），因此只能作为条件性House价值。
为什么用这组概率：30\%/50\%/20\%是统一筛选先验：基准情景权重最高，同时保留熊市和牛市的实质性下行/上行贡献；它是House情景权重，不是经过统计估计的客观事件概率。
证据如何参与：官方半年报预告、Q1财务包和行情/K线提供事实输入；研报/来源用于分母或可比背景。外部目标价14.44元，赋予正权重10\%。
\begin{align*}
V_s = \frac{R_{2026E}}{Shares}\times PS_s=\frac{95.91}{69.9956}\times PS_s\\
V_{b/m/b} &= 10.28/13.70/15.07\ \mathrm{CNY/share}\\
V_{prob} &= 0.30\times 10.28+0.50\times 13.70+0.20\times 15.07=12.95\\
P_{target} &= 12.95\times 0.90+14.44\times 0.10=13.10\\
\mathrm{Upside} &= \frac{13.10}{9.00}-1=45.6\%\\
\end{align*}

\paragraph{600918 中泰证券：金融 PB/ROE + 盈利混合}
分母：Q1 BPS plus H1 deducted EPS bridged to H2；证据质量：中；审计状态：PASS。
为什么选这个方法：金融公司主要通过净资产、ROE和盈利质量估值；PB保留资产负债表/ROE锚，PE提供独立盈利交叉检查，混合使用可以避免只依赖投资收益驱动的单季EPS。
为什么选这个分母：Q1 BPS作为资产分母，H1扣非EPS通过H2情景转换为熊/基准/牛盈利；这样将资产价值和盈利能力分开处理。
为什么用这组倍数：PB 0.75/1.05/1.35倍分别代表资产质量压力/正常/重估；行业PE区间提供盈利交叉检查；这些是House区间，不是外部目标价。
为什么用这组概率：30\%/50\%/20\%是统一筛选先验：基准情景权重最高，同时保留熊市和牛市的实质性下行/上行贡献；它是House情景权重，不是经过统计估计的客观事件概率。
证据如何参与：官方半年报预告、Q1财务包和行情/K线提供事实输入；研报/来源用于分母或可比背景。没有正权重外部目标锚，House价值不伪装为Street共识。
\begin{align*}
\mathrm{EPS}_{H1}^d &= \frac{16.98}{79.1839} = 0.2144\ \mathrm{CNY/share}\\
V_s &= 0.65\left(\mathrm{BPS}_1\times \mathrm{PB}_s\right)+0.35\left(\mathrm{EPS}_s\times \mathrm{PE}_s\right)\\
V_{b/m/b} &= 3.77/5.50/7.32\ \mathrm{CNY/share}\quad(\mathrm{PB}=4.23/5.93/7.62,\ \mathrm{PE}=2.92/4.72/6.75)\\
V_{prob} &= 0.30\times 3.77+0.50\times 5.50+0.20\times 7.32=5.35\\
P_{target} &= V_{prob}=5.35\quad(\text{外部锚权重}=0)\\
\mathrm{Upside} &= \frac{5.35}{5.71}-1=-6.3\%\\
\end{align*}

\paragraph{000975 山金国际：周期调整扣非 EPS / PE}
分母：H1 deducted EPS bridged with H2/H1 conversion；证据质量：中高；审计状态：PASS。
为什么选这个方法：有色金属盈利受商品价格、运价、价差或利用率周期影响，因此使用周期调整EPS/PE区间，而不是把某一个季度的PE当作全年估值。H2转换比例用于检验当前周期能否延续。
为什么选这个分母：使用H1扣非利润以减少非经常性收益污染；股本由总市值/当前价反推；H2/H1转换比例是明确的季节性和交付假设，不是机械年化。
为什么用这组倍数：熊/基准/牛PE区间是针对该周期业务类型的House假设，不是外部券商目标价；熊值较低反映周期反转，牛值较高必须由价格、销量和现金流共同确认。
为什么用这组概率：30\%/50\%/20\%是统一筛选先验：基准情景权重最高，同时保留熊市和牛市的实质性下行/上行贡献；它是House情景权重，不是经过统计估计的客观事件概率。
证据如何参与：官方半年报预告、Q1财务包和行情/K线提供事实输入；研报/来源用于分母或可比背景。没有正权重外部目标锚，House价值不伪装为Street共识。
\begin{align*}
\mathrm{Shares} &= \frac{536.21}{19.34}=27.7254\ \mathrm{100m\ shares}\\
\mathrm{EPS}_{H1}^d &= \frac{23.45}{27.7254}=0.8458\\
\mathrm{EPS}_s &= \mathrm{EPS}_{H1}^d(1+r_s)=1.31/1.89/1.82\\
V_s &= \mathrm{EPS}_s\times \mathrm{PE}_s=10.49/22.68/29.10\\
V_{prob} &= 0.30\times 10.49+0.50\times 22.68+0.20\times 29.10=20.31\\
P_{target} &= V_{prob}=20.31\quad(\text{外部锚权重}=0)\\
\mathrm{Upside} &= \frac{20.31}{19.34}-1=5.0\%\\
\end{align*}

\paragraph{000623 吉林敖东：成长盈利情景 PE}
分母：H1 deducted EPS bridged with H2/H1 conversion；证据质量：中；审计状态：PASS。
为什么选这个方法：医药生物具有成长或盈利持续期特征，因此模型给盈利桥接配置较高PE区间，而不是直接给主题收入估值；价格和现金流检查防止把泛化成长叙事直接变成EPS。
为什么选这个分母：使用H1扣非利润以减少非经常性收益污染；股本由总市值/当前价反推；H2/H1转换比例是明确的季节性和交付假设，不是机械年化。
为什么用这组倍数：PE区间是House对成长持续期的情景判断；牛市情景要求业绩桥接和经营证据持续，不是市场一致目标价。
为什么用这组概率：30\%/50\%/20\%是统一筛选先验：基准情景权重最高，同时保留熊市和牛市的实质性下行/上行贡献；它是House情景权重，不是经过统计估计的客观事件概率。
证据如何参与：官方半年报预告、Q1财务包和行情/K线提供事实输入；研报/来源用于分母或可比背景。没有正权重外部目标锚，House价值不伪装为Street共识。
\begin{align*}
\mathrm{Shares} &= \frac{223.15}{18.66}=11.9587\ \mathrm{100m\ shares}\\
\mathrm{EPS}_{H1}^d &= \frac{22.16}{11.9587}=1.8534\\
\mathrm{EPS}_s &= \mathrm{EPS}_{H1}^d(1+r_s)=3.15/3.71/4.26\\
V_s &= \mathrm{EPS}_s\times \mathrm{PE}_s=14.00/92.67/136.41\\
V_{prob} &= 0.30\times 14.00+0.50\times 92.67+0.20\times 136.41=77.82\\
P_{target} &= V_{prob}=77.82\quad(\text{外部锚权重}=0)\\
\mathrm{Upside} &= \frac{77.82}{18.66}-1=317.0\%\\
\end{align*}

\paragraph{600292 电投水电：常规扣非 EPS / PE}
分母：H1 deducted EPS bridged with H2/H1 conversion；证据质量：中；审计状态：PASS。
为什么选这个方法：该标的H1扣非利润为正且未触发恢复估值阻断，因此扣非EPS是最透明的筛选分母；情景桥接避免把一个H1数字直接当作全年点预测。
为什么选这个分母：使用H1扣非利润以减少非经常性收益污染；股本由总市值/当前价反推；H2/H1转换比例是明确的季节性和交付假设，不是机械年化。
为什么用这组倍数：PE区间按公司所属行业匹配，表达估值压缩/正常/重估三种情景；这是House筛选区间，正式升级前仍需当前可审计的Street区间。
为什么用这组概率：30\%/50\%/20\%是统一筛选先验：基准情景权重最高，同时保留熊市和牛市的实质性下行/上行贡献；它是House情景权重，不是经过统计估计的客观事件概率。
证据如何参与：官方半年报预告、Q1财务包和行情/K线提供事实输入；研报/来源用于分母或可比背景。没有正权重外部目标锚，House价值不伪装为Street共识。
\begin{align*}
\mathrm{Shares} &= \frac{547.53}{12.50}=43.8024\ \mathrm{100m\ shares}\\
\mathrm{EPS}_{H1}^d &= \frac{11.83}{43.8024}=0.2701\\
\mathrm{EPS}_s &= \mathrm{EPS}_{H1}^d(1+r_s)=0.45/0.51/0.59\\
V_s &= \mathrm{EPS}_s\times \mathrm{PE}_s=4.01/6.67/10.10\\
V_{prob} &= 0.30\times 4.01+0.50\times 6.67+0.20\times 10.10=6.56\\
P_{target} &= V_{prob}=6.56\quad(\text{外部锚权重}=0)\\
\mathrm{Upside} &= \frac{6.56}{12.50}-1=-47.5\%\\
\end{align*}

\paragraph{000987 越秀资本：金融 PB/ROE + 盈利混合}
分母：Q1 BPS plus H1 deducted EPS bridged to H2；证据质量：中；审计状态：PASS。
为什么选这个方法：金融公司主要通过净资产、ROE和盈利质量估值；PB保留资产负债表/ROE锚，PE提供独立盈利交叉检查，混合使用可以避免只依赖投资收益驱动的单季EPS。
为什么选这个分母：Q1 BPS作为资产分母，H1扣非EPS通过H2情景转换为熊/基准/牛盈利；这样将资产价值和盈利能力分开处理。
为什么用这组倍数：PB 0.75/1.05/1.35倍分别代表资产质量压力/正常/重估；行业PE区间提供盈利交叉检查；这些是House区间，不是外部目标价。
为什么用这组概率：30\%/50\%/20\%是统一筛选先验：基准情景权重最高，同时保留熊市和牛市的实质性下行/上行贡献；它是House情景权重，不是经过统计估计的客观事件概率。
证据如何参与：官方半年报预告、Q1财务包和行情/K线提供事实输入；研报/来源用于分母或可比背景。外部目标价9.44元，赋予正权重5\%。
\begin{align*}
\mathrm{EPS}_{H1}^d &= \frac{23.32}{50.1712} = 0.4648\ \mathrm{CNY/share}\\
V_s &= 0.65\left(\mathrm{BPS}_1\times \mathrm{PB}_s\right)+0.35\left(\mathrm{EPS}_s\times \mathrm{PE}_s\right)\\
V_{b/m/b} &= 5.45/8.11/10.95\ \mathrm{CNY/share}\quad(\mathrm{PB}=4.98/6.97/8.96,\ \mathrm{PE}=6.32/10.23/14.64)\\
V_{prob} &= 0.30\times 5.45+0.50\times 8.11+0.20\times 10.95=7.88\\
P_{target} &= 7.88\times 0.95+9.44\times 0.05=7.96\\
\mathrm{Upside} &= \frac{7.96}{7.77}-1=2.5\%\\
\end{align*}

\paragraph{000155 川能动力：常规扣非 EPS / PE}
分母：H1 deducted EPS bridged with H2/H1 conversion；证据质量：高；审计状态：PASS。
为什么选这个方法：该标的H1扣非利润为正且未触发恢复估值阻断，因此扣非EPS是最透明的筛选分母；情景桥接避免把一个H1数字直接当作全年点预测。
为什么选这个分母：使用H1扣非利润以减少非经常性收益污染；股本由总市值/当前价反推；H2/H1转换比例是明确的季节性和交付假设，不是机械年化。
为什么用这组倍数：PE区间按公司所属行业匹配，表达估值压缩/正常/重估三种情景；这是House筛选区间，正式升级前仍需当前可审计的Street区间。
为什么用这组概率：30\%/50\%/20\%是统一筛选先验：基准情景权重最高，同时保留熊市和牛市的实质性下行/上行贡献；它是House情景权重，不是经过统计估计的客观事件概率。
证据如何参与：官方半年报预告、Q1财务包和行情/K线提供事实输入；研报/来源用于分母或可比背景。外部目标价20.43元，赋予正权重10\%。
\begin{align*}
\mathrm{Shares} &= \frac{204.74}{11.09}=18.4617\ \mathrm{100m\ shares}\\
\mathrm{EPS}_{H1}^d &= \frac{6.75}{18.4617}=0.3656\\
\mathrm{EPS}_s &= \mathrm{EPS}_{H1}^d(1+r_s)=0.60/0.69/0.80\\
V_s &= \mathrm{EPS}_s\times \mathrm{PE}_s=5.43/9.03/13.67\\
V_{prob} &= 0.30\times 5.43+0.50\times 9.03+0.20\times 13.67=8.88\\
P_{target} &= 8.88\times 0.90+20.43\times 0.10=10.04\\
\mathrm{Upside} &= \frac{10.04}{11.09}-1=-9.5\%\\
\end{align*}

\paragraph{601377 兴业证券：金融 PB/ROE + 盈利混合}
分母：Q1 BPS plus H1 deducted EPS bridged to H2；证据质量：中高；审计状态：PASS。
为什么选这个方法：金融公司主要通过净资产、ROE和盈利质量估值；PB保留资产负债表/ROE锚，PE提供独立盈利交叉检查，混合使用可以避免只依赖投资收益驱动的单季EPS。
为什么选这个分母：Q1 BPS作为资产分母，H1扣非EPS通过H2情景转换为熊/基准/牛盈利；这样将资产价值和盈利能力分开处理。
为什么用这组倍数：PB 0.75/1.05/1.35倍分别代表资产质量压力/正常/重估；行业PE区间提供盈利交叉检查；这些是House区间，不是外部目标价。
为什么用这组概率：30\%/50\%/20\%是统一筛选先验：基准情景权重最高，同时保留熊市和牛市的实质性下行/上行贡献；它是House情景权重，不是经过统计估计的客观事件概率。
证据如何参与：官方半年报预告、Q1财务包和行情/K线提供事实输入；研报/来源用于分母或可比背景。没有正权重外部目标锚，House价值不伪装为Street共识。
\begin{align*}
\mathrm{EPS}_{H1}^d &= \frac{20.95}{86.3591} = 0.2426\ \mathrm{CNY/share}\\
V_s &= 0.65\left(\mathrm{BPS}_1\times \mathrm{PB}_s\right)+0.35\left(\mathrm{EPS}_s\times \mathrm{PE}_s\right)\\
V_{b/m/b} &= 4.35/6.34/8.43\ \mathrm{CNY/share}\quad(\mathrm{PB}=4.92/6.89/8.85,\ \mathrm{PE}=3.30/5.34/7.64)\\
V_{prob} &= 0.30\times 4.35+0.50\times 6.34+0.20\times 8.43=6.16\\
P_{target} &= V_{prob}=6.16\quad(\text{外部锚权重}=0)\\
\mathrm{Upside} &= \frac{6.16}{6.21}-1=-0.8\%\\
\end{align*}

\paragraph{600120 浙江东方：金融 PB/ROE + 盈利混合}
分母：Q1 BPS plus H1 deducted EPS bridged to H2；证据质量：中低；审计状态：PASS。
为什么选这个方法：金融公司主要通过净资产、ROE和盈利质量估值；PB保留资产负债表/ROE锚，PE提供独立盈利交叉检查，混合使用可以避免只依赖投资收益驱动的单季EPS。
为什么选这个分母：Q1 BPS作为资产分母，H1扣非EPS通过H2情景转换为熊/基准/牛盈利；这样将资产价值和盈利能力分开处理。
为什么用这组倍数：PB 0.75/1.05/1.35倍分别代表资产质量压力/正常/重估；行业PE区间提供盈利交叉检查；这些是House区间，不是外部目标价。
为什么用这组概率：30\%/50\%/20\%是统一筛选先验：基准情景权重最高，同时保留熊市和牛市的实质性下行/上行贡献；它是House情景权重，不是经过统计估计的客观事件概率。
证据如何参与：官方半年报预告、Q1财务包和行情/K线提供事实输入；研报/来源用于分母或可比背景。没有正权重外部目标锚，House价值不伪装为Street共识。
\begin{align*}
\mathrm{EPS}_{H1}^d &= \frac{8.40}{34.1548} = 0.2459\ \mathrm{CNY/share}\\
V_s &= 0.65\left(\mathrm{BPS}_1\times \mathrm{PB}_s\right)+0.35\left(\mathrm{EPS}_s\times \mathrm{PE}_s\right)\\
V_{b/m/b} &= 3.51/5.17/6.93\ \mathrm{CNY/share}\quad(\mathrm{PB}=3.60/5.04/6.48,\ \mathrm{PE}=3.34/5.41/7.75)\\
V_{prob} &= 0.30\times 3.51+0.50\times 5.17+0.20\times 6.93=5.02\\
P_{target} &= V_{prob}=5.02\quad(\text{外部锚权重}=0)\\
\mathrm{Upside} &= \frac{5.02}{4.91}-1=2.2\%\\
\end{align*}

\paragraph{000703 恒逸石化：周期调整扣非 EPS / PE}
分母：H1 deducted EPS bridged with H2/H1 conversion；证据质量：高；审计状态：PASS。
为什么选这个方法：石油石化盈利受商品价格、运价、价差或利用率周期影响，因此使用周期调整EPS/PE区间，而不是把某一个季度的PE当作全年估值。H2转换比例用于检验当前周期能否延续。
为什么选这个分母：使用H1扣非利润以减少非经常性收益污染；股本由总市值/当前价反推；H2/H1转换比例是明确的季节性和交付假设，不是机械年化。
为什么用这组倍数：熊/基准/牛PE区间是针对该周期业务类型的House假设，不是外部券商目标价；熊值较低反映周期反转，牛值较高必须由价格、销量和现金流共同确认。
为什么用这组概率：30\%/50\%/20\%是统一筛选先验：基准情景权重最高，同时保留熊市和牛市的实质性下行/上行贡献；它是House情景权重，不是经过统计估计的客观事件概率。
证据如何参与：官方半年报预告、Q1财务包和行情/K线提供事实输入；研报/来源用于分母或可比背景。外部目标价17.05元，赋予正权重10\%。
\begin{align*}
\mathrm{Shares} &= \frac{553.36}{14.48}=38.2155\ \mathrm{100m\ shares}\\
\mathrm{EPS}_{H1}^d &= \frac{57.20}{38.2155}=1.4968\\
\mathrm{EPS}_s &= \mathrm{EPS}_{H1}^d(1+r_s)=2.32/2.84/3.37\\
V_s &= \mathrm{EPS}_s\times \mathrm{PE}_s=10.86/25.59/40.41\\
V_{prob} &= 0.30\times 10.86+0.50\times 25.59+0.20\times 40.41=24.13\\
P_{target} &= 24.13\times 0.90+17.05\times 0.10=23.42\\
\mathrm{Upside} &= \frac{23.42}{14.48}-1=61.7\%\\
\end{align*}

\paragraph{000301 东方盛虹：周期调整扣非 EPS / PE}
分母：H1 deducted EPS bridged with H2/H1 conversion；证据质量：中高；审计状态：PASS。
为什么选这个方法：石油石化盈利受商品价格、运价、价差或利用率周期影响，因此使用周期调整EPS/PE区间，而不是把某一个季度的PE当作全年估值。H2转换比例用于检验当前周期能否延续。
为什么选这个分母：使用H1扣非利润以减少非经常性收益污染；股本由总市值/当前价反推；H2/H1转换比例是明确的季节性和交付假设，不是机械年化。
为什么用这组倍数：熊/基准/牛PE区间是针对该周期业务类型的House假设，不是外部券商目标价；熊值较低反映周期反转，牛值较高必须由价格、销量和现金流共同确认。
为什么用这组概率：30\%/50\%/20\%是统一筛选先验：基准情景权重最高，同时保留熊市和牛市的实质性下行/上行贡献；它是House情景权重，不是经过统计估计的客观事件概率。
证据如何参与：官方半年报预告、Q1财务包和行情/K线提供事实输入；研报/来源用于分母或可比背景。没有正权重外部目标锚，House价值不伪装为Street共识。
\begin{align*}
\mathrm{Shares} &= \frac{767.56}{11.61}=66.1120\ \mathrm{100m\ shares}\\
\mathrm{EPS}_{H1}^d &= \frac{44.16}{66.1120}=0.6680\\
\mathrm{EPS}_s &= \mathrm{EPS}_{H1}^d(1+r_s)=1.04/1.27/1.50\\
V_s &= \mathrm{EPS}_s\times \mathrm{PE}_s=6.21/11.42/18.03\\
V_{prob} &= 0.30\times 6.21+0.50\times 11.42+0.20\times 18.03=11.18\\
P_{target} &= V_{prob}=11.18\quad(\text{外部锚权重}=0)\\
\mathrm{Upside} &= \frac{11.18}{11.61}-1=-3.7\%\\
\end{align*}

\paragraph{001339 智微智能：成长盈利情景 PE}
分母：H1 deducted EPS bridged with H2/H1 conversion；证据质量：中高；审计状态：PASS。
为什么选这个方法：计算机具有成长或盈利持续期特征，因此模型给盈利桥接配置较高PE区间，而不是直接给主题收入估值；价格和现金流检查防止把泛化成长叙事直接变成EPS。
为什么选这个分母：使用H1扣非利润以减少非经常性收益污染；股本由总市值/当前价反推；H2/H1转换比例是明确的季节性和交付假设，不是机械年化。
为什么用这组倍数：PE区间是House对成长持续期的情景判断；牛市情景要求业绩桥接和经营证据持续，不是市场一致目标价。
为什么用这组概率：30\%/50\%/20\%是统一筛选先验：基准情景权重最高，同时保留熊市和牛市的实质性下行/上行贡献；它是House情景权重，不是经过统计估计的客观事件概率。
证据如何参与：官方半年报预告、Q1财务包和行情/K线提供事实输入；研报/来源用于分母或可比背景。没有正权重外部目标锚，House价值不伪装为Street共识。
\begin{align*}
\mathrm{Shares} &= \frac{292.59}{89.21}=3.2798\ \mathrm{100m\ shares}\\
\mathrm{EPS}_{H1}^d &= \frac{3.67}{3.2798}=1.1199\\
\mathrm{EPS}_s &= \mathrm{EPS}_{H1}^d(1+r_s)=1.90/2.48/2.69\\
V_s &= \mathrm{EPS}_s\times \mathrm{PE}_s=34.27/69.44/102.13\\
V_{prob} &= 0.30\times 34.27+0.50\times 69.44+0.20\times 102.13=65.43\\
P_{target} &= V_{prob}=65.43\quad(\text{外部锚权重}=0)\\
\mathrm{Upside} &= \frac{65.43}{89.21}-1=-26.7\%\\
\end{align*}

\paragraph{002558 巨人网络：成长盈利情景 PE}
分母：H1 deducted EPS bridged with H2/H1 conversion；证据质量：中高；审计状态：PASS。
为什么选这个方法：传媒具有成长或盈利持续期特征，因此模型给盈利桥接配置较高PE区间，而不是直接给主题收入估值；价格和现金流检查防止把泛化成长叙事直接变成EPS。
为什么选这个分母：使用H1扣非利润以减少非经常性收益污染；股本由总市值/当前价反推；H2/H1转换比例是明确的季节性和交付假设，不是机械年化。
为什么用这组倍数：PE区间是House对成长持续期的情景判断；牛市情景要求业绩桥接和经营证据持续，不是市场一致目标价。
为什么用这组概率：30\%/50\%/20\%是统一筛选先验：基准情景权重最高，同时保留熊市和牛市的实质性下行/上行贡献；它是House情景权重，不是经过统计估计的客观事件概率。
证据如何参与：官方半年报预告、Q1财务包和行情/K线提供事实输入；研报/来源用于分母或可比背景。没有正权重外部目标锚，House价值不伪装为Street共识。
\begin{align*}
\mathrm{Shares} &= \frac{561.81}{29.56}=19.0058\ \mathrm{100m\ shares}\\
\mathrm{EPS}_{H1}^d &= \frac{23.00}{19.0058}=1.2102\\
\mathrm{EPS}_s &= \mathrm{EPS}_{H1}^d(1+r_s)=2.06/2.48/2.90\\
V_s &= \mathrm{EPS}_s\times \mathrm{PE}_s=20.57/34.73/52.28\\
V_{prob} &= 0.30\times 20.57+0.50\times 34.73+0.20\times 52.28=33.99\\
P_{target} &= V_{prob}=33.99\quad(\text{外部锚权重}=0)\\
\mathrm{Upside} &= \frac{33.99}{29.56}-1=15.0\%\\
\end{align*}

\paragraph{301308 江波龙：成长盈利情景 PE}
分母：H1 deducted EPS bridged with H2/H1 conversion；证据质量：中高；审计状态：PASS。
为什么选这个方法：电子具有成长或盈利持续期特征，因此模型给盈利桥接配置较高PE区间，而不是直接给主题收入估值；价格和现金流检查防止把泛化成长叙事直接变成EPS。
为什么选这个分母：使用H1扣非利润以减少非经常性收益污染；股本由总市值/当前价反推；H2/H1转换比例是明确的季节性和交付假设，不是机械年化。
为什么用这组倍数：PE区间是House对成长持续期的情景判断；牛市情景要求业绩桥接和经营证据持续，不是市场一致目标价。
为什么用这组概率：30\%/50\%/20\%是统一筛选先验：基准情景权重最高，同时保留熊市和牛市的实质性下行/上行贡献；它是House情景权重，不是经过统计估计的客观事件概率。
证据如何参与：官方半年报预告、Q1财务包和行情/K线提供事实输入；研报/来源用于分母或可比背景。没有正权重外部目标锚，House价值不伪装为Street共识。
\begin{align*}
\mathrm{Shares} &= \frac{2064.54}{488.00}=4.2306\ \mathrm{100m\ shares}\\
\mathrm{EPS}_{H1}^d &= \frac{97.50}{4.2306}=23.0463\\
\mathrm{EPS}_s &= \mathrm{EPS}_{H1}^d(1+r_s)=39.18/46.09/53.01\\
V_s &= \mathrm{EPS}_s\times \mathrm{PE}_s=366.00/1198.41/1802.22\\
V_{prob} &= 0.30\times 366.00+0.50\times 1198.41+0.20\times 1802.22=1069.45\\
P_{target} &= V_{prob}=1069.45\quad(\text{外部锚权重}=0)\\
\mathrm{Upside} &= \frac{1069.45}{488.00}-1=119.2\%\\
\end{align*}

\paragraph{601628 中国人寿：金融 PB/ROE + 盈利混合}
分母：Q1 BPS plus H1 deducted EPS bridged to H2；证据质量：中高；审计状态：PASS。
为什么选这个方法：金融公司主要通过净资产、ROE和盈利质量估值；PB保留资产负债表/ROE锚，PE提供独立盈利交叉检查，混合使用可以避免只依赖投资收益驱动的单季EPS。
为什么选这个分母：Q1 BPS作为资产分母，H1扣非EPS通过H2情景转换为熊/基准/牛盈利；这样将资产价值和盈利能力分开处理。
为什么用这组倍数：PB 0.75/1.05/1.35倍分别代表资产质量压力/正常/重估；行业PE区间提供盈利交叉检查；这些是House区间，不是外部目标价。
为什么用这组概率：30\%/50\%/20\%是统一筛选先验：基准情景权重最高，同时保留熊市和牛市的实质性下行/上行贡献；它是House情景权重，不是经过统计估计的客观事件概率。
证据如何参与：官方半年报预告、Q1财务包和行情/K线提供事实输入；研报/来源用于分母或可比背景。没有正权重外部目标锚，House价值不伪装为Street共识。
\begin{align*}
\mathrm{EPS}_{H1}^d &= \frac{1332.76}{282.6470} = 4.7153\ \mathrm{CNY/share}\\
V_s &= 0.65\left(\mathrm{BPS}_1\times \mathrm{PB}_s\right)+0.35\left(\mathrm{EPS}_s\times \mathrm{PE}_s\right)\\
V_{b/m/b} &= 32.76/50.75/70.56\ \mathrm{CNY/share}\quad(\mathrm{PB}=15.87/22.22/28.57,\ \mathrm{PE}=64.13/103.74/148.53)\\
V_{prob} &= 0.30\times 30.32+0.50\times 50.75+0.20\times 70.56=48.58\\
P_{target} &= V_{prob}=48.58\quad(\text{外部锚权重}=0)\\
\mathrm{Upside} &= \frac{48.58}{40.42}-1=20.2\%\\
\end{align*}

\paragraph{002493 荣盛石化：周期调整扣非 EPS / PE}
分母：H1 deducted EPS bridged with H2/H1 conversion；证据质量：中高；审计状态：PASS。
为什么选这个方法：石油石化盈利受商品价格、运价、价差或利用率周期影响，因此使用周期调整EPS/PE区间，而不是把某一个季度的PE当作全年估值。H2转换比例用于检验当前周期能否延续。
为什么选这个分母：使用H1扣非利润以减少非经常性收益污染；股本由总市值/当前价反推；H2/H1转换比例是明确的季节性和交付假设，不是机械年化。
为什么用这组倍数：熊/基准/牛PE区间是针对该周期业务类型的House假设，不是外部券商目标价；熊值较低反映周期反转，牛值较高必须由价格、销量和现金流共同确认。
为什么用这组概率：30\%/50\%/20\%是统一筛选先验：基准情景权重最高，同时保留熊市和牛市的实质性下行/上行贡献；它是House情景权重，不是经过统计估计的客观事件概率。
证据如何参与：官方半年报预告、Q1财务包和行情/K线提供事实输入；研报/来源用于分母或可比背景。没有正权重外部目标锚，House价值不伪装为Street共识。
\begin{align*}
\mathrm{Shares} &= \frac{1108.83}{11.10}=99.8946\ \mathrm{100m\ shares}\\
\mathrm{EPS}_{H1}^d &= \frac{53.00}{99.8946}=0.5306\\
\mathrm{EPS}_s &= \mathrm{EPS}_{H1}^d(1+r_s)=0.82/1.01/1.19\\
V_s &= \mathrm{EPS}_s\times \mathrm{PE}_s=4.93/9.07/14.33\\
V_{prob} &= 0.30\times 4.93+0.50\times 9.07+0.20\times 14.33=8.88\\
P_{target} &= V_{prob}=8.88\quad(\text{外部锚权重}=0)\\
\mathrm{Upside} &= \frac{8.88}{11.10}-1=-20.0\%\\
\end{align*}

\paragraph{001309 德明利：存储周期前瞻 EPS / PE}
分母：2026E memory-cycle-adjusted EPS, not H1 annualized EPS；证据质量：中高；审计状态：PASS。
为什么选这个方法：最新H1利润存在低基数、扭亏或周期扭曲，因此先正常化EPS，再套业务匹配PE区间；业务匹配可比包括江波龙、兆易创新、北京君正。
为什么选这个分母：正常化分母为：2026E memory-cycle-adjusted EPS, not H1 annualized EPS；H1机械年化指标只作为市场隐含压力测试，不作为最终目标价分母。
为什么用这组倍数：熊/基准/牛倍数分别表达下行、正常兑现和上行重估；区间通过可比方法和当前价格隐含指标交叉检查（存储价格和企业级产品景气强，但周期反转后利润回撤风险高），因此只能作为条件性House价值。
为什么用这组概率：30\%/50\%/20\%是统一筛选先验：基准情景权重最高，同时保留熊市和牛市的实质性下行/上行贡献；它是House情景权重，不是经过统计估计的客观事件概率。
证据如何参与：官方半年报预告、Q1财务包和行情/K线提供事实输入；研报/来源用于分母或可比背景。没有正权重外部目标锚，House价值不伪装为Street共识。
\begin{align*}
V_s=\mathrm{EPS}_{normalized,s}\times \mathrm{PE}_s\\
V_{b/m/b} &= 330.00/513.00/630.00\ \mathrm{CNY/share}\\
V_{prob} &= 0.30\times 330.00+0.50\times 513.00+0.20\times 630.00=481.50\\
P_{target} &= V_{prob}=481.50\quad(\text{外部锚权重}=0)\\
\mathrm{Upside} &= \frac{481.50}{662.40}-1=-27.3\%\\
\end{align*}

\paragraph{601233 桐昆股份：周期调整扣非 EPS / PE}
分母：H1 deducted EPS bridged with H2/H1 conversion；证据质量：中高；审计状态：PASS。
为什么选这个方法：石油石化盈利受商品价格、运价、价差或利用率周期影响，因此使用周期调整EPS/PE区间，而不是把某一个季度的PE当作全年估值。H2转换比例用于检验当前周期能否延续。
为什么选这个分母：使用H1扣非利润以减少非经常性收益污染；股本由总市值/当前价反推；H2/H1转换比例是明确的季节性和交付假设，不是机械年化。
为什么用这组倍数：熊/基准/牛PE区间是针对该周期业务类型的House假设，不是外部券商目标价；熊值较低反映周期反转，牛值较高必须由价格、销量和现金流共同确认。
为什么用这组概率：30\%/50\%/20\%是统一筛选先验：基准情景权重最高，同时保留熊市和牛市的实质性下行/上行贡献；它是House情景权重，不是经过统计估计的客观事件概率。
证据如何参与：官方半年报预告、Q1财务包和行情/K线提供事实输入；研报/来源用于分母或可比背景。没有正权重外部目标锚，House价值不伪装为Street共识。
\begin{align*}
\mathrm{Shares} &= \frac{492.93}{20.72}=23.7901\ \mathrm{100m\ shares}\\
\mathrm{EPS}_{H1}^d &= \frac{40.80}{23.7901}=1.7150\\
\mathrm{EPS}_s &= \mathrm{EPS}_{H1}^d(1+r_s)=2.66/3.26/3.86\\
V_s &= \mathrm{EPS}_s\times \mathrm{PE}_s=15.54/29.33/46.31\\
V_{prob} &= 0.30\times 15.54+0.50\times 29.33+0.20\times 46.31=28.59\\
P_{target} &= V_{prob}=28.59\quad(\text{外部锚权重}=0)\\
\mathrm{Upside} &= \frac{28.59}{20.72}-1=38.0\%\\
\end{align*}

\paragraph{603986 兆易创新：成长盈利情景 PE}
分母：H1 deducted EPS bridged with H2/H1 conversion；证据质量：中高；审计状态：PASS。
为什么选这个方法：电子具有成长或盈利持续期特征，因此模型给盈利桥接配置较高PE区间，而不是直接给主题收入估值；价格和现金流检查防止把泛化成长叙事直接变成EPS。
为什么选这个分母：使用H1扣非利润以减少非经常性收益污染；股本由总市值/当前价反推；H2/H1转换比例是明确的季节性和交付假设，不是机械年化。
为什么用这组倍数：PE区间是House对成长持续期的情景判断；牛市情景要求业绩桥接和经营证据持续，不是市场一致目标价。
为什么用这组概率：30\%/50\%/20\%是统一筛选先验：基准情景权重最高，同时保留熊市和牛市的实质性下行/上行贡献；它是House情景权重，不是经过统计估计的客观事件概率。
证据如何参与：官方半年报预告、Q1财务包和行情/K线提供事实输入；研报/来源用于分母或可比背景。没有正权重外部目标锚，House价值不伪装为Street共识。
\begin{align*}
\mathrm{Shares} &= \frac{4012.51}{571.79}=7.0175\ \mathrm{100m\ shares}\\
\mathrm{EPS}_{H1}^d &= \frac{48.50}{7.0175}=6.9113\\
\mathrm{EPS}_s &= \mathrm{EPS}_{H1}^d(1+r_s)=11.75/13.82/15.90\\
V_s &= \mathrm{EPS}_s\times \mathrm{PE}_s=211.49/359.39/540.47\\
V_{prob} &= 0.30\times 211.49+0.50\times 359.39+0.20\times 540.47=351.24\\
P_{target} &= V_{prob}=351.24\quad(\text{外部锚权重}=0)\\
\mathrm{Upside} &= \frac{351.24}{571.79}-1=-38.6\%\\
\end{align*}

\paragraph{601225 陕西煤业：周期调整扣非 EPS / PE}
分母：H1 deducted EPS bridged with H2/H1 conversion；证据质量：中高；审计状态：PASS。
为什么选这个方法：煤炭盈利受商品价格、运价、价差或利用率周期影响，因此使用周期调整EPS/PE区间，而不是把某一个季度的PE当作全年估值。H2转换比例用于检验当前周期能否延续。
为什么选这个分母：使用H1扣非利润以减少非经常性收益污染；股本由总市值/当前价反推；H2/H1转换比例是明确的季节性和交付假设，不是机械年化。
为什么用这组倍数：熊/基准/牛PE区间是针对该周期业务类型的House假设，不是外部券商目标价；熊值较低反映周期反转，牛值较高必须由价格、销量和现金流共同确认。
为什么用这组概率：30\%/50\%/20\%是统一筛选先验：基准情景权重最高，同时保留熊市和牛市的实质性下行/上行贡献；它是House情景权重，不是经过统计估计的客观事件概率。
证据如何参与：官方半年报预告、Q1财务包和行情/K线提供事实输入；研报/来源用于分母或可比背景。没有正权重外部目标锚，House价值不伪装为Street共识。
\begin{align*}
\mathrm{Shares} &= \frac{2385.94}{24.61}=96.9500\ \mathrm{100m\ shares}\\
\mathrm{EPS}_{H1}^d &= \frac{97.46}{96.9500}=1.0053\\
\mathrm{EPS}_s &= \mathrm{EPS}_{H1}^d(1+r_s)=1.56/2.08/2.26\\
V_s &= \mathrm{EPS}_s\times \mathrm{PE}_s=10.91/18.72/24.88\\
V_{prob} &= 0.30\times 10.91+0.50\times 18.72+0.20\times 24.88=17.61\\
P_{target} &= V_{prob}=17.61\quad(\text{外部锚权重}=0)\\
\mathrm{Upside} &= \frac{17.61}{24.61}-1=-28.4\%\\
\end{align*}

\paragraph{601138 工业富联：成长盈利情景 PE}
分母：H1 deducted EPS bridged with H2/H1 conversion；证据质量：中高；审计状态：PASS。
为什么选这个方法：电子具有成长或盈利持续期特征，因此模型给盈利桥接配置较高PE区间，而不是直接给主题收入估值；价格和现金流检查防止把泛化成长叙事直接变成EPS。
为什么选这个分母：使用H1扣非利润以减少非经常性收益污染；股本由总市值/当前价反推；H2/H1转换比例是明确的季节性和交付假设，不是机械年化。
为什么用这组倍数：PE区间是House对成长持续期的情景判断；牛市情景要求业绩桥接和经营证据持续，不是市场一致目标价。
为什么用这组概率：30\%/50\%/20\%是统一筛选先验：基准情景权重最高，同时保留熊市和牛市的实质性下行/上行贡献；它是House情景权重，不是经过统计估计的客观事件概率。
证据如何参与：官方半年报预告、Q1财务包和行情/K线提供事实输入；研报/来源用于分母或可比背景。没有正权重外部目标锚，House价值不伪装为Street共识。
\begin{align*}
\mathrm{Shares} &= \frac{12624.81}{63.62}=198.4409\ \mathrm{100m\ shares}\\
\mathrm{EPS}_{H1}^d &= \frac{232.00}{198.4409}=1.1691\\
\mathrm{EPS}_s &= \mathrm{EPS}_{H1}^d(1+r_s)=1.99/3.06/2.69\\
V_s &= \mathrm{EPS}_s\times \mathrm{PE}_s=35.77/79.56/91.42\\
V_{prob} &= 0.30\times 35.77+0.50\times 79.56+0.20\times 91.42=68.80\\
P_{target} &= V_{prob}=68.80\quad(\text{外部锚权重}=0)\\
\mathrm{Upside} &= \frac{68.80}{63.62}-1=8.1\%\\
\end{align*}

\paragraph{000807 云铝股份：周期调整扣非 EPS / PE}
分母：H1 deducted EPS bridged with H2/H1 conversion；证据质量：中高；审计状态：PASS。
为什么选这个方法：有色金属盈利受商品价格、运价、价差或利用率周期影响，因此使用周期调整EPS/PE区间，而不是把某一个季度的PE当作全年估值。H2转换比例用于检验当前周期能否延续。
为什么选这个分母：使用H1扣非利润以减少非经常性收益污染；股本由总市值/当前价反推；H2/H1转换比例是明确的季节性和交付假设，不是机械年化。
为什么用这组倍数：熊/基准/牛PE区间是针对该周期业务类型的House假设，不是外部券商目标价；熊值较低反映周期反转，牛值较高必须由价格、销量和现金流共同确认。
为什么用这组概率：30\%/50\%/20\%是统一筛选先验：基准情景权重最高，同时保留熊市和牛市的实质性下行/上行贡献；它是House情景权重，不是经过统计估计的客观事件概率。
证据如何参与：官方半年报预告、Q1财务包和行情/K线提供事实输入；研报/来源用于分母或可比背景。没有正权重外部目标锚，House价值不伪装为Street共识。
\begin{align*}
\mathrm{Shares} &= \frac{905.14}{26.10}=34.6797\ \mathrm{100m\ shares}\\
\mathrm{EPS}_{H1}^d &= \frac{75.50}{34.6797}=2.1771\\
\mathrm{EPS}_s &= \mathrm{EPS}_{H1}^d(1+r_s)=3.37/4.03/4.68\\
V_s &= \mathrm{EPS}_s\times \mathrm{PE}_s=19.58/48.33/74.89\\
V_{prob} &= 0.30\times 19.58+0.50\times 48.33+0.20\times 74.89=45.02\\
P_{target} &= V_{prob}=45.02\quad(\text{外部锚权重}=0)\\
\mathrm{Upside} &= \frac{45.02}{26.10}-1=72.5\%\\
\end{align*}

\paragraph{603993 洛阳钼业：周期调整扣非 EPS / PE}
分母：H1 deducted EPS bridged with H2/H1 conversion；证据质量：中高；审计状态：PASS。
为什么选这个方法：有色金属盈利受商品价格、运价、价差或利用率周期影响，因此使用周期调整EPS/PE区间，而不是把某一个季度的PE当作全年估值。H2转换比例用于检验当前周期能否延续。
为什么选这个分母：使用H1扣非利润以减少非经常性收益污染；股本由总市值/当前价反推；H2/H1转换比例是明确的季节性和交付假设，不是机械年化。
为什么用这组倍数：熊/基准/牛PE区间是针对该周期业务类型的House假设，不是外部券商目标价；熊值较低反映周期反转，牛值较高必须由价格、销量和现金流共同确认。
为什么用这组概率：30\%/50\%/20\%是统一筛选先验：基准情景权重最高，同时保留熊市和牛市的实质性下行/上行贡献；它是House情景权重，不是经过统计估计的客观事件概率。
证据如何参与：官方半年报预告、Q1财务包和行情/K线提供事实输入；研报/来源用于分母或可比背景。没有正权重外部目标锚，House价值不伪装为Street共识。
\begin{align*}
\mathrm{Shares} &= \frac{3831.72}{17.91}=213.9430\ \mathrm{100m\ shares}\\
\mathrm{EPS}_{H1}^d &= \frac{155.00}{213.9430}=0.7245\\
\mathrm{EPS}_s &= \mathrm{EPS}_{H1}^d(1+r_s)=1.12/1.46/1.56\\
V_s &= \mathrm{EPS}_s\times \mathrm{PE}_s=8.98/17.52/24.92\\
V_{prob} &= 0.30\times 8.98+0.50\times 17.52+0.20\times 24.92=16.44\\
P_{target} &= V_{prob}=16.44\quad(\text{外部锚权重}=0)\\
\mathrm{Upside} &= \frac{16.44}{17.91}-1=-8.2\%\\
\end{align*}

\paragraph{601872 招商轮船：周期调整扣非 EPS / PE}
分母：H1 deducted EPS bridged with H2/H1 conversion；证据质量：中高；审计状态：PASS。
为什么选这个方法：交通运输盈利受商品价格、运价、价差或利用率周期影响，因此使用周期调整EPS/PE区间，而不是把某一个季度的PE当作全年估值。H2转换比例用于检验当前周期能否延续。
为什么选这个分母：使用H1扣非利润以减少非经常性收益污染；股本由总市值/当前价反推；H2/H1转换比例是明确的季节性和交付假设，不是机械年化。
为什么用这组倍数：熊/基准/牛PE区间是针对该周期业务类型的House假设，不是外部券商目标价；熊值较低反映周期反转，牛值较高必须由价格、销量和现金流共同确认。
为什么用这组概率：30\%/50\%/20\%是统一筛选先验：基准情景权重最高，同时保留熊市和牛市的实质性下行/上行贡献；它是House情景权重，不是经过统计估计的客观事件概率。
证据如何参与：官方半年报预告、Q1财务包和行情/K线提供事实输入；研报/来源用于分母或可比背景。没有正权重外部目标锚，House价值不伪装为Street共识。
\begin{align*}
\mathrm{Shares} &= \frac{1232.17}{15.26}=80.7451\ \mathrm{100m\ shares}\\
\mathrm{EPS}_{H1}^d &= \frac{68.90}{80.7451}=0.8533\\
\mathrm{EPS}_s &= \mathrm{EPS}_{H1}^d(1+r_s)=1.32/2.09/1.83\\
V_s &= \mathrm{EPS}_s\times \mathrm{PE}_s=10.58/25.08/29.35\\
V_{prob} &= 0.30\times 10.58+0.50\times 25.08+0.20\times 29.35=21.58\\
P_{target} &= V_{prob}=21.58\quad(\text{外部锚权重}=0)\\
\mathrm{Upside} &= \frac{21.58}{15.26}-1=41.4\%\\
\end{align*}

\paragraph{600999 招商证券：金融 PB/ROE + 盈利混合}
分母：Q1 BPS plus H1 deducted EPS bridged to H2；证据质量：中高；审计状态：PASS。
为什么选这个方法：金融公司主要通过净资产、ROE和盈利质量估值；PB保留资产负债表/ROE锚，PE提供独立盈利交叉检查，混合使用可以避免只依赖投资收益驱动的单季EPS。
为什么选这个分母：Q1 BPS作为资产分母，H1扣非EPS通过H2情景转换为熊/基准/牛盈利；这样将资产价值和盈利能力分开处理。
为什么用这组倍数：PB 0.75/1.05/1.35倍分别代表资产质量压力/正常/重估；行业PE区间提供盈利交叉检查；这些是House区间，不是外部目标价。
为什么用这组概率：30\%/50\%/20\%是统一筛选先验：基准情景权重最高，同时保留熊市和牛市的实质性下行/上行贡献；它是House情景权重，不是经过统计估计的客观事件概率。
证据如何参与：官方半年报预告、Q1财务包和行情/K线提供事实输入；研报/来源用于分母或可比背景。没有正权重外部目标锚，House价值不伪装为Street共识。
\begin{align*}
\mathrm{EPS}_{H1}^d &= \frac{105.00}{86.9651} = 1.2074\ \mathrm{CNY/share}\\
V_s &= 0.65\left(\mathrm{BPS}_1\times \mathrm{PB}_s\right)+0.35\left(\mathrm{EPS}_s\times \mathrm{PE}_s\right)\\
V_{b/m/b} &= 12.72/19.05/25.86\ \mathrm{CNY/share}\quad(\mathrm{PB}=10.72/15.01/19.30,\ \mathrm{PE}=16.42/26.56/38.03)\\
V_{prob} &= 0.30\times 12.72+0.50\times 19.05+0.20\times 25.86=18.51\\
P_{target} &= V_{prob}=18.51\quad(\text{外部锚权重}=0)\\
\mathrm{Upside} &= \frac{18.51}{20.90}-1=-11.4\%\\
\end{align*}

\paragraph{601899 紫金矿业：周期调整扣非 EPS / PE}
分母：H1 deducted EPS bridged with H2/H1 conversion；证据质量：中高；审计状态：PASS。
为什么选这个方法：有色金属盈利受商品价格、运价、价差或利用率周期影响，因此使用周期调整EPS/PE区间，而不是把某一个季度的PE当作全年估值。H2转换比例用于检验当前周期能否延续。
为什么选这个分母：使用H1扣非利润以减少非经常性收益污染；股本由总市值/当前价反推；H2/H1转换比例是明确的季节性和交付假设，不是机械年化。
为什么用这组倍数：熊/基准/牛PE区间是针对该周期业务类型的House假设，不是外部券商目标价；熊值较低反映周期反转，牛值较高必须由价格、销量和现金流共同确认。
为什么用这组概率：30\%/50\%/20\%是统一筛选先验：基准情景权重最高，同时保留熊市和牛市的实质性下行/上行贡献；它是House情景权重，不是经过统计估计的客观事件概率。
证据如何参与：官方半年报预告、Q1财务包和行情/K线提供事实输入；研报/来源用于分母或可比背景。没有正权重外部目标锚，House价值不伪装为Street共识。
\begin{align*}
\mathrm{Shares} &= \frac{7666.10}{28.83}=265.9070\ \mathrm{100m\ shares}\\
\mathrm{EPS}_{H1}^d &= \frac{379.00}{265.9070}=1.4253\\
\mathrm{EPS}_s &= \mathrm{EPS}_{H1}^d(1+r_s)=2.21/3.34/3.06\\
V_s &= \mathrm{EPS}_s\times \mathrm{PE}_s=17.67/40.08/49.03\\
V_{prob} &= 0.30\times 17.67+0.50\times 40.08+0.20\times 49.03=35.15\\
P_{target} &= V_{prob}=35.15\quad(\text{外部锚权重}=0)\\
\mathrm{Upside} &= \frac{35.15}{28.83}-1=21.9\%\\
\end{align*}

\paragraph{600309 万华化学：周期调整扣非 EPS / PE}
分母：H1 deducted EPS bridged with H2/H1 conversion；证据质量：中高；审计状态：PASS。
为什么选这个方法：基础化工盈利受商品价格、运价、价差或利用率周期影响，因此使用周期调整EPS/PE区间，而不是把某一个季度的PE当作全年估值。H2转换比例用于检验当前周期能否延续。
为什么选这个分母：使用H1扣非利润以减少非经常性收益污染；股本由总市值/当前价反推；H2/H1转换比例是明确的季节性和交付假设，不是机械年化。
为什么用这组倍数：熊/基准/牛PE区间是针对该周期业务类型的House假设，不是外部券商目标价；熊值较低反映周期反转，牛值较高必须由价格、销量和现金流共同确认。
为什么用这组概率：30\%/50\%/20\%是统一筛选先验：基准情景权重最高，同时保留熊市和牛市的实质性下行/上行贡献；它是House情景权重，不是经过统计估计的客观事件概率。
证据如何参与：官方半年报预告、Q1财务包和行情/K线提供事实输入；研报/来源用于分母或可比背景。没有正权重外部目标锚，House价值不伪装为Street共识。
\begin{align*}
\mathrm{Shares} &= \frac{2194.15}{70.09}=31.3048\ \mathrm{100m\ shares}\\
\mathrm{EPS}_{H1}^d &= \frac{99.00}{31.3048}=3.1625\\
\mathrm{EPS}_s &= \mathrm{EPS}_{H1}^d(1+r_s)=4.90/5.85/6.80\\
V_s &= \mathrm{EPS}_s\times \mathrm{PE}_s=49.02/81.91/122.39\\
V_{prob} &= 0.30\times 49.02+0.50\times 81.91+0.20\times 122.39=80.14\\
P_{target} &= V_{prob}=80.14\quad(\text{外部锚权重}=0)\\
\mathrm{Upside} &= \frac{80.14}{70.09}-1=14.3\%\\
\end{align*}

\paragraph{600989 宝丰能源：周期调整扣非 EPS / PE}
分母：H1 deducted EPS bridged with H2/H1 conversion；证据质量：高；审计状态：PASS。
为什么选这个方法：基础化工盈利受商品价格、运价、价差或利用率周期影响，因此使用周期调整EPS/PE区间，而不是把某一个季度的PE当作全年估值。H2转换比例用于检验当前周期能否延续。
为什么选这个分母：使用H1扣非利润以减少非经常性收益污染；股本由总市值/当前价反推；H2/H1转换比例是明确的季节性和交付假设，不是机械年化。
为什么用这组倍数：熊/基准/牛PE区间是针对该周期业务类型的House假设，不是外部券商目标价；熊值较低反映周期反转，牛值较高必须由价格、销量和现金流共同确认。
为什么用这组概率：30\%/50\%/20\%是统一筛选先验：基准情景权重最高，同时保留熊市和牛市的实质性下行/上行贡献；它是House情景权重，不是经过统计估计的客观事件概率。
证据如何参与：官方半年报预告、Q1财务包和行情/K线提供事实输入；研报/来源用于分母或可比背景。外部目标价39.60元，赋予正权重10\%。
\begin{align*}
\mathrm{Shares} &= \frac{1618.47}{22.07}=73.3335\ \mathrm{100m\ shares}\\
\mathrm{EPS}_{H1}^d &= \frac{91.50}{73.3335}=1.2477\\
\mathrm{EPS}_s &= \mathrm{EPS}_{H1}^d(1+r_s)=1.93/2.31/2.68\\
V_s &= \mathrm{EPS}_s\times \mathrm{PE}_s=16.55/32.32/48.29\\
V_{prob} &= 0.30\times 16.55+0.50\times 32.32+0.20\times 48.29=30.78\\
P_{target} &= 30.78\times 0.90+39.60\times 0.10=31.66\\
\mathrm{Upside} &= \frac{31.66}{22.07}-1=43.5\%\\
\end{align*}

\paragraph{300475 香农芯创：成长盈利情景 PE}
分母：H1 deducted EPS bridged with H2/H1 conversion；证据质量：中；审计状态：PASS。
为什么选这个方法：电子具有成长或盈利持续期特征，因此模型给盈利桥接配置较高PE区间，而不是直接给主题收入估值；价格和现金流检查防止把泛化成长叙事直接变成EPS。
为什么选这个分母：使用H1扣非利润以减少非经常性收益污染；股本由总市值/当前价反推；H2/H1转换比例是明确的季节性和交付假设，不是机械年化。
为什么用这组倍数：PE区间是House对成长持续期的情景判断；牛市情景要求业绩桥接和经营证据持续，不是市场一致目标价。
为什么用这组概率：30\%/50\%/20\%是统一筛选先验：基准情景权重最高，同时保留熊市和牛市的实质性下行/上行贡献；它是House情景权重，不是经过统计估计的客观事件概率。
证据如何参与：官方半年报预告、Q1财务包和行情/K线提供事实输入；研报/来源用于分母或可比背景。没有正权重外部目标锚，House价值不伪装为Street共识。
\begin{align*}
\mathrm{Shares} &= \frac{1038.63}{221.20}=4.6954\ \mathrm{100m\ shares}\\
\mathrm{EPS}_{H1}^d &= \frac{36.40}{4.6954}=7.7522\\
\mathrm{EPS}_s &= \mathrm{EPS}_{H1}^d(1+r_s)=13.18/15.50/17.83\\
V_s &= \mathrm{EPS}_s\times \mathrm{PE}_s=165.90/403.12/606.22\\
V_{prob} &= 0.30\times 165.90+0.50\times 403.12+0.20\times 606.22=372.57\\
P_{target} &= V_{prob}=372.57\quad(\text{外部锚权重}=0)\\
\mathrm{Upside} &= \frac{372.57}{221.20}-1=68.4\%\\
\end{align*}

\paragraph{000776 广发证券：金融 PB/ROE + 盈利混合}
分母：Q1 BPS plus H1 deducted EPS bridged to H2；证据质量：中高；审计状态：PASS。
为什么选这个方法：金融公司主要通过净资产、ROE和盈利质量估值；PB保留资产负债表/ROE锚，PE提供独立盈利交叉检查，混合使用可以避免只依赖投资收益驱动的单季EPS。
为什么选这个分母：Q1 BPS作为资产分母，H1扣非EPS通过H2情景转换为熊/基准/牛盈利；这样将资产价值和盈利能力分开处理。
为什么用这组倍数：PB 0.75/1.05/1.35倍分别代表资产质量压力/正常/重估；行业PE区间提供盈利交叉检查；这些是House区间，不是外部目标价。
为什么用这组概率：30\%/50\%/20\%是统一筛选先验：基准情景权重最高，同时保留熊市和牛市的实质性下行/上行贡献；它是House情景权重，不是经过统计估计的客观事件概率。
证据如何参与：官方半年报预告、Q1财务包和行情/K线提供事实输入；研报/来源用于分母或可比背景。没有正权重外部目标锚，House价值不伪装为Street共识。
\begin{align*}
\mathrm{EPS}_{H1}^d &= \frac{116.90}{78.2485} = 1.4940\ \mathrm{CNY/share}\\
V_s &= 0.65\left(\mathrm{BPS}_1\times \mathrm{PB}_s\right)+0.35\left(\mathrm{EPS}_s\times \mathrm{PE}_s\right)\\
V_{b/m/b} &= 15.69/23.52/31.92\ \mathrm{CNY/share}\quad(\mathrm{PB}=13.20/18.49/23.77,\ \mathrm{PE}=20.32/32.87/47.06)\\
V_{prob} &= 0.30\times 15.69+0.50\times 23.52+0.20\times 31.92=22.85\\
P_{target} &= V_{prob}=22.85\quad(\text{外部锚权重}=0)\\
\mathrm{Upside} &= \frac{22.85}{22.78}-1=0.3\%\\
\end{align*}

\paragraph{600030 中信证券：金融 PB/ROE + 盈利混合}
分母：Q1 BPS plus H1 deducted EPS bridged to H2；证据质量：中高；审计状态：PASS。
为什么选这个方法：金融公司主要通过净资产、ROE和盈利质量估值；PB保留资产负债表/ROE锚，PE提供独立盈利交叉检查，混合使用可以避免只依赖投资收益驱动的单季EPS。
为什么选这个分母：Q1 BPS作为资产分母，H1扣非EPS通过H2情景转换为熊/基准/牛盈利；这样将资产价值和盈利能力分开处理。
为什么用这组倍数：PB 0.75/1.05/1.35倍分别代表资产质量压力/正常/重估；行业PE区间提供盈利交叉检查；这些是House区间，不是外部目标价。
为什么用这组概率：30\%/50\%/20\%是统一筛选先验：基准情景权重最高，同时保留熊市和牛市的实质性下行/上行贡献；它是House情景权重，不是经过统计估计的客观事件概率。
证据如何参与：官方半年报预告、Q1财务包和行情/K线提供事实输入；研报/来源用于分母或可比背景。没有正权重外部目标锚，House价值不伪装为Street共识。
\begin{align*}
\mathrm{EPS}_{H1}^d &= \frac{236.10}{148.2056} = 1.5931\ \mathrm{CNY/share}\\
V_s &= 0.65\left(\mathrm{BPS}_1\times \mathrm{PB}_s\right)+0.35\left(\mathrm{EPS}_s\times \mathrm{PE}_s\right)\\
V_{b/m/b} &= 17.17/25.69/34.82\ \mathrm{CNY/share}\quad(\mathrm{PB}=14.75/20.65/26.55,\ \mathrm{PE}=21.67/35.05/50.18)\\
V_{prob} &= 0.30\times 17.17+0.50\times 25.69+0.20\times 34.82=24.96\\
P_{target} &= V_{prob}=24.96\quad(\text{外部锚权重}=0)\\
\mathrm{Upside} &= \frac{24.96}{28.60}-1=-12.7\%\\
\end{align*}

\paragraph{601088 中国神华：周期调整扣非 EPS / PE}
分母：H1 deducted EPS bridged with H2/H1 conversion；证据质量：中；审计状态：PASS。
为什么选这个方法：煤炭盈利受商品价格、运价、价差或利用率周期影响，因此使用周期调整EPS/PE区间，而不是把某一个季度的PE当作全年估值。H2转换比例用于检验当前周期能否延续。
为什么选这个分母：使用H1扣非利润以减少非经常性收益污染；股本由总市值/当前价反推；H2/H1转换比例是明确的季节性和交付假设，不是机械年化。
为什么用这组倍数：熊/基准/牛PE区间是针对该周期业务类型的House假设，不是外部券商目标价；熊值较低反映周期反转，牛值较高必须由价格、销量和现金流共同确认。
为什么用这组概率：30\%/50\%/20\%是统一筛选先验：基准情景权重最高，同时保留熊市和牛市的实质性下行/上行贡献；它是House情景权重，不是经过统计估计的客观事件概率。
证据如何参与：官方半年报预告、Q1财务包和行情/K线提供事实输入；研报/来源用于分母或可比背景。没有正权重外部目标锚，House价值不伪装为Street共识。
\begin{align*}
\mathrm{Shares} &= \frac{9476.11}{43.69}=216.8943\ \mathrm{100m\ shares}\\
\mathrm{EPS}_{H1}^d &= \frac{271.00}{216.8943}=1.2495\\
\mathrm{EPS}_s &= \mathrm{EPS}_{H1}^d(1+r_s)=1.94/2.54/2.81\\
V_s &= \mathrm{EPS}_s\times \mathrm{PE}_s=13.56/22.86/30.92\\
V_{prob} &= 0.30\times 13.56+0.50\times 22.86+0.20\times 30.92=21.68\\
P_{target} &= V_{prob}=21.68\quad(\text{外部锚权重}=0)\\
\mathrm{Upside} &= \frac{21.68}{43.69}-1=-50.4\%\\
\end{align*}

\paragraph{600362 江西铜业：周期调整扣非 EPS / PE}
分母：H1 deducted EPS bridged with H2/H1 conversion；证据质量：中；审计状态：PASS。
为什么选这个方法：有色金属盈利受商品价格、运价、价差或利用率周期影响，因此使用周期调整EPS/PE区间，而不是把某一个季度的PE当作全年估值。H2转换比例用于检验当前周期能否延续。
为什么选这个分母：使用H1扣非利润以减少非经常性收益污染；股本由总市值/当前价反推；H2/H1转换比例是明确的季节性和交付假设，不是机械年化。
为什么用这组倍数：熊/基准/牛PE区间是针对该周期业务类型的House假设，不是外部券商目标价；熊值较低反映周期反转，牛值较高必须由价格、销量和现金流共同确认。
为什么用这组概率：30\%/50\%/20\%是统一筛选先验：基准情景权重最高，同时保留熊市和牛市的实质性下行/上行贡献；它是House情景权重，不是经过统计估计的客观事件概率。
证据如何参与：官方半年报预告、Q1财务包和行情/K线提供事实输入；研报/来源用于分母或可比背景。没有正权重外部目标锚，House价值不伪装为Street共识。
\begin{align*}
\mathrm{Shares} &= \frac{1464.04}{42.28}=34.6272\ \mathrm{100m\ shares}\\
\mathrm{EPS}_{H1}^d &= \frac{81.05}{34.6272}=2.3406\\
\mathrm{EPS}_s &= \mathrm{EPS}_{H1}^d(1+r_s)=3.63/4.33/5.03\\
V_s &= \mathrm{EPS}_s\times \mathrm{PE}_s=29.02/51.96/80.52\\
V_{prob} &= 0.30\times 29.02+0.50\times 51.96+0.20\times 80.52=50.79\\
P_{target} &= V_{prob}=50.79\quad(\text{外部锚权重}=0)\\
\mathrm{Upside} &= \frac{50.79}{42.28}-1=20.1\%\\
\end{align*}

\paragraph{002709 天赐材料：成长盈利情景 PE}
分母：H1 deducted EPS bridged with H2/H1 conversion；证据质量：高；审计状态：PASS。
为什么选这个方法：电力设备具有成长或盈利持续期特征，因此模型给盈利桥接配置较高PE区间，而不是直接给主题收入估值；价格和现金流检查防止把泛化成长叙事直接变成EPS。
为什么选这个分母：使用H1扣非利润以减少非经常性收益污染；股本由总市值/当前价反推；H2/H1转换比例是明确的季节性和交付假设，不是机械年化。
为什么用这组倍数：PE区间是House对成长持续期的情景判断；牛市情景要求业绩桥接和经营证据持续，不是市场一致目标价。
为什么用这组概率：30\%/50\%/20\%是统一筛选先验：基准情景权重最高，同时保留熊市和牛市的实质性下行/上行贡献；它是House情景权重，不是经过统计估计的客观事件概率。
证据如何参与：官方半年报预告、Q1财务包和行情/K线提供事实输入；研报/来源用于分母或可比背景。外部目标价66.40元，赋予正权重10\%。
\begin{align*}
\mathrm{Shares} &= \frac{765.48}{37.55}=20.3856\ \mathrm{100m\ shares}\\
\mathrm{EPS}_{H1}^d &= \frac{28.00}{20.3856}=1.3735\\
\mathrm{EPS}_s &= \mathrm{EPS}_{H1}^d(1+r_s)=2.33/3.06/3.16\\
V_s &= \mathrm{EPS}_s\times \mathrm{PE}_s=28.16/67.32/88.45\\
V_{prob} &= 0.30\times 28.16+0.50\times 67.32+0.20\times 88.45=59.80\\
P_{target} &= 59.80\times 0.90+66.40\times 0.10=60.46\\
\mathrm{Upside} &= \frac{60.46}{37.55}-1=61.0\%\\
\end{align*}

\paragraph{002414 高德红外：成长盈利情景 PE}
分母：H1 deducted EPS bridged with H2/H1 conversion；证据质量：中；审计状态：PASS。
为什么选这个方法：国防军工具有成长或盈利持续期特征，因此模型给盈利桥接配置较高PE区间，而不是直接给主题收入估值；价格和现金流检查防止把泛化成长叙事直接变成EPS。
为什么选这个分母：使用H1扣非利润以减少非经常性收益污染；股本由总市值/当前价反推；H2/H1转换比例是明确的季节性和交付假设，不是机械年化。
为什么用这组倍数：PE区间是House对成长持续期的情景判断；牛市情景要求业绩桥接和经营证据持续，不是市场一致目标价。
为什么用这组概率：30\%/50\%/20\%是统一筛选先验：基准情景权重最高，同时保留熊市和牛市的实质性下行/上行贡献；它是House情景权重，不是经过统计估计的客观事件概率。
证据如何参与：官方半年报预告、Q1财务包和行情/K线提供事实输入；研报/来源用于分母或可比背景。外部目标价20.00元，赋予正权重5\%。
\begin{align*}
\mathrm{Shares} &= \frac{568.86}{13.32}=42.7072\ \mathrm{100m\ shares}\\
\mathrm{EPS}_{H1}^d &= \frac{13.25}{42.7072}=0.3103\\
\mathrm{EPS}_s &= \mathrm{EPS}_{H1}^d(1+r_s)=0.53/0.64/0.74\\
V_s &= \mathrm{EPS}_s\times \mathrm{PE}_s=9.99/22.26/33.51\\
V_{prob} &= 0.30\times 9.99+0.50\times 22.26+0.20\times 33.51=20.83\\
P_{target} &= 20.83\times 0.95+20.00\times 0.05=20.79\\
\mathrm{Upside} &= \frac{20.79}{13.32}-1=56.1\%\\
\end{align*}

\paragraph{300390 天华新能：周期调整 EPS / PE}
分母：2026E cycle-adjusted EPS and lithium-price sensitivity；证据质量：中高；审计状态：PASS。
为什么选这个方法：最新H1利润存在低基数、扭亏或周期扭曲，因此先正常化EPS，再套业务匹配PE区间；业务匹配可比包括赣锋锂业、盛新锂能、天齐锂业。
为什么选这个分母：正常化分母为：2026E cycle-adjusted EPS and lithium-price sensitivity；H1机械年化指标只作为市场隐含压力测试，不作为最终目标价分母。
为什么用这组倍数：熊/基准/牛倍数分别表达下行、正常兑现和上行重估；区间通过可比方法和当前价格隐含指标交叉检查（锂盐价格恢复和大客户绑定带来弹性，但周期延续性仍需验证），因此只能作为条件性House价值。
为什么用这组概率：30\%/50\%/20\%是统一筛选先验：基准情景权重最高，同时保留熊市和牛市的实质性下行/上行贡献；它是House情景权重，不是经过统计估计的客观事件概率。
证据如何参与：官方半年报预告、Q1财务包和行情/K线提供事实输入；研报/来源用于分母或可比背景。没有正权重外部目标锚，House价值不伪装为Street共识。
\begin{align*}
V_s=\mathrm{EPS}_{normalized,s}\times \mathrm{PE}_s\\
V_{b/m/b} &= 56.40/79.80/102.40\ \mathrm{CNY/share}\\
V_{prob} &= 0.30\times 56.40+0.50\times 79.80+0.20\times 102.40=77.30\\
P_{target} &= V_{prob}=77.30\quad(\text{外部锚权重}=0)\\
\mathrm{Upside} &= \frac{77.30}{60.01}-1=28.8\%\\
\end{align*}

\paragraph{002648 卫星化学：周期调整扣非 EPS / PE}
分母：H1 deducted EPS bridged with H2/H1 conversion；证据质量：中高；审计状态：PASS。
为什么选这个方法：基础化工盈利受商品价格、运价、价差或利用率周期影响，因此使用周期调整EPS/PE区间，而不是把某一个季度的PE当作全年估值。H2转换比例用于检验当前周期能否延续。
为什么选这个分母：使用H1扣非利润以减少非经常性收益污染；股本由总市值/当前价反推；H2/H1转换比例是明确的季节性和交付假设，不是机械年化。
为什么用这组倍数：熊/基准/牛PE区间是针对该周期业务类型的House假设，不是外部券商目标价；熊值较低反映周期反转，牛值较高必须由价格、销量和现金流共同确认。
为什么用这组概率：30\%/50\%/20\%是统一筛选先验：基准情景权重最高，同时保留熊市和牛市的实质性下行/上行贡献；它是House情景权重，不是经过统计估计的客观事件概率。
证据如何参与：官方半年报预告、Q1财务包和行情/K线提供事实输入；研报/来源用于分母或可比背景。没有正权重外部目标锚，House价值不伪装为Street共识。
\begin{align*}
\mathrm{Shares} &= \frac{789.27}{23.43}=33.6863\ \mathrm{100m\ shares}\\
\mathrm{EPS}_{H1}^d &= \frac{61.87}{33.6863}=1.8367\\
\mathrm{EPS}_s &= \mathrm{EPS}_{H1}^d(1+r_s)=2.85/3.43/3.95\\
V_s &= \mathrm{EPS}_s\times \mathrm{PE}_s=17.57/48.02/71.08\\
V_{prob} &= 0.30\times 17.57+0.50\times 48.02+0.20\times 71.08=43.50\\
P_{target} &= V_{prob}=43.50\quad(\text{外部锚权重}=0)\\
\mathrm{Upside} &= \frac{43.50}{23.43}-1=85.7\%\\
\end{align*}

\paragraph{000792 盐湖股份：周期调整扣非 EPS / PE}
分母：H1 deducted EPS bridged with H2/H1 conversion；证据质量：中；审计状态：PASS。
为什么选这个方法：基础化工盈利受商品价格、运价、价差或利用率周期影响，因此使用周期调整EPS/PE区间，而不是把某一个季度的PE当作全年估值。H2转换比例用于检验当前周期能否延续。
为什么选这个分母：使用H1扣非利润以减少非经常性收益污染；股本由总市值/当前价反推；H2/H1转换比例是明确的季节性和交付假设，不是机械年化。
为什么用这组倍数：熊/基准/牛PE区间是针对该周期业务类型的House假设，不是外部券商目标价；熊值较低反映周期反转，牛值较高必须由价格、销量和现金流共同确认。
为什么用这组概率：30\%/50\%/20\%是统一筛选先验：基准情景权重最高，同时保留熊市和牛市的实质性下行/上行贡献；它是House情景权重，不是经过统计估计的客观事件概率。
证据如何参与：官方半年报预告、Q1财务包和行情/K线提供事实输入；研报/来源用于分母或可比背景。没有正权重外部目标锚，House价值不伪装为Street共识。
\begin{align*}
\mathrm{Shares} &= \frac{1337.18}{25.27}=52.9157\ \mathrm{100m\ shares}\\
\mathrm{EPS}_{H1}^d &= \frac{61.50}{52.9157}=1.1622\\
\mathrm{EPS}_s &= \mathrm{EPS}_{H1}^d(1+r_s)=1.80/2.15/2.50\\
V_s &= \mathrm{EPS}_s\times \mathrm{PE}_s=18.01/30.10/44.98\\
V_{prob} &= 0.30\times 18.01+0.50\times 30.10+0.20\times 44.98=29.45\\
P_{target} &= V_{prob}=29.45\quad(\text{外部锚权重}=0)\\
\mathrm{Upside} &= \frac{29.45}{25.27}-1=16.5\%\\
\end{align*}

\paragraph{601869 长飞光纤：成长盈利情景 PE}
分母：H1 deducted EPS bridged with H2/H1 conversion；证据质量：中；审计状态：PASS。
为什么选这个方法：通信具有成长或盈利持续期特征，因此模型给盈利桥接配置较高PE区间，而不是直接给主题收入估值；价格和现金流检查防止把泛化成长叙事直接变成EPS。
为什么选这个分母：使用H1扣非利润以减少非经常性收益污染；股本由总市值/当前价反推；H2/H1转换比例是明确的季节性和交付假设，不是机械年化。
为什么用这组倍数：PE区间是House对成长持续期的情景判断；牛市情景要求业绩桥接和经营证据持续，不是市场一致目标价。
为什么用这组概率：30\%/50\%/20\%是统一筛选先验：基准情景权重最高，同时保留熊市和牛市的实质性下行/上行贡献；它是House情景权重，不是经过统计估计的客观事件概率。
证据如何参与：官方半年报预告、Q1财务包和行情/K线提供事实输入；研报/来源用于分母或可比背景。没有正权重外部目标锚，House价值不伪装为Street共识。
\begin{align*}
\mathrm{Shares} &= \frac{3522.32}{425.45}=8.2790\ \mathrm{100m\ shares}\\
\mathrm{EPS}_{H1}^d &= \frac{23.00}{8.2790}=2.7781\\
\mathrm{EPS}_s &= \mathrm{EPS}_{H1}^d(1+r_s)=4.72/6.05/6.39\\
V_s &= \mathrm{EPS}_s\times \mathrm{PE}_s=85.01/157.30/217.25\\
V_{prob} &= 0.30\times 85.01+0.50\times 157.30+0.20\times 217.25=147.60\\
P_{target} &= V_{prob}=147.60\quad(\text{外部锚权重}=0)\\
\mathrm{Upside} &= \frac{147.60}{425.45}-1=-65.3\%\\
\end{align*}

\paragraph{002432 九安医疗：成长盈利情景 PE}
分母：H1 deducted EPS bridged with H2/H1 conversion；证据质量：中；审计状态：PASS。
为什么选这个方法：医药生物具有成长或盈利持续期特征，因此模型给盈利桥接配置较高PE区间，而不是直接给主题收入估值；价格和现金流检查防止把泛化成长叙事直接变成EPS。
为什么选这个分母：使用H1扣非利润以减少非经常性收益污染；股本由总市值/当前价反推；H2/H1转换比例是明确的季节性和交付假设，不是机械年化。
为什么用这组倍数：PE区间是House对成长持续期的情景判断；牛市情景要求业绩桥接和经营证据持续，不是市场一致目标价。
为什么用这组概率：30\%/50\%/20\%是统一筛选先验：基准情景权重最高，同时保留熊市和牛市的实质性下行/上行贡献；它是House情景权重，不是经过统计估计的客观事件概率。
证据如何参与：官方半年报预告、Q1财务包和行情/K线提供事实输入；研报/来源用于分母或可比背景。没有正权重外部目标锚，House价值不伪装为Street共识。
\begin{align*}
\mathrm{Shares} &= \frac{312.10}{66.99}=4.6589\ \mathrm{100m\ shares}\\
\mathrm{EPS}_{H1}^d &= \frac{31.50}{4.6589}=6.7612\\
\mathrm{EPS}_s &= \mathrm{EPS}_{H1}^d(1+r_s)=11.49/13.52/15.55\\
V_s &= \mathrm{EPS}_s\times \mathrm{PE}_s=50.24/338.06/497.63\\
V_{prob} &= 0.30\times 50.24+0.50\times 338.06+0.20\times 497.63=283.63\\
P_{target} &= V_{prob}=283.63\quad(\text{外部锚权重}=0)\\
\mathrm{Upside} &= \frac{283.63}{66.99}-1=323.4\%\\
\end{align*}

\paragraph{000657 中钨高新：周期调整扣非 EPS / PE}
分母：H1 deducted EPS bridged with H2/H1 conversion；证据质量：中高；审计状态：PASS。
为什么选这个方法：有色金属盈利受商品价格、运价、价差或利用率周期影响，因此使用周期调整EPS/PE区间，而不是把某一个季度的PE当作全年估值。H2转换比例用于检验当前周期能否延续。
为什么选这个分母：使用H1扣非利润以减少非经常性收益污染；股本由总市值/当前价反推；H2/H1转换比例是明确的季节性和交付假设，不是机械年化。
为什么用这组倍数：熊/基准/牛PE区间是针对该周期业务类型的House假设，不是外部券商目标价；熊值较低反映周期反转，牛值较高必须由价格、销量和现金流共同确认。
为什么用这组概率：30\%/50\%/20\%是统一筛选先验：基准情景权重最高，同时保留熊市和牛市的实质性下行/上行贡献；它是House情景权重，不是经过统计估计的客观事件概率。
证据如何参与：官方半年报预告、Q1财务包和行情/K线提供事实输入；研报/来源用于分母或可比背景。没有正权重外部目标锚，House价值不伪装为Street共识。
\begin{align*}
\mathrm{Shares} &= \frac{1590.47}{69.80}=22.7861\ \mathrm{100m\ shares}\\
\mathrm{EPS}_{H1}^d &= \frac{20.55}{22.7861}=0.9019\\
\mathrm{EPS}_s &= \mathrm{EPS}_{H1}^d(1+r_s)=1.40/2.25/1.94\\
V_s &= \mathrm{EPS}_s\times \mathrm{PE}_s=11.18/27.00/31.02\\
V_{prob} &= 0.30\times 11.18+0.50\times 27.00+0.20\times 31.02=23.06\\
P_{target} &= V_{prob}=23.06\quad(\text{外部锚权重}=0)\\
\mathrm{Upside} &= \frac{23.06}{69.80}-1=-67.0\%\\
\end{align*}

\paragraph{601995 中金公司：金融 PB/ROE + 盈利混合}
分母：Q1 BPS plus H1 deducted EPS bridged to H2；证据质量：中高；审计状态：PASS。
为什么选这个方法：金融公司主要通过净资产、ROE和盈利质量估值；PB保留资产负债表/ROE锚，PE提供独立盈利交叉检查，混合使用可以避免只依赖投资收益驱动的单季EPS。
为什么选这个分母：Q1 BPS作为资产分母，H1扣非EPS通过H2情景转换为熊/基准/牛盈利；这样将资产价值和盈利能力分开处理。
为什么用这组倍数：PB 0.75/1.05/1.35倍分别代表资产质量压力/正常/重估；行业PE区间提供盈利交叉检查；这些是House区间，不是外部目标价。
为什么用这组概率：30\%/50\%/20\%是统一筛选先验：基准情景权重最高，同时保留熊市和牛市的实质性下行/上行贡献；它是House情景权重，不是经过统计估计的客观事件概率。
证据如何参与：官方半年报预告、Q1财务包和行情/K线提供事实输入；研报/来源用于分母或可比背景。没有正权重外部目标锚，House价值不伪装为Street共识。
\begin{align*}
\mathrm{EPS}_{H1}^d &= \frac{78.07}{48.2726} = 1.6173\ \mathrm{CNY/share}\\
V_s &= 0.65\left(\mathrm{BPS}_1\times \mathrm{PB}_s\right)+0.35\left(\mathrm{EPS}_s\times \mathrm{PE}_s\right)\\
V_{b/m/b} &= 18.07/26.97/36.49\ \mathrm{CNY/share}\quad(\mathrm{PB}=15.95/22.33/28.71,\ \mathrm{PE}=22.00/35.58/50.94)\\
V_{prob} &= 0.30\times 18.07+0.50\times 26.97+0.20\times 36.49=26.20\\
P_{target} &= V_{prob}=26.20\quad(\text{外部锚权重}=0)\\
\mathrm{Upside} &= \frac{26.20}{38.34}-1=-31.7\%\\
\end{align*}

\paragraph{002497 雅化集团：周期调整扣非 EPS / PE}
分母：H1 deducted EPS bridged with H2/H1 conversion；证据质量：高；审计状态：PASS。
为什么选这个方法：基础化工盈利受商品价格、运价、价差或利用率周期影响，因此使用周期调整EPS/PE区间，而不是把某一个季度的PE当作全年估值。H2转换比例用于检验当前周期能否延续。
为什么选这个分母：使用H1扣非利润以减少非经常性收益污染；股本由总市值/当前价反推；H2/H1转换比例是明确的季节性和交付假设，不是机械年化。
为什么用这组倍数：熊/基准/牛PE区间是针对该周期业务类型的House假设，不是外部券商目标价；熊值较低反映周期反转，牛值较高必须由价格、销量和现金流共同确认。
为什么用这组概率：30\%/50\%/20\%是统一筛选先验：基准情景权重最高，同时保留熊市和牛市的实质性下行/上行贡献；它是House情景权重，不是经过统计估计的客观事件概率。
证据如何参与：官方半年报预告、Q1财务包和行情/K线提供事实输入；研报/来源用于分母或可比背景。外部目标价42.00元，赋予正权重10\%。
\begin{align*}
\mathrm{Shares} &= \frac{208.27}{18.07}=11.5257\ \mathrm{100m\ shares}\\
\mathrm{EPS}_{H1}^d &= \frac{12.20}{11.5257}=1.0585\\
\mathrm{EPS}_s &= \mathrm{EPS}_{H1}^d(1+r_s)=1.64/2.63/2.28\\
V_s &= \mathrm{EPS}_s\times \mathrm{PE}_s=13.55/36.82/40.96\\
V_{prob} &= 0.30\times 13.55+0.50\times 36.82+0.20\times 40.96=30.67\\
P_{target} &= 30.67\times 0.90+42.00\times 0.10=31.80\\
\mathrm{Upside} &= \frac{31.80}{18.07}-1=76.0\%\\
\end{align*}

\paragraph{000933 神火股份：周期调整扣非 EPS / PE}
分母：H1 deducted EPS bridged with H2/H1 conversion；证据质量：中；审计状态：PASS。
为什么选这个方法：有色金属盈利受商品价格、运价、价差或利用率周期影响，因此使用周期调整EPS/PE区间，而不是把某一个季度的PE当作全年估值。H2转换比例用于检验当前周期能否延续。
为什么选这个分母：使用H1扣非利润以减少非经常性收益污染；股本由总市值/当前价反推；H2/H1转换比例是明确的季节性和交付假设，不是机械年化。
为什么用这组倍数：熊/基准/牛PE区间是针对该周期业务类型的House假设，不是外部券商目标价；熊值较低反映周期反转，牛值较高必须由价格、销量和现金流共同确认。
为什么用这组概率：30\%/50\%/20\%是统一筛选先验：基准情景权重最高，同时保留熊市和牛市的实质性下行/上行贡献；它是House情景权重，不是经过统计估计的客观事件概率。
证据如何参与：官方半年报预告、Q1财务包和行情/K线提供事实输入；研报/来源用于分母或可比背景。没有正权重外部目标锚，House价值不伪装为Street共识。
\begin{align*}
\mathrm{Shares} &= \frac{573.72}{25.51}=22.4900\ \mathrm{100m\ shares}\\
\mathrm{EPS}_{H1}^d &= \frac{48.06}{22.4900}=2.1369\\
\mathrm{EPS}_s &= \mathrm{EPS}_{H1}^d(1+r_s)=3.31/3.96/4.59\\
V_s &= \mathrm{EPS}_s\times \mathrm{PE}_s=19.13/47.54/73.51\\
V_{prob} &= 0.30\times 19.13+0.50\times 47.54+0.20\times 73.51=44.21\\
P_{target} &= V_{prob}=44.21\quad(\text{外部锚权重}=0)\\
\mathrm{Upside} &= \frac{44.21}{25.51}-1=73.3\%\\
\end{align*}

\paragraph{300223 北京君正：成长盈利情景 PE}
分母：H1 deducted EPS bridged with H2/H1 conversion；证据质量：中高；审计状态：PASS。
为什么选这个方法：电子具有成长或盈利持续期特征，因此模型给盈利桥接配置较高PE区间，而不是直接给主题收入估值；价格和现金流检查防止把泛化成长叙事直接变成EPS。
为什么选这个分母：使用H1扣非利润以减少非经常性收益污染；股本由总市值/当前价反推；H2/H1转换比例是明确的季节性和交付假设，不是机械年化。
为什么用这组倍数：PE区间是House对成长持续期的情景判断；牛市情景要求业绩桥接和经营证据持续，不是市场一致目标价。
为什么用这组概率：30\%/50\%/20\%是统一筛选先验：基准情景权重最高，同时保留熊市和牛市的实质性下行/上行贡献；它是House情景权重，不是经过统计估计的客观事件概率。
证据如何参与：官方半年报预告、Q1财务包和行情/K线提供事实输入；研报/来源用于分母或可比背景。没有正权重外部目标锚，House价值不伪装为Street共识。
\begin{align*}
\mathrm{Shares} &= \frac{919.83}{190.18}=4.8366\ \mathrm{100m\ shares}\\
\mathrm{EPS}_{H1}^d &= \frac{11.51}{4.8366}=2.3797\\
\mathrm{EPS}_s &= \mathrm{EPS}_{H1}^d(1+r_s)=4.05/4.76/5.47\\
V_s &= \mathrm{EPS}_s\times \mathrm{PE}_s=72.82/123.75/186.09\\
V_{prob} &= 0.30\times 72.82+0.50\times 123.75+0.20\times 186.09=120.94\\
P_{target} &= V_{prob}=120.94\quad(\text{外部锚权重}=0)\\
\mathrm{Upside} &= \frac{120.94}{190.18}-1=-36.4\%\\
\end{align*}

\paragraph{002192 融捷股份：周期调整扣非 EPS / PE}
分母：H1 deducted EPS bridged with H2/H1 conversion；证据质量：中；审计状态：PASS。
为什么选这个方法：有色金属盈利受商品价格、运价、价差或利用率周期影响，因此使用周期调整EPS/PE区间，而不是把某一个季度的PE当作全年估值。H2转换比例用于检验当前周期能否延续。
为什么选这个分母：使用H1扣非利润以减少非经常性收益污染；股本由总市值/当前价反推；H2/H1转换比例是明确的季节性和交付假设，不是机械年化。
为什么用这组倍数：熊/基准/牛PE区间是针对该周期业务类型的House假设，不是外部券商目标价；熊值较低反映周期反转，牛值较高必须由价格、销量和现金流共同确认。
为什么用这组概率：30\%/50\%/20\%是统一筛选先验：基准情景权重最高，同时保留熊市和牛市的实质性下行/上行贡献；它是House情景权重，不是经过统计估计的客观事件概率。
证据如何参与：官方半年报预告、Q1财务包和行情/K线提供事实输入；研报/来源用于分母或可比背景。没有正权重外部目标锚，House价值不伪装为Street共识。
\begin{align*}
\mathrm{Shares} &= \frac{172.46}{66.42}=2.5965\ \mathrm{100m\ shares}\\
\mathrm{EPS}_{H1}^d &= \frac{10.03}{2.5965}=3.8629\\
\mathrm{EPS}_s &= \mathrm{EPS}_{H1}^d(1+r_s)=5.99/7.15/8.31\\
V_s &= \mathrm{EPS}_s\times \mathrm{PE}_s=47.90/85.76/132.88\\
V_{prob} &= 0.30\times 47.90+0.50\times 85.76+0.20\times 132.88=83.83\\
P_{target} &= V_{prob}=83.83\quad(\text{外部锚权重}=0)\\
\mathrm{Upside} &= \frac{83.83}{66.42}-1=26.2\%\\
\end{align*}

\paragraph{601066 中信建投：金融 PB/ROE + 盈利混合}
分母：Q1 BPS plus H1 deducted EPS bridged to H2；证据质量：中高；审计状态：PASS。
为什么选这个方法：金融公司主要通过净资产、ROE和盈利质量估值；PB保留资产负债表/ROE锚，PE提供独立盈利交叉检查，混合使用可以避免只依赖投资收益驱动的单季EPS。
为什么选这个分母：Q1 BPS作为资产分母，H1扣非EPS通过H2情景转换为熊/基准/牛盈利；这样将资产价值和盈利能力分开处理。
为什么用这组倍数：PB 0.75/1.05/1.35倍分别代表资产质量压力/正常/重估；行业PE区间提供盈利交叉检查；这些是House区间，不是外部目标价。
为什么用这组概率：30\%/50\%/20\%是统一筛选先验：基准情景权重最高，同时保留熊市和牛市的实质性下行/上行贡献；它是House情景权重，不是经过统计估计的客观事件概率。
证据如何参与：官方半年报预告、Q1财务包和行情/K线提供事实输入；研报/来源用于分母或可比背景。没有正权重外部目标锚，House价值不伪装为Street共识。
\begin{align*}
\mathrm{EPS}_{H1}^d &= \frac{77.90}{77.5669} = 1.0043\ \mathrm{CNY/share}\\
V_s &= 0.65\left(\mathrm{BPS}_1\times \mathrm{PB}_s\right)+0.35\left(\mathrm{EPS}_s\times \mathrm{PE}_s\right)\\
V_{b/m/b} &= 10.20/15.32/20.83\ \mathrm{CNY/share}\quad(\mathrm{PB}=8.34/11.68/15.02,\ \mathrm{PE}=13.66/22.09/31.64)\\
V_{prob} &= 0.30\times 10.20+0.50\times 15.32+0.20\times 20.83=14.89\\
P_{target} &= V_{prob}=14.89\quad(\text{外部锚权重}=0)\\
\mathrm{Upside} &= \frac{14.89}{27.82}-1=-46.5\%\\
\end{align*}

\paragraph{002812 恩捷股份：产能周期正常化前瞻 EPS / PE}
分母：2027E normalized EPS after separator price and utilization recovery；证据质量：高；审计状态：PASS。
为什么选这个方法：最新H1利润存在低基数、扭亏或周期扭曲，因此先正常化EPS，再套业务匹配PE区间；业务匹配可比包括星源材质、中材科技、沧州明珠。
为什么选这个分母：正常化分母为：2027E normalized EPS after separator price and utilization recovery；H1机械年化指标只作为市场隐含压力测试，不作为最终目标价分母。
为什么用这组倍数：熊/基准/牛倍数分别表达下行、正常兑现和上行重估；区间通过可比方法和当前价格隐含指标交叉检查（隔膜供需改善和涨价预期强化，但2027E利润跳升尚需订单和价格验证），因此只能作为条件性House价值。
为什么用这组概率：30\%/50\%/20\%是统一筛选先验：基准情景权重最高，同时保留熊市和牛市的实质性下行/上行贡献；它是House情景权重，不是经过统计估计的客观事件概率。
证据如何参与：官方半年报预告、Q1财务包和行情/K线提供事实输入；研报/来源用于分母或可比背景。外部目标价103.00元，赋予正权重10\%。
\begin{align*}
V_s=\mathrm{EPS}_{normalized,s}\times \mathrm{PE}_s\\
V_{b/m/b} &= 21.60/103.20/126.20\ \mathrm{CNY/share}\\
V_{prob} &= 0.30\times 21.60+0.50\times 103.20+0.20\times 126.20=83.32\\
P_{target} &= 83.32\times 0.90+103.00\times 0.10=85.29\\
\mathrm{Upside} &= \frac{85.29}{54.96}-1=55.2\%\\
\end{align*}

\paragraph{002636 金安国纪：成长盈利情景 PE}
分母：H1 deducted EPS bridged with H2/H1 conversion；证据质量：中；审计状态：PASS。
为什么选这个方法：电子具有成长或盈利持续期特征，因此模型给盈利桥接配置较高PE区间，而不是直接给主题收入估值；价格和现金流检查防止把泛化成长叙事直接变成EPS。
为什么选这个分母：使用H1扣非利润以减少非经常性收益污染；股本由总市值/当前价反推；H2/H1转换比例是明确的季节性和交付假设，不是机械年化。
为什么用这组倍数：PE区间是House对成长持续期的情景判断；牛市情景要求业绩桥接和经营证据持续，不是市场一致目标价。
为什么用这组概率：30\%/50\%/20\%是统一筛选先验：基准情景权重最高，同时保留熊市和牛市的实质性下行/上行贡献；它是House情景权重，不是经过统计估计的客观事件概率。
证据如何参与：官方半年报预告、Q1财务包和行情/K线提供事实输入；研报/来源用于分母或可比背景。没有正权重外部目标锚，House价值不伪装为Street共识。
\begin{align*}
\mathrm{Shares} &= \frac{690.07}{94.79}=7.2800\ \mathrm{100m\ shares}\\
\mathrm{EPS}_{H1}^d &= \frac{7.80}{7.2800}=1.0721\\
\mathrm{EPS}_s &= \mathrm{EPS}_{H1}^d(1+r_s)=1.82/2.14/2.47\\
V_s &= \mathrm{EPS}_s\times \mathrm{PE}_s=32.81/55.75/83.84\\
V_{prob} &= 0.30\times 32.81+0.50\times 55.75+0.20\times 83.84=54.49\\
P_{target} &= V_{prob}=54.49\quad(\text{外部锚权重}=0)\\
\mathrm{Upside} &= \frac{54.49}{94.79}-1=-42.5\%\\
\end{align*}

\paragraph{002378 章源钨业：周期调整扣非 EPS / PE}
分母：H1 deducted EPS bridged with H2/H1 conversion；证据质量：中；审计状态：PASS。
为什么选这个方法：有色金属盈利受商品价格、运价、价差或利用率周期影响，因此使用周期调整EPS/PE区间，而不是把某一个季度的PE当作全年估值。H2转换比例用于检验当前周期能否延续。
为什么选这个分母：使用H1扣非利润以减少非经常性收益污染；股本由总市值/当前价反推；H2/H1转换比例是明确的季节性和交付假设，不是机械年化。
为什么用这组倍数：熊/基准/牛PE区间是针对该周期业务类型的House假设，不是外部券商目标价；熊值较低反映周期反转，牛值较高必须由价格、销量和现金流共同确认。
为什么用这组概率：30\%/50\%/20\%是统一筛选先验：基准情景权重最高，同时保留熊市和牛市的实质性下行/上行贡献；它是House情景权重，不是经过统计估计的客观事件概率。
证据如何参与：官方半年报预告、Q1财务包和行情/K线提供事实输入；研报/来源用于分母或可比背景。没有正权重外部目标锚，House价值不伪装为Street共识。
\begin{align*}
\mathrm{Shares} &= \frac{310.45}{25.84}=12.0143\ \mathrm{100m\ shares}\\
\mathrm{EPS}_{H1}^d &= \frac{6.80}{12.0143}=0.5660\\
\mathrm{EPS}_s &= \mathrm{EPS}_{H1}^d(1+r_s)=0.88/1.05/1.22\\
V_s &= \mathrm{EPS}_s\times \mathrm{PE}_s=7.02/12.57/19.47\\
V_{prob} &= 0.30\times 7.02+0.50\times 12.57+0.20\times 19.47=12.29\\
P_{target} &= V_{prob}=12.29\quad(\text{外部锚权重}=0)\\
\mathrm{Upside} &= \frac{12.29}{25.84}-1=-52.4\%\\
\end{align*}

\paragraph{601211 国泰海通：金融 PB/ROE + 盈利混合}
分母：Q1 BPS plus H1 deducted EPS bridged to H2；证据质量：中高；审计状态：PASS。
为什么选这个方法：金融公司主要通过净资产、ROE和盈利质量估值；PB保留资产负债表/ROE锚，PE提供独立盈利交叉检查，混合使用可以避免只依赖投资收益驱动的单季EPS。
为什么选这个分母：Q1 BPS作为资产分母，H1扣非EPS通过H2情景转换为熊/基准/牛盈利；这样将资产价值和盈利能力分开处理。
为什么用这组倍数：PB 0.75/1.05/1.35倍分别代表资产质量压力/正常/重估；行业PE区间提供盈利交叉检查；这些是House区间，不是外部目标价。
为什么用这组概率：30\%/50\%/20\%是统一筛选先验：基准情景权重最高，同时保留熊市和牛市的实质性下行/上行贡献；它是House情景权重，不是经过统计估计的客观事件概率。
证据如何参与：官方半年报预告、Q1财务包和行情/K线提供事实输入；研报/来源用于分母或可比背景。没有正权重外部目标锚，House价值不伪装为Street共识。
\begin{align*}
\mathrm{EPS}_{H1}^d &= \frac{195.03}{176.2890} = 1.1063\ \mathrm{CNY/share}\\
V_s &= 0.65\left(\mathrm{BPS}_1\times \mathrm{PB}_s\right)+0.35\left(\mathrm{EPS}_s\times \mathrm{PE}_s\right)\\
V_{b/m/b} &= 14.29/21.16/28.45\ \mathrm{CNY/share}\quad(\mathrm{PB}=13.89/19.44/25.00,\ \mathrm{PE}=15.05/24.34/34.85)\\
V_{prob} &= 0.30\times 13.93+0.50\times 21.16+0.20\times 28.45=20.45\\
P_{target} &= V_{prob}=20.45\quad(\text{外部锚权重}=0)\\
\mathrm{Upside} &= \frac{20.45}{18.58}-1=10.1\%\\
\end{align*}

\paragraph{600884 杉杉股份：成长盈利情景 PE}
分母：H1 deducted EPS bridged with H2/H1 conversion；证据质量：中；审计状态：PASS。
为什么选这个方法：电力设备具有成长或盈利持续期特征，因此模型给盈利桥接配置较高PE区间，而不是直接给主题收入估值；价格和现金流检查防止把泛化成长叙事直接变成EPS。
为什么选这个分母：使用H1扣非利润以减少非经常性收益污染；股本由总市值/当前价反推；H2/H1转换比例是明确的季节性和交付假设，不是机械年化。
为什么用这组倍数：PE区间是House对成长持续期的情景判断；牛市情景要求业绩桥接和经营证据持续，不是市场一致目标价。
为什么用这组概率：30\%/50\%/20\%是统一筛选先验：基准情景权重最高，同时保留熊市和牛市的实质性下行/上行贡献；它是House情景权重，不是经过统计估计的客观事件概率。
证据如何参与：官方半年报预告、Q1财务包和行情/K线提供事实输入；研报/来源用于分母或可比背景。没有正权重外部目标锚，House价值不伪装为Street共识。
\begin{align*}
\mathrm{Shares} &= \frac{273.08}{12.14}=22.4942\ \mathrm{100m\ shares}\\
\mathrm{EPS}_{H1}^d &= \frac{7.75}{22.4942}=0.3445\\
\mathrm{EPS}_s &= \mathrm{EPS}_{H1}^d(1+r_s)=0.59/0.69/0.79\\
V_s &= \mathrm{EPS}_s\times \mathrm{PE}_s=9.11/15.16/22.19\\
V_{prob} &= 0.30\times 9.11+0.50\times 15.16+0.20\times 22.19=14.75\\
P_{target} &= V_{prob}=14.75\quad(\text{外部锚权重}=0)\\
\mathrm{Upside} &= \frac{14.75}{12.14}-1=21.5\%\\
\end{align*}

\paragraph{002842 翔鹭钨业：周期调整扣非 EPS / PE}
分母：H1 deducted EPS bridged with H2/H1 conversion；证据质量：中；审计状态：PASS。
为什么选这个方法：有色金属盈利受商品价格、运价、价差或利用率周期影响，因此使用周期调整EPS/PE区间，而不是把某一个季度的PE当作全年估值。H2转换比例用于检验当前周期能否延续。
为什么选这个分母：使用H1扣非利润以减少非经常性收益污染；股本由总市值/当前价反推；H2/H1转换比例是明确的季节性和交付假设，不是机械年化。
为什么用这组倍数：熊/基准/牛PE区间是针对该周期业务类型的House假设，不是外部券商目标价；熊值较低反映周期反转，牛值较高必须由价格、销量和现金流共同确认。
为什么用这组概率：30\%/50\%/20\%是统一筛选先验：基准情景权重最高，同时保留熊市和牛市的实质性下行/上行贡献；它是House情景权重，不是经过统计估计的客观事件概率。
证据如何参与：官方半年报预告、Q1财务包和行情/K线提供事实输入；研报/来源用于分母或可比背景。没有正权重外部目标锚，House价值不伪装为Street共识。
\begin{align*}
\mathrm{Shares} &= \frac{117.65}{35.96}=3.2717\ \mathrm{100m\ shares}\\
\mathrm{EPS}_{H1}^d &= \frac{5.35}{3.2717}=1.6352\\
\mathrm{EPS}_s &= \mathrm{EPS}_{H1}^d(1+r_s)=2.53/3.03/3.52\\
V_s &= \mathrm{EPS}_s\times \mathrm{PE}_s=20.28/36.30/56.25\\
V_{prob} &= 0.30\times 20.28+0.50\times 36.30+0.20\times 56.25=35.48\\
P_{target} &= V_{prob}=35.48\quad(\text{外部锚权重}=0)\\
\mathrm{Upside} &= \frac{35.48}{35.96}-1=-1.3\%\\
\end{align*}

\paragraph{002407 多氟多：周期调整扣非 EPS / PE}
分母：H1 deducted EPS bridged with H2/H1 conversion；证据质量：中；审计状态：PASS。
为什么选这个方法：基础化工盈利受商品价格、运价、价差或利用率周期影响，因此使用周期调整EPS/PE区间，而不是把某一个季度的PE当作全年估值。H2转换比例用于检验当前周期能否延续。
为什么选这个分母：使用H1扣非利润以减少非经常性收益污染；股本由总市值/当前价反推；H2/H1转换比例是明确的季节性和交付假设，不是机械年化。
为什么用这组倍数：熊/基准/牛PE区间是针对该周期业务类型的House假设，不是外部券商目标价；熊值较低反映周期反转，牛值较高必须由价格、销量和现金流共同确认。
为什么用这组概率：30\%/50\%/20\%是统一筛选先验：基准情景权重最高，同时保留熊市和牛市的实质性下行/上行贡献；它是House情景权重，不是经过统计估计的客观事件概率。
证据如何参与：官方半年报预告、Q1财务包和行情/K线提供事实输入；研报/来源用于分母或可比背景。没有正权重外部目标锚，House价值不伪装为Street共识。
\begin{align*}
\mathrm{Shares} &= \frac{403.32}{33.88}=11.9044\ \mathrm{100m\ shares}\\
\mathrm{EPS}_{H1}^d &= \frac{4.89}{11.9044}=0.4106\\
\mathrm{EPS}_s &= \mathrm{EPS}_{H1}^d(1+r_s)=0.64/0.76/0.88\\
V_s &= \mathrm{EPS}_s\times \mathrm{PE}_s=6.37/10.64/15.89\\
V_{prob} &= 0.30\times 6.37+0.50\times 10.64+0.20\times 15.89=10.41\\
P_{target} &= V_{prob}=10.41\quad(\text{外部锚权重}=0)\\
\mathrm{Upside} &= \frac{10.41}{33.88}-1=-69.3\%\\
\end{align*}

\paragraph{002653 海思科：成长盈利情景 PE}
分母：H1 deducted EPS bridged with H2/H1 conversion；证据质量：中高；审计状态：PASS。
为什么选这个方法：医药生物具有成长或盈利持续期特征，因此模型给盈利桥接配置较高PE区间，而不是直接给主题收入估值；价格和现金流检查防止把泛化成长叙事直接变成EPS。
为什么选这个分母：使用H1扣非利润以减少非经常性收益污染；股本由总市值/当前价反推；H2/H1转换比例是明确的季节性和交付假设，不是机械年化。
为什么用这组倍数：PE区间是House对成长持续期的情景判断；牛市情景要求业绩桥接和经营证据持续，不是市场一致目标价。
为什么用这组概率：30\%/50\%/20\%是统一筛选先验：基准情景权重最高，同时保留熊市和牛市的实质性下行/上行贡献；它是House情景权重，不是经过统计估计的客观事件概率。
证据如何参与：官方半年报预告、Q1财务包和行情/K线提供事实输入；研报/来源用于分母或可比背景。没有正权重外部目标锚，House价值不伪装为Street共识。
\begin{align*}
\mathrm{Shares} &= \frac{830.20}{74.13}=11.1992\ \mathrm{100m\ shares}\\
\mathrm{EPS}_{H1}^d &= \frac{7.30}{11.1992}=0.6518\\
\mathrm{EPS}_s &= \mathrm{EPS}_{H1}^d(1+r_s)=1.11/1.30/1.50\\
V_s &= \mathrm{EPS}_s\times \mathrm{PE}_s=19.95/32.59/47.97\\
V_{prob} &= 0.30\times 19.95+0.50\times 32.59+0.20\times 47.97=31.87\\
P_{target} &= V_{prob}=31.87\quad(\text{外部锚权重}=0)\\
\mathrm{Upside} &= \frac{31.87}{74.13}-1=-57.0\%\\
\end{align*}

\paragraph{603629 利通电子：成长盈利情景 PE}
分母：H1 deducted EPS bridged with H2/H1 conversion；证据质量：中；审计状态：PASS。
为什么选这个方法：电子具有成长或盈利持续期特征，因此模型给盈利桥接配置较高PE区间，而不是直接给主题收入估值；价格和现金流检查防止把泛化成长叙事直接变成EPS。
为什么选这个分母：使用H1扣非利润以减少非经常性收益污染；股本由总市值/当前价反推；H2/H1转换比例是明确的季节性和交付假设，不是机械年化。
为什么用这组倍数：PE区间是House对成长持续期的情景判断；牛市情景要求业绩桥接和经营证据持续，不是市场一致目标价。
为什么用这组概率：30\%/50\%/20\%是统一筛选先验：基准情景权重最高，同时保留熊市和牛市的实质性下行/上行贡献；它是House情景权重，不是经过统计估计的客观事件概率。
证据如何参与：官方半年报预告、Q1财务包和行情/K线提供事实输入；研报/来源用于分母或可比背景。没有正权重外部目标锚，House价值不伪装为Street共识。
\begin{align*}
\mathrm{Shares} &= \frac{411.07}{112.09}=3.6673\ \mathrm{100m\ shares}\\
\mathrm{EPS}_{H1}^d &= \frac{6.12}{3.6673}=1.6674\\
\mathrm{EPS}_s &= \mathrm{EPS}_{H1}^d(1+r_s)=2.83/3.33/3.84\\
V_s &= \mathrm{EPS}_s\times \mathrm{PE}_s=51.02/86.71/130.39\\
V_{prob} &= 0.30\times 51.02+0.50\times 86.71+0.20\times 130.39=84.74\\
P_{target} &= V_{prob}=84.74\quad(\text{外部锚权重}=0)\\
\mathrm{Upside} &= \frac{84.74}{112.09}-1=-24.4\%\\
\end{align*}

\paragraph{600961 株冶集团：周期调整扣非 EPS / PE}
分母：H1 deducted EPS bridged with H2/H1 conversion；证据质量：中高；审计状态：PASS。
为什么选这个方法：有色金属盈利受商品价格、运价、价差或利用率周期影响，因此使用周期调整EPS/PE区间，而不是把某一个季度的PE当作全年估值。H2转换比例用于检验当前周期能否延续。
为什么选这个分母：使用H1扣非利润以减少非经常性收益污染；股本由总市值/当前价反推；H2/H1转换比例是明确的季节性和交付假设，不是机械年化。
为什么用这组倍数：熊/基准/牛PE区间是针对该周期业务类型的House假设，不是外部券商目标价；熊值较低反映周期反转，牛值较高必须由价格、销量和现金流共同确认。
为什么用这组概率：30\%/50\%/20\%是统一筛选先验：基准情景权重最高，同时保留熊市和牛市的实质性下行/上行贡献；它是House情景权重，不是经过统计估计的客观事件概率。
证据如何参与：官方半年报预告、Q1财务包和行情/K线提供事实输入；研报/来源用于分母或可比背景。没有正权重外部目标锚，House价值不伪装为Street共识。
\begin{align*}
\mathrm{Shares} &= \frac{265.64}{24.76}=10.7286\ \mathrm{100m\ shares}\\
\mathrm{EPS}_{H1}^d &= \frac{19.10}{10.7286}=1.7803\\
\mathrm{EPS}_s &= \mathrm{EPS}_{H1}^d(1+r_s)=2.76/3.29/3.83\\
V_s &= \mathrm{EPS}_s\times \mathrm{PE}_s=18.57/39.52/61.24\\
V_{prob} &= 0.30\times 18.57+0.50\times 39.52+0.20\times 61.24=37.58\\
P_{target} &= V_{prob}=37.58\quad(\text{外部锚权重}=0)\\
\mathrm{Upside} &= \frac{37.58}{24.76}-1=51.8\%\\
\end{align*}

\paragraph{002128 电投能源：周期调整扣非 EPS / PE}
分母：H1 deducted EPS bridged with H2/H1 conversion；证据质量：中高；审计状态：PASS。
为什么选这个方法：煤炭盈利受商品价格、运价、价差或利用率周期影响，因此使用周期调整EPS/PE区间，而不是把某一个季度的PE当作全年估值。H2转换比例用于检验当前周期能否延续。
为什么选这个分母：使用H1扣非利润以减少非经常性收益污染；股本由总市值/当前价反推；H2/H1转换比例是明确的季节性和交付假设，不是机械年化。
为什么用这组倍数：熊/基准/牛PE区间是针对该周期业务类型的House假设，不是外部券商目标价；熊值较低反映周期反转，牛值较高必须由价格、销量和现金流共同确认。
为什么用这组概率：30\%/50\%/20\%是统一筛选先验：基准情景权重最高，同时保留熊市和牛市的实质性下行/上行贡献；它是House情景权重，不是经过统计估计的客观事件概率。
证据如何参与：官方半年报预告、Q1财务包和行情/K线提供事实输入；研报/来源用于分母或可比背景。没有正权重外部目标锚，House价值不伪装为Street共识。
\begin{align*}
\mathrm{Shares} &= \frac{782.95}{26.51}=29.5341\ \mathrm{100m\ shares}\\
\mathrm{EPS}_{H1}^d &= \frac{40.47}{29.5341}=1.3701\\
\mathrm{EPS}_s &= \mathrm{EPS}_{H1}^d(1+r_s)=2.12/2.83/3.08\\
V_s &= \mathrm{EPS}_s\times \mathrm{PE}_s=14.87/25.47/33.91\\
V_{prob} &= 0.30\times 14.87+0.50\times 25.47+0.20\times 33.91=23.98\\
P_{target} &= V_{prob}=23.98\quad(\text{外部锚权重}=0)\\
\mathrm{Upside} &= \frac{23.98}{26.51}-1=-9.5\%\\
\end{align*}

\paragraph{601168 西部矿业：周期调整扣非 EPS / PE}
分母：H1 deducted EPS bridged with H2/H1 conversion；证据质量：高；审计状态：PASS。
为什么选这个方法：有色金属盈利受商品价格、运价、价差或利用率周期影响，因此使用周期调整EPS/PE区间，而不是把某一个季度的PE当作全年估值。H2转换比例用于检验当前周期能否延续。
为什么选这个分母：使用H1扣非利润以减少非经常性收益污染；股本由总市值/当前价反推；H2/H1转换比例是明确的季节性和交付假设，不是机械年化。
为什么用这组倍数：熊/基准/牛PE区间是针对该周期业务类型的House假设，不是外部券商目标价；熊值较低反映周期反转，牛值较高必须由价格、销量和现金流共同确认。
为什么用这组概率：30\%/50\%/20\%是统一筛选先验：基准情景权重最高，同时保留熊市和牛市的实质性下行/上行贡献；它是House情景权重，不是经过统计估计的客观事件概率。
证据如何参与：官方半年报预告、Q1财务包和行情/K线提供事实输入；研报/来源用于分母或可比背景。外部目标价38.50元，赋予正权重10\%。
\begin{align*}
\mathrm{Shares} &= \frac{790.68}{33.18}=23.8300\ \mathrm{100m\ shares}\\
\mathrm{EPS}_{H1}^d &= \frac{41.90}{23.8300}=1.7583\\
\mathrm{EPS}_s &= \mathrm{EPS}_{H1}^d(1+r_s)=2.73/3.25/3.78\\
V_s &= \mathrm{EPS}_s\times \mathrm{PE}_s=21.80/39.03/60.49\\
V_{prob} &= 0.30\times 21.80+0.50\times 39.03+0.20\times 60.49=38.15\\
P_{target} &= 38.15\times 0.90+38.50\times 0.10=38.19\\
\mathrm{Upside} &= \frac{38.19}{33.18}-1=15.1\%\\
\end{align*}

\paragraph{600188 兖矿能源：周期调整扣非 EPS / PE}
分母：H1 deducted EPS bridged with H2/H1 conversion；证据质量：中高；审计状态：PASS。
为什么选这个方法：煤炭盈利受商品价格、运价、价差或利用率周期影响，因此使用周期调整EPS/PE区间，而不是把某一个季度的PE当作全年估值。H2转换比例用于检验当前周期能否延续。
为什么选这个分母：使用H1扣非利润以减少非经常性收益污染；股本由总市值/当前价反推；H2/H1转换比例是明确的季节性和交付假设，不是机械年化。
为什么用这组倍数：熊/基准/牛PE区间是针对该周期业务类型的House假设，不是外部券商目标价；熊值较低反映周期反转，牛值较高必须由价格、销量和现金流共同确认。
为什么用这组概率：30\%/50\%/20\%是统一筛选先验：基准情景权重最高，同时保留熊市和牛市的实质性下行/上行贡献；它是House情景权重，不是经过统计估计的客观事件概率。
证据如何参与：官方半年报预告、Q1财务包和行情/K线提供事实输入；研报/来源用于分母或可比背景。没有正权重外部目标锚，House价值不伪装为Street共识。
\begin{align*}
\mathrm{Shares} &= \frac{2030.46}{20.23}=100.3688\ \mathrm{100m\ shares}\\
\mathrm{EPS}_{H1}^d &= \frac{45.00}{100.3688}=0.4483\\
\mathrm{EPS}_s &= \mathrm{EPS}_{H1}^d(1+r_s)=0.69/1.64/1.01\\
V_s &= \mathrm{EPS}_s\times \mathrm{PE}_s=4.86/11.10/14.76\\
V_{prob} &= 0.30\times 4.86+0.50\times 11.10+0.20\times 14.76=9.96\\
P_{target} &= V_{prob}=9.96\quad(\text{外部锚权重}=0)\\
\mathrm{Upside} &= \frac{9.96}{20.23}-1=-50.8\%\\
\end{align*}

\paragraph{603162 海通发展：周期调整扣非 EPS / PE}
分母：H1 deducted EPS bridged with H2/H1 conversion；证据质量：中高；审计状态：PASS。
为什么选这个方法：交通运输盈利受商品价格、运价、价差或利用率周期影响，因此使用周期调整EPS/PE区间，而不是把某一个季度的PE当作全年估值。H2转换比例用于检验当前周期能否延续。
为什么选这个分母：使用H1扣非利润以减少非经常性收益污染；股本由总市值/当前价反推；H2/H1转换比例是明确的季节性和交付假设，不是机械年化。
为什么用这组倍数：熊/基准/牛PE区间是针对该周期业务类型的House假设，不是外部券商目标价；熊值较低反映周期反转，牛值较高必须由价格、销量和现金流共同确认。
为什么用这组概率：30\%/50\%/20\%是统一筛选先验：基准情景权重最高，同时保留熊市和牛市的实质性下行/上行贡献；它是House情景权重，不是经过统计估计的客观事件概率。
证据如何参与：官方半年报预告、Q1财务包和行情/K线提供事实输入；研报/来源用于分母或可比背景。没有正权重外部目标锚，House价值不伪装为Street共识。
\begin{align*}
\mathrm{Shares} &= \frac{154.62}{11.24}=13.7562\ \mathrm{100m\ shares}\\
\mathrm{EPS}_{H1}^d &= \frac{5.23}{13.7562}=0.3802\\
\mathrm{EPS}_s &= \mathrm{EPS}_{H1}^d(1+r_s)=0.59/0.75/0.82\\
V_s &= \mathrm{EPS}_s\times \mathrm{PE}_s=4.71/9.00/13.08\\
V_{prob} &= 0.30\times 4.71+0.50\times 9.00+0.20\times 13.08=8.53\\
P_{target} &= V_{prob}=8.53\quad(\text{外部锚权重}=0)\\
\mathrm{Upside} &= \frac{8.53}{11.24}-1=-24.1\%\\
\end{align*}

\paragraph{002240 盛新锂能：周期调整 EPS / PE}
分母：2026E cycle-adjusted EPS with lithium-price and cash-flow bridge；证据质量：中高；审计状态：PASS。
为什么选这个方法：最新H1利润存在低基数、扭亏或周期扭曲，因此先正常化EPS，再套业务匹配PE区间；业务匹配可比包括赣锋锂业、天齐锂业、天华新能。
为什么选这个分母：正常化分母为：2026E cycle-adjusted EPS with lithium-price and cash-flow bridge；H1机械年化指标只作为市场隐含压力测试，不作为最终目标价分母。
为什么用这组倍数：熊/基准/牛倍数分别表达下行、正常兑现和上行重估；区间通过可比方法和当前价格隐含指标交叉检查（锂价和自有矿贡献改善，但H1经营现金流为负、周期持续性仍是核心变量），因此只能作为条件性House价值。
为什么用这组概率：30\%/50\%/20\%是统一筛选先验：基准情景权重最高，同时保留熊市和牛市的实质性下行/上行贡献；它是House情景权重，不是经过统计估计的客观事件概率。
证据如何参与：官方半年报预告、Q1财务包和行情/K线提供事实输入；研报/来源用于分母或可比背景。没有正权重外部目标锚，House价值不伪装为Street共识。
\begin{align*}
V_s=\mathrm{EPS}_{normalized,s}\times \mathrm{PE}_s\\
V_{b/m/b} &= 21.60/39.20/59.04\ \mathrm{CNY/share}\\
V_{prob} &= 0.30\times 21.60+0.50\times 39.20+0.20\times 59.04=37.89\\
P_{target} &= V_{prob}=37.89\quad(\text{外部锚权重}=0)\\
\mathrm{Upside} &= \frac{37.89}{31.11}-1=21.8\%\\
\end{align*}

\paragraph{002440 闰土股份：周期调整扣非 EPS / PE}
分母：H1 deducted EPS bridged with H2/H1 conversion；证据质量：中；审计状态：PASS。
为什么选这个方法：基础化工盈利受商品价格、运价、价差或利用率周期影响，因此使用周期调整EPS/PE区间，而不是把某一个季度的PE当作全年估值。H2转换比例用于检验当前周期能否延续。
为什么选这个分母：使用H1扣非利润以减少非经常性收益污染；股本由总市值/当前价反推；H2/H1转换比例是明确的季节性和交付假设，不是机械年化。
为什么用这组倍数：熊/基准/牛PE区间是针对该周期业务类型的House假设，不是外部券商目标价；熊值较低反映周期反转，牛值较高必须由价格、销量和现金流共同确认。
为什么用这组概率：30\%/50\%/20\%是统一筛选先验：基准情景权重最高，同时保留熊市和牛市的实质性下行/上行贡献；它是House情景权重，不是经过统计估计的客观事件概率。
证据如何参与：官方半年报预告、Q1财务包和行情/K线提供事实输入；研报/来源用于分母或可比背景。没有正权重外部目标锚，House价值不伪装为Street共识。
\begin{align*}
\mathrm{Shares} &= \frac{126.90}{11.29}=11.2400\ \mathrm{100m\ shares}\\
\mathrm{EPS}_{H1}^d &= \frac{4.40}{11.2400}=0.3915\\
\mathrm{EPS}_s &= \mathrm{EPS}_{H1}^d(1+r_s)=0.61/0.72/0.84\\
V_s &= \mathrm{EPS}_s\times \mathrm{PE}_s=6.07/10.14/15.15\\
V_{prob} &= 0.30\times 6.07+0.50\times 10.14+0.20\times 15.15=9.92\\
P_{target} &= V_{prob}=9.92\quad(\text{外部锚权重}=0)\\
\mathrm{Upside} &= \frac{9.92}{11.29}-1=-12.1\%\\
\end{align*}

\paragraph{002532 天山铝业：周期调整扣非 EPS / PE}
分母：H1 deducted EPS bridged with H2/H1 conversion；证据质量：中高；审计状态：PASS。
为什么选这个方法：有色金属盈利受商品价格、运价、价差或利用率周期影响，因此使用周期调整EPS/PE区间，而不是把某一个季度的PE当作全年估值。H2转换比例用于检验当前周期能否延续。
为什么选这个分母：使用H1扣非利润以减少非经常性收益污染；股本由总市值/当前价反推；H2/H1转换比例是明确的季节性和交付假设，不是机械年化。
为什么用这组倍数：熊/基准/牛PE区间是针对该周期业务类型的House假设，不是外部券商目标价；熊值较低反映周期反转，牛值较高必须由价格、销量和现金流共同确认。
为什么用这组概率：30\%/50\%/20\%是统一筛选先验：基准情景权重最高，同时保留熊市和牛市的实质性下行/上行贡献；它是House情景权重，不是经过统计估计的客观事件概率。
证据如何参与：官方半年报预告、Q1财务包和行情/K线提供事实输入；研报/来源用于分母或可比背景。没有正权重外部目标锚，House价值不伪装为Street共识。
\begin{align*}
\mathrm{Shares} &= \frac{590.16}{12.75}=46.2871\ \mathrm{100m\ shares}\\
\mathrm{EPS}_{H1}^d &= \frac{40.90}{46.2871}=0.8836\\
\mathrm{EPS}_s &= \mathrm{EPS}_{H1}^d(1+r_s)=1.37/1.92/1.90\\
V_s &= \mathrm{EPS}_s\times \mathrm{PE}_s=9.56/23.03/30.40\\
V_{prob} &= 0.30\times 9.56+0.50\times 23.03+0.20\times 30.40=20.46\\
P_{target} &= V_{prob}=20.46\quad(\text{外部锚权重}=0)\\
\mathrm{Upside} &= \frac{20.46}{12.75}-1=60.5\%\\
\end{align*}

\paragraph{301377 鼎泰高科：常规扣非 EPS / PE}
分母：H1 deducted EPS bridged with H2/H1 conversion；证据质量：中高；审计状态：PASS。
为什么选这个方法：该标的H1扣非利润为正且未触发恢复估值阻断，因此扣非EPS是最透明的筛选分母；情景桥接避免把一个H1数字直接当作全年点预测。
为什么选这个分母：使用H1扣非利润以减少非经常性收益污染；股本由总市值/当前价反推；H2/H1转换比例是明确的季节性和交付假设，不是机械年化。
为什么用这组倍数：PE区间按公司所属行业匹配，表达估值压缩/正常/重估三种情景；这是House筛选区间，正式升级前仍需当前可审计的Street区间。
为什么用这组概率：30\%/50\%/20\%是统一筛选先验：基准情景权重最高，同时保留熊市和牛市的实质性下行/上行贡献；它是House情景权重，不是经过统计估计的客观事件概率。
证据如何参与：官方半年报预告、Q1财务包和行情/K线提供事实输入；研报/来源用于分母或可比背景。没有正权重外部目标锚，House价值不伪装为Street共识。
\begin{align*}
\mathrm{Shares} &= \frac{1984.02}{467.88}=4.2404\ \mathrm{100m\ shares}\\
\mathrm{EPS}_{H1}^d &= \frac{6.60}{4.2404}=1.5564\\
\mathrm{EPS}_s &= \mathrm{EPS}_{H1}^d(1+r_s)=2.57/2.96/3.42\\
V_s &= \mathrm{EPS}_s\times \mathrm{PE}_s=35.95/59.14/95.88\\
V_{prob} &= 0.30\times 35.95+0.50\times 59.14+0.20\times 95.88=59.53\\
P_{target} &= V_{prob}=59.53\quad(\text{外部锚权重}=0)\\
\mathrm{Upside} &= \frac{59.53}{467.88}-1=-87.3\%\\
\end{align*}

\paragraph{688257 新锐股份：常规扣非 EPS / PE}
分母：H1 deducted EPS bridged with H2/H1 conversion；证据质量：中高；审计状态：PASS。
为什么选这个方法：该标的H1扣非利润为正且未触发恢复估值阻断，因此扣非EPS是最透明的筛选分母；情景桥接避免把一个H1数字直接当作全年点预测。
为什么选这个分母：使用H1扣非利润以减少非经常性收益污染；股本由总市值/当前价反推；H2/H1转换比例是明确的季节性和交付假设，不是机械年化。
为什么用这组倍数：PE区间按公司所属行业匹配，表达估值压缩/正常/重估三种情景；这是House筛选区间，正式升级前仍需当前可审计的Street区间。
为什么用这组概率：30\%/50\%/20\%是统一筛选先验：基准情景权重最高，同时保留熊市和牛市的实质性下行/上行贡献；它是House情景权重，不是经过统计估计的客观事件概率。
证据如何参与：官方半年报预告、Q1财务包和行情/K线提供事实输入；研报/来源用于分母或可比背景。没有正权重外部目标锚，House价值不伪装为Street共识。
\begin{align*}
\mathrm{Shares} &= \frac{284.42}{80.06}=3.5526\ \mathrm{100m\ shares}\\
\mathrm{EPS}_{H1}^d &= \frac{5.90}{3.5526}=1.6608\\
\mathrm{EPS}_s &= \mathrm{EPS}_{H1}^d(1+r_s)=2.74/3.16/3.65\\
V_s &= \mathrm{EPS}_s\times \mathrm{PE}_s=38.36/63.11/102.30\\
V_{prob} &= 0.30\times 38.36+0.50\times 63.11+0.20\times 102.30=63.52\\
P_{target} &= V_{prob}=63.52\quad(\text{外部锚权重}=0)\\
\mathrm{Upside} &= \frac{63.52}{80.06}-1=-20.7\%\\
\end{align*}

\paragraph{000408 藏格矿业：周期调整扣非 EPS / PE}
分母：H1 deducted EPS bridged with H2/H1 conversion；证据质量：中高；审计状态：PASS。
为什么选这个方法：基础化工盈利受商品价格、运价、价差或利用率周期影响，因此使用周期调整EPS/PE区间，而不是把某一个季度的PE当作全年估值。H2转换比例用于检验当前周期能否延续。
为什么选这个分母：使用H1扣非利润以减少非经常性收益污染；股本由总市值/当前价反推；H2/H1转换比例是明确的季节性和交付假设，不是机械年化。
为什么用这组倍数：熊/基准/牛PE区间是针对该周期业务类型的House假设，不是外部券商目标价；熊值较低反映周期反转，牛值较高必须由价格、销量和现金流共同确认。
为什么用这组概率：30\%/50\%/20\%是统一筛选先验：基准情景权重最高，同时保留熊市和牛市的实质性下行/上行贡献；它是House情景权重，不是经过统计估计的客观事件概率。
证据如何参与：官方半年报预告、Q1财务包和行情/K线提供事实输入；研报/来源用于分母或可比背景。没有正权重外部目标锚，House价值不伪装为Street共识。
\begin{align*}
\mathrm{Shares} &= \frac{1121.77}{71.50}=15.6891\ \mathrm{100m\ shares}\\
\mathrm{EPS}_{H1}^d &= \frac{37.13}{15.6891}=2.3666\\
\mathrm{EPS}_s &= \mathrm{EPS}_{H1}^d(1+r_s)=3.67/5.12/5.09\\
V_s &= \mathrm{EPS}_s\times \mathrm{PE}_s=36.68/71.68/91.59\\
V_{prob} &= 0.30\times 36.68+0.50\times 71.68+0.20\times 91.59=65.16\\
P_{target} &= V_{prob}=65.16\quad(\text{外部锚权重}=0)\\
\mathrm{Upside} &= \frac{65.16}{71.50}-1=-8.9\%\\
\end{align*}

\paragraph{600601 方正科技：成长盈利情景 PE}
分母：H1 deducted EPS bridged with H2/H1 conversion；证据质量：中；审计状态：PASS。
为什么选这个方法：电子具有成长或盈利持续期特征，因此模型给盈利桥接配置较高PE区间，而不是直接给主题收入估值；价格和现金流检查防止把泛化成长叙事直接变成EPS。
为什么选这个分母：使用H1扣非利润以减少非经常性收益污染；股本由总市值/当前价反推；H2/H1转换比例是明确的季节性和交付假设，不是机械年化。
为什么用这组倍数：PE区间是House对成长持续期的情景判断；牛市情景要求业绩桥接和经营证据持续，不是市场一致目标价。
为什么用这组概率：30\%/50\%/20\%是统一筛选先验：基准情景权重最高，同时保留熊市和牛市的实质性下行/上行贡献；它是House情景权重，不是经过统计估计的客观事件概率。
证据如何参与：官方半年报预告、Q1财务包和行情/K线提供事实输入；研报/来源用于分母或可比背景。没有正权重外部目标锚，House价值不伪装为Street共识。
\begin{align*}
\mathrm{Shares} &= \frac{500.46}{11.71}=42.7378\ \mathrm{100m\ shares}\\
\mathrm{EPS}_{H1}^d &= \frac{5.55}{42.7378}=0.1299\\
\mathrm{EPS}_s &= \mathrm{EPS}_{H1}^d(1+r_s)=0.22/0.26/0.30\\
V_s &= \mathrm{EPS}_s\times \mathrm{PE}_s=3.97/6.75/10.16\\
V_{prob} &= 0.30\times 3.97+0.50\times 6.75+0.20\times 10.16=6.60\\
P_{target} &= V_{prob}=6.60\quad(\text{外部锚权重}=0)\\
\mathrm{Upside} &= \frac{6.60}{11.71}-1=-43.6\%\\
\end{align*}

\paragraph{002475 立讯精密：成长盈利情景 PE}
分母：H1 deducted EPS bridged with H2/H1 conversion；证据质量：中；审计状态：PASS。
为什么选这个方法：电子具有成长或盈利持续期特征，因此模型给盈利桥接配置较高PE区间，而不是直接给主题收入估值；价格和现金流检查防止把泛化成长叙事直接变成EPS。
为什么选这个分母：使用H1扣非利润以减少非经常性收益污染；股本由总市值/当前价反推；H2/H1转换比例是明确的季节性和交付假设，不是机械年化。
为什么用这组倍数：PE区间是House对成长持续期的情景判断；牛市情景要求业绩桥接和经营证据持续，不是市场一致目标价。
为什么用这组概率：30\%/50\%/20\%是统一筛选先验：基准情景权重最高，同时保留熊市和牛市的实质性下行/上行贡献；它是House情景权重，不是经过统计估计的客观事件概率。
证据如何参与：官方半年报预告、Q1财务包和行情/K线提供事实输入；研报/来源用于分母或可比背景。没有正权重外部目标锚，House价值不伪装为Street共识。
\begin{align*}
\mathrm{Shares} &= \frac{4671.48}{60.47}=77.2529\ \mathrm{100m\ shares}\\
\mathrm{EPS}_{H1}^d &= \frac{64.90}{77.2529}=0.8400\\
\mathrm{EPS}_s &= \mathrm{EPS}_{H1}^d(1+r_s)=1.43/1.68/1.93\\
V_s &= \mathrm{EPS}_s\times \mathrm{PE}_s=25.71/43.68/65.69\\
V_{prob} &= 0.30\times 25.71+0.50\times 43.68+0.20\times 65.69=42.69\\
P_{target} &= V_{prob}=42.69\quad(\text{外部锚权重}=0)\\
\mathrm{Upside} &= \frac{42.69}{60.47}-1=-29.4\%\\
\end{align*}

\paragraph{600595 中孚实业：周期调整扣非 EPS / PE}
分母：H1 deducted EPS bridged with H2/H1 conversion；证据质量：中高；审计状态：PASS。
为什么选这个方法：有色金属盈利受商品价格、运价、价差或利用率周期影响，因此使用周期调整EPS/PE区间，而不是把某一个季度的PE当作全年估值。H2转换比例用于检验当前周期能否延续。
为什么选这个分母：使用H1扣非利润以减少非经常性收益污染；股本由总市值/当前价反推；H2/H1转换比例是明确的季节性和交付假设，不是机械年化。
为什么用这组倍数：熊/基准/牛PE区间是针对该周期业务类型的House假设，不是外部券商目标价；熊值较低反映周期反转，牛值较高必须由价格、销量和现金流共同确认。
为什么用这组概率：30\%/50\%/20\%是统一筛选先验：基准情景权重最高，同时保留熊市和牛市的实质性下行/上行贡献；它是House情景权重，不是经过统计估计的客观事件概率。
证据如何参与：官方半年报预告、Q1财务包和行情/K线提供事实输入；研报/来源用于分母或可比背景。没有正权重外部目标锚，House价值不伪装为Street共识。
\begin{align*}
\mathrm{Shares} &= \frac{258.91}{6.46}=40.0789\ \mathrm{100m\ shares}\\
\mathrm{EPS}_{H1}^d &= \frac{19.25}{40.0789}=0.4803\\
\mathrm{EPS}_s &= \mathrm{EPS}_{H1}^d(1+r_s)=0.74/1.00/1.03\\
V_s &= \mathrm{EPS}_s\times \mathrm{PE}_s=4.84/12.00/16.52\\
V_{prob} &= 0.30\times 4.84+0.50\times 12.00+0.20\times 16.52=10.76\\
P_{target} &= V_{prob}=10.76\quad(\text{外部锚权重}=0)\\
\mathrm{Upside} &= \frac{10.76}{6.46}-1=66.6\%\\
\end{align*}

\paragraph{600487 亨通光电：成长盈利情景 PE}
分母：H1 deducted EPS bridged with H2/H1 conversion；证据质量：中高；审计状态：PASS。
为什么选这个方法：通信具有成长或盈利持续期特征，因此模型给盈利桥接配置较高PE区间，而不是直接给主题收入估值；价格和现金流检查防止把泛化成长叙事直接变成EPS。
为什么选这个分母：使用H1扣非利润以减少非经常性收益污染；股本由总市值/当前价反推；H2/H1转换比例是明确的季节性和交付假设，不是机械年化。
为什么用这组倍数：PE区间是House对成长持续期的情景判断；牛市情景要求业绩桥接和经营证据持续，不是市场一致目标价。
为什么用这组概率：30\%/50\%/20\%是统一筛选先验：基准情景权重最高，同时保留熊市和牛市的实质性下行/上行贡献；它是House情景权重，不是经过统计估计的客观事件概率。
证据如何参与：官方半年报预告、Q1财务包和行情/K线提供事实输入；研报/来源用于分母或可比背景。没有正权重外部目标锚，House价值不伪装为Street共识。
\begin{align*}
\mathrm{Shares} &= \frac{1752.37}{71.05}=24.6639\ \mathrm{100m\ shares}\\
\mathrm{EPS}_{H1}^d &= \frac{33.26}{24.6639}=1.3485\\
\mathrm{EPS}_s &= \mathrm{EPS}_{H1}^d(1+r_s)=2.29/2.70/3.10\\
V_s &= \mathrm{EPS}_s\times \mathrm{PE}_s=41.26/70.12/105.45\\
V_{prob} &= 0.30\times 41.26+0.50\times 70.12+0.20\times 105.45=68.53\\
P_{target} &= V_{prob}=68.53\quad(\text{外部锚权重}=0)\\
\mathrm{Upside} &= \frac{68.53}{71.05}-1=-3.5\%\\
\end{align*}

\paragraph{600183 生益科技：成长盈利情景 PE}
分母：H1 deducted EPS bridged with H2/H1 conversion；证据质量：中高；审计状态：PASS。
为什么选这个方法：电子具有成长或盈利持续期特征，因此模型给盈利桥接配置较高PE区间，而不是直接给主题收入估值；价格和现金流检查防止把泛化成长叙事直接变成EPS。
为什么选这个分母：使用H1扣非利润以减少非经常性收益污染；股本由总市值/当前价反推；H2/H1转换比例是明确的季节性和交付假设，不是机械年化。
为什么用这组倍数：PE区间是House对成长持续期的情景判断；牛市情景要求业绩桥接和经营证据持续，不是市场一致目标价。
为什么用这组概率：30\%/50\%/20\%是统一筛选先验：基准情景权重最高，同时保留熊市和牛市的实质性下行/上行贡献；它是House情景权重，不是经过统计估计的客观事件概率。
证据如何参与：官方半年报预告、Q1财务包和行情/K线提供事实输入；研报/来源用于分母或可比背景。没有正权重外部目标锚，House价值不伪装为Street共识。
\begin{align*}
\mathrm{Shares} &= \frac{3702.53}{152.43}=24.2900\ \mathrm{100m\ shares}\\
\mathrm{EPS}_{H1}^d &= \frac{28.18}{24.2900}=1.1604\\
\mathrm{EPS}_s &= \mathrm{EPS}_{H1}^d(1+r_s)=1.97/2.32/2.67\\
V_s &= \mathrm{EPS}_s\times \mathrm{PE}_s=35.51/60.34/90.74\\
V_{prob} &= 0.30\times 35.51+0.50\times 60.34+0.20\times 90.74=58.97\\
P_{target} &= V_{prob}=58.97\quad(\text{外部锚权重}=0)\\
\mathrm{Upside} &= \frac{58.97}{152.43}-1=-61.3\%\\
\end{align*}

\paragraph{000630 铜陵有色：周期调整扣非 EPS / PE}
分母：H1 deducted EPS bridged with H2/H1 conversion；证据质量：中高；审计状态：PASS。
为什么选这个方法：有色金属盈利受商品价格、运价、价差或利用率周期影响，因此使用周期调整EPS/PE区间，而不是把某一个季度的PE当作全年估值。H2转换比例用于检验当前周期能否延续。
为什么选这个分母：使用H1扣非利润以减少非经常性收益污染；股本由总市值/当前价反推；H2/H1转换比例是明确的季节性和交付假设，不是机械年化。
为什么用这组倍数：熊/基准/牛PE区间是针对该周期业务类型的House假设，不是外部券商目标价；熊值较低反映周期反转，牛值较高必须由价格、销量和现金流共同确认。
为什么用这组概率：30\%/50\%/20\%是统一筛选先验：基准情景权重最高，同时保留熊市和牛市的实质性下行/上行贡献；它是House情景权重，不是经过统计估计的客观事件概率。
证据如何参与：官方半年报预告、Q1财务包和行情/K线提供事实输入；研报/来源用于分母或可比背景。没有正权重外部目标锚，House价值不伪装为Street共识。
\begin{align*}
\mathrm{Shares} &= \frac{812.61}{6.06}=134.0941\ \mathrm{100m\ shares}\\
\mathrm{EPS}_{H1}^d &= \frac{28.50}{134.0941}=0.2125\\
\mathrm{EPS}_s &= \mathrm{EPS}_{H1}^d(1+r_s)=0.33/0.56/0.46\\
V_s &= \mathrm{EPS}_s\times \mathrm{PE}_s=2.64/6.72/7.31\\
V_{prob} &= 0.30\times 2.64+0.50\times 6.72+0.20\times 7.31=5.61\\
P_{target} &= V_{prob}=5.61\quad(\text{外部锚权重}=0)\\
\mathrm{Upside} &= \frac{5.61}{6.06}-1=-7.4\%\\
\end{align*}

\paragraph{600549 厦门钨业：周期调整扣非 EPS / PE}
分母：H1 deducted EPS bridged with H2/H1 conversion；证据质量：中高；审计状态：PASS。
为什么选这个方法：有色金属盈利受商品价格、运价、价差或利用率周期影响，因此使用周期调整EPS/PE区间，而不是把某一个季度的PE当作全年估值。H2转换比例用于检验当前周期能否延续。
为什么选这个分母：使用H1扣非利润以减少非经常性收益污染；股本由总市值/当前价反推；H2/H1转换比例是明确的季节性和交付假设，不是机械年化。
为什么用这组倍数：熊/基准/牛PE区间是针对该周期业务类型的House假设，不是外部券商目标价；熊值较低反映周期反转，牛值较高必须由价格、销量和现金流共同确认。
为什么用这组概率：30\%/50\%/20\%是统一筛选先验：基准情景权重最高，同时保留熊市和牛市的实质性下行/上行贡献；它是House情景权重，不是经过统计估计的客观事件概率。
证据如何参与：官方半年报预告、Q1财务包和行情/K线提供事实输入；研报/来源用于分母或可比背景。没有正权重外部目标锚，House价值不伪装为Street共识。
\begin{align*}
\mathrm{Shares} &= \frac{922.07}{58.08}=15.8759\ \mathrm{100m\ shares}\\
\mathrm{EPS}_{H1}^d &= \frac{21.76}{15.8759}=1.3707\\
\mathrm{EPS}_s &= \mathrm{EPS}_{H1}^d(1+r_s)=2.12/2.93/2.95\\
V_s &= \mathrm{EPS}_s\times \mathrm{PE}_s=17.00/35.16/47.15\\
V_{prob} &= 0.30\times 17.00+0.50\times 35.16+0.20\times 47.15=32.11\\
P_{target} &= V_{prob}=32.11\quad(\text{外部锚权重}=0)\\
\mathrm{Upside} &= \frac{32.11}{58.08}-1=-44.7\%\\
\end{align*}

\paragraph{000100 TCL科技：成长盈利情景 PE}
分母：H1 deducted EPS bridged with H2/H1 conversion；证据质量：中高；审计状态：PASS。
为什么选这个方法：电子具有成长或盈利持续期特征，因此模型给盈利桥接配置较高PE区间，而不是直接给主题收入估值；价格和现金流检查防止把泛化成长叙事直接变成EPS。
为什么选这个分母：使用H1扣非利润以减少非经常性收益污染；股本由总市值/当前价反推；H2/H1转换比例是明确的季节性和交付假设，不是机械年化。
为什么用这组倍数：PE区间是House对成长持续期的情景判断；牛市情景要求业绩桥接和经营证据持续，不是市场一致目标价。
为什么用这组概率：30\%/50\%/20\%是统一筛选先验：基准情景权重最高，同时保留熊市和牛市的实质性下行/上行贡献；它是House情景权重，不是经过统计估计的客观事件概率。
证据如何参与：官方半年报预告、Q1财务包和行情/K线提供事实输入；研报/来源用于分母或可比背景。没有正权重外部目标锚，House价值不伪装为Street共识。
\begin{align*}
\mathrm{Shares} &= \frac{1044.20}{5.02}=208.0080\ \mathrm{100m\ shares}\\
\mathrm{EPS}_{H1}^d &= \frac{30.30}{208.0080}=0.1457\\
\mathrm{EPS}_s &= \mathrm{EPS}_{H1}^d(1+r_s)=0.25/0.36/0.34\\
V_s &= \mathrm{EPS}_s\times \mathrm{PE}_s=3.76/9.36/11.39\\
V_{prob} &= 0.30\times 3.76+0.50\times 9.36+0.20\times 11.39=8.09\\
P_{target} &= V_{prob}=8.09\quad(\text{外部锚权重}=0)\\
\mathrm{Upside} &= \frac{8.09}{5.02}-1=61.2\%\\
\end{align*}

\paragraph{001287 中电港：成长盈利情景 PE}
分母：H1 deducted EPS bridged with H2/H1 conversion；证据质量：中高；审计状态：PASS。
为什么选这个方法：电子具有成长或盈利持续期特征，因此模型给盈利桥接配置较高PE区间，而不是直接给主题收入估值；价格和现金流检查防止把泛化成长叙事直接变成EPS。
为什么选这个分母：使用H1扣非利润以减少非经常性收益污染；股本由总市值/当前价反推；H2/H1转换比例是明确的季节性和交付假设，不是机械年化。
为什么用这组倍数：PE区间是House对成长持续期的情景判断；牛市情景要求业绩桥接和经营证据持续，不是市场一致目标价。
为什么用这组概率：30\%/50\%/20\%是统一筛选先验：基准情景权重最高，同时保留熊市和牛市的实质性下行/上行贡献；它是House情景权重，不是经过统计估计的客观事件概率。
证据如何参与：官方半年报预告、Q1财务包和行情/K线提供事实输入；研报/来源用于分母或可比背景。没有正权重外部目标锚，House价值不伪装为Street共识。
\begin{align*}
\mathrm{Shares} &= \frac{202.51}{26.65}=7.5989\ \mathrm{100m\ shares}\\
\mathrm{EPS}_{H1}^d &= \frac{5.15}{7.5989}=0.6777\\
\mathrm{EPS}_s &= \mathrm{EPS}_{H1}^d(1+r_s)=1.15/1.84/1.56\\
V_s &= \mathrm{EPS}_s\times \mathrm{PE}_s=19.99/47.84/53.00\\
V_{prob} &= 0.30\times 19.99+0.50\times 47.84+0.20\times 53.00=40.52\\
P_{target} &= V_{prob}=40.52\quad(\text{外部锚权重}=0)\\
\mathrm{Upside} &= \frac{40.52}{26.65}-1=52.0\%\\
\end{align*}

\paragraph{600111 北方稀土：周期调整扣非 EPS / PE}
分母：H1 deducted EPS bridged with H2/H1 conversion；证据质量：中高；审计状态：PASS。
为什么选这个方法：有色金属盈利受商品价格、运价、价差或利用率周期影响，因此使用周期调整EPS/PE区间，而不是把某一个季度的PE当作全年估值。H2转换比例用于检验当前周期能否延续。
为什么选这个分母：使用H1扣非利润以减少非经常性收益污染；股本由总市值/当前价反推；H2/H1转换比例是明确的季节性和交付假设，不是机械年化。
为什么用这组倍数：熊/基准/牛PE区间是针对该周期业务类型的House假设，不是外部券商目标价；熊值较低反映周期反转，牛值较高必须由价格、销量和现金流共同确认。
为什么用这组概率：30\%/50\%/20\%是统一筛选先验：基准情景权重最高，同时保留熊市和牛市的实质性下行/上行贡献；它是House情景权重，不是经过统计估计的客观事件概率。
证据如何参与：官方半年报预告、Q1财务包和行情/K线提供事实输入；研报/来源用于分母或可比背景。没有正权重外部目标锚，House价值不伪装为Street共识。
\begin{align*}
\mathrm{Shares} &= \frac{1430.12}{39.56}=36.1507\ \mathrm{100m\ shares}\\
\mathrm{EPS}_{H1}^d &= \frac{20.30}{36.1507}=0.5615\\
\mathrm{EPS}_s &= \mathrm{EPS}_{H1}^d(1+r_s)=0.87/1.12/1.21\\
V_s &= \mathrm{EPS}_s\times \mathrm{PE}_s=6.96/13.44/19.32\\
V_{prob} &= 0.30\times 6.96+0.50\times 13.44+0.20\times 19.32=12.67\\
P_{target} &= V_{prob}=12.67\quad(\text{外部锚权重}=0)\\
\mathrm{Upside} &= \frac{12.67}{39.56}-1=-68.0\%\\
\end{align*}

\paragraph{000725 京东方A：成长盈利情景 PE}
分母：H1 deducted EPS bridged with H2/H1 conversion；证据质量：中高；审计状态：PASS。
为什么选这个方法：电子具有成长或盈利持续期特征，因此模型给盈利桥接配置较高PE区间，而不是直接给主题收入估值；价格和现金流检查防止把泛化成长叙事直接变成EPS。
为什么选这个分母：使用H1扣非利润以减少非经常性收益污染；股本由总市值/当前价反推；H2/H1转换比例是明确的季节性和交付假设，不是机械年化。
为什么用这组倍数：PE区间是House对成长持续期的情景判断；牛市情景要求业绩桥接和经营证据持续，不是市场一致目标价。
为什么用这组概率：30\%/50\%/20\%是统一筛选先验：基准情景权重最高，同时保留熊市和牛市的实质性下行/上行贡献；它是House情景权重，不是经过统计估计的客观事件概率。
证据如何参与：官方半年报预告、Q1财务包和行情/K线提供事实输入；研报/来源用于分母或可比背景。没有正权重外部目标锚，House价值不伪装为Street共识。
\begin{align*}
\mathrm{Shares} &= \frac{2363.43}{6.38}=370.4436\ \mathrm{100m\ shares}\\
\mathrm{EPS}_{H1}^d &= \frac{31.95}{370.4436}=0.0862\\
\mathrm{EPS}_s &= \mathrm{EPS}_{H1}^d(1+r_s)=0.15/0.26/0.20\\
V_s &= \mathrm{EPS}_s\times \mathrm{PE}_s=2.64/6.74/6.76\\
V_{prob} &= 0.30\times 2.64+0.50\times 6.74+0.20\times 6.76=5.51\\
P_{target} &= V_{prob}=5.51\quad(\text{外部锚权重}=0)\\
\mathrm{Upside} &= \frac{5.51}{6.38}-1=-13.6\%\\
\end{align*}

\paragraph{600176 中国巨石：常规扣非 EPS / PE}
分母：H1 deducted EPS bridged with H2/H1 conversion；证据质量：中高；审计状态：PASS。
为什么选这个方法：该标的H1扣非利润为正且未触发恢复估值阻断，因此扣非EPS是最透明的筛选分母；情景桥接避免把一个H1数字直接当作全年点预测。
为什么选这个分母：使用H1扣非利润以减少非经常性收益污染；股本由总市值/当前价反推；H2/H1转换比例是明确的季节性和交付假设，不是机械年化。
为什么用这组倍数：PE区间按公司所属行业匹配，表达估值压缩/正常/重估三种情景；这是House筛选区间，正式升级前仍需当前可审计的Street区间。
为什么用这组概率：30\%/50\%/20\%是统一筛选先验：基准情景权重最高，同时保留熊市和牛市的实质性下行/上行贡献；它是House情景权重，不是经过统计估计的客观事件概率。
证据如何参与：官方半年报预告、Q1财务包和行情/K线提供事实输入；研报/来源用于分母或可比背景。没有正权重外部目标锚，House价值不伪装为Street共识。
\begin{align*}
\mathrm{Shares} &= \frac{2121.66}{53.00}=40.0313\ \mathrm{100m\ shares}\\
\mathrm{EPS}_{H1}^d &= \frac{29.77}{40.0313}=0.7437\\
\mathrm{EPS}_s &= \mathrm{EPS}_{H1}^d(1+r_s)=1.23/1.67/1.64\\
V_s &= \mathrm{EPS}_s\times \mathrm{PE}_s=14.73/30.06/39.27\\
V_{prob} &= 0.30\times 14.73+0.50\times 30.06+0.20\times 39.27=27.30\\
P_{target} &= V_{prob}=27.30\quad(\text{外部锚权重}=0)\\
\mathrm{Upside} &= \frac{27.30}{53.00}-1=-48.5\%\\
\end{align*}

\paragraph{002266 浙富控股：常规扣非 EPS / PE}
分母：H1 deducted EPS bridged with H2/H1 conversion；证据质量：中高；审计状态：PASS。
为什么选这个方法：该标的H1扣非利润为正且未触发恢复估值阻断，因此扣非EPS是最透明的筛选分母；情景桥接避免把一个H1数字直接当作全年点预测。
为什么选这个分母：使用H1扣非利润以减少非经常性收益污染；股本由总市值/当前价反推；H2/H1转换比例是明确的季节性和交付假设，不是机械年化。
为什么用这组倍数：PE区间按公司所属行业匹配，表达估值压缩/正常/重估三种情景；这是House筛选区间，正式升级前仍需当前可审计的Street区间。
为什么用这组概率：30\%/50\%/20\%是统一筛选先验：基准情景权重最高，同时保留熊市和牛市的实质性下行/上行贡献；它是House情景权重，不是经过统计估计的客观事件概率。
证据如何参与：官方半年报预告、Q1财务包和行情/K线提供事实输入；研报/来源用于分母或可比背景。没有正权重外部目标锚，House价值不伪装为Street共识。
\begin{align*}
\mathrm{Shares} &= \frac{265.66}{5.09}=52.1925\ \mathrm{100m\ shares}\\
\mathrm{EPS}_{H1}^d &= \frac{12.50}{52.1925}=0.2395\\
\mathrm{EPS}_s &= \mathrm{EPS}_{H1}^d(1+r_s)=0.40/0.46/0.53\\
V_s &= \mathrm{EPS}_s\times \mathrm{PE}_s=3.82/8.19/12.65\\
V_{prob} &= 0.30\times 3.82+0.50\times 8.19+0.20\times 12.65=7.77\\
P_{target} &= V_{prob}=7.77\quad(\text{外部锚权重}=0)\\
\mathrm{Upside} &= \frac{7.77}{5.09}-1=52.6\%\\
\end{align*}

\paragraph{000612 焦作万方：周期调整扣非 EPS / PE}
分母：H1 deducted EPS bridged with H2/H1 conversion；证据质量：中；审计状态：PASS。
为什么选这个方法：有色金属盈利受商品价格、运价、价差或利用率周期影响，因此使用周期调整EPS/PE区间，而不是把某一个季度的PE当作全年估值。H2转换比例用于检验当前周期能否延续。
为什么选这个分母：使用H1扣非利润以减少非经常性收益污染；股本由总市值/当前价反推；H2/H1转换比例是明确的季节性和交付假设，不是机械年化。
为什么用这组倍数：熊/基准/牛PE区间是针对该周期业务类型的House假设，不是外部券商目标价；熊值较低反映周期反转，牛值较高必须由价格、销量和现金流共同确认。
为什么用这组概率：30\%/50\%/20\%是统一筛选先验：基准情景权重最高，同时保留熊市和牛市的实质性下行/上行贡献；它是House情景权重，不是经过统计估计的客观事件概率。
证据如何参与：官方半年报预告、Q1财务包和行情/K线提供事实输入；研报/来源用于分母或可比背景。没有正权重外部目标锚，House价值不伪装为Street共识。
\begin{align*}
\mathrm{Shares} &= \frac{129.23}{10.84}=11.9216\ \mathrm{100m\ shares}\\
\mathrm{EPS}_{H1}^d &= \frac{12.08}{11.9216}=1.0133\\
\mathrm{EPS}_s &= \mathrm{EPS}_{H1}^d(1+r_s)=1.57/1.87/2.18\\
V_s &= \mathrm{EPS}_s\times \mathrm{PE}_s=8.13/22.50/34.86\\
V_{prob} &= 0.30\times 8.13+0.50\times 22.50+0.20\times 34.86=20.66\\
P_{target} &= V_{prob}=20.66\quad(\text{外部锚权重}=0)\\
\mathrm{Upside} &= \frac{20.66}{10.84}-1=90.6\%\\
\end{align*}

\paragraph{605117 德业股份：成长盈利情景 PE}
分母：H1 deducted EPS bridged with H2/H1 conversion；证据质量：中高；审计状态：PASS。
为什么选这个方法：电力设备具有成长或盈利持续期特征，因此模型给盈利桥接配置较高PE区间，而不是直接给主题收入估值；价格和现金流检查防止把泛化成长叙事直接变成EPS。
为什么选这个分母：使用H1扣非利润以减少非经常性收益污染；股本由总市值/当前价反推；H2/H1转换比例是明确的季节性和交付假设，不是机械年化。
为什么用这组倍数：PE区间是House对成长持续期的情景判断；牛市情景要求业绩桥接和经营证据持续，不是市场一致目标价。
为什么用这组概率：30\%/50\%/20\%是统一筛选先验：基准情景权重最高，同时保留熊市和牛市的实质性下行/上行贡献；它是House情景权重，不是经过统计估计的客观事件概率。
证据如何参与：官方半年报预告、Q1财务包和行情/K线提供事实输入；研报/来源用于分母或可比背景。没有正权重外部目标锚，House价值不伪装为Street共识。
\begin{align*}
\mathrm{Shares} &= \frac{1087.34}{85.40}=12.7323\ \mathrm{100m\ shares}\\
\mathrm{EPS}_{H1}^d &= \frac{25.23}{12.7323}=1.9815\\
\mathrm{EPS}_s &= \mathrm{EPS}_{H1}^d(1+r_s)=3.37/5.08/4.56\\
V_s &= \mathrm{EPS}_s\times \mathrm{PE}_s=53.90/111.76/127.61\\
V_{prob} &= 0.30\times 53.90+0.50\times 111.76+0.20\times 127.61=97.57\\
P_{target} &= V_{prob}=97.57\quad(\text{外部锚权重}=0)\\
\mathrm{Upside} &= \frac{97.57}{85.40}-1=14.2\%\\
\end{align*}

\paragraph{000783 长江证券：金融 PB/ROE + 盈利混合}
分母：Q1 BPS plus H1 deducted EPS bridged to H2；证据质量：中；审计状态：PASS。
为什么选这个方法：金融公司主要通过净资产、ROE和盈利质量估值；PB保留资产负债表/ROE锚，PE提供独立盈利交叉检查，混合使用可以避免只依赖投资收益驱动的单季EPS。
为什么选这个分母：Q1 BPS作为资产分母，H1扣非EPS通过H2情景转换为熊/基准/牛盈利；这样将资产价值和盈利能力分开处理。
为什么用这组倍数：PB 0.75/1.05/1.35倍分别代表资产质量压力/正常/重估；行业PE区间提供盈利交叉检查；这些是House区间，不是外部目标价。
为什么用这组概率：30\%/50\%/20\%是统一筛选先验：基准情景权重最高，同时保留熊市和牛市的实质性下行/上行贡献；它是House情景权重，不是经过统计估计的客观事件概率。
证据如何参与：官方半年报预告、Q1财务包和行情/K线提供事实输入；研报/来源用于分母或可比背景。没有正权重外部目标锚，House价值不伪装为Street共识。
\begin{align*}
\mathrm{EPS}_{H1}^d &= \frac{31.84}{55.3002} = 0.5758\ \mathrm{CNY/share}\\
V_s &= 0.65\left(\mathrm{BPS}_1\times \mathrm{PB}_s\right)+0.35\left(\mathrm{EPS}_s\times \mathrm{PE}_s\right)\\
V_{b/m/b} &= 5.99/8.98/12.19\ \mathrm{CNY/share}\quad(\mathrm{PB}=5.00/6.99/8.99,\ \mathrm{PE}=7.83/12.67/18.14)\\
V_{prob} &= 0.30\times 5.99+0.50\times 8.98+0.20\times 12.19=8.72\\
P_{target} &= V_{prob}=8.72\quad(\text{外部锚权重}=0)\\
\mathrm{Upside} &= \frac{8.72}{9.06}-1=-3.8\%\\
\end{align*}

\paragraph{600233 圆通速递：周期调整扣非 EPS / PE}
分母：H1 deducted EPS bridged with H2/H1 conversion；证据质量：中高；审计状态：PASS。
为什么选这个方法：交通运输盈利受商品价格、运价、价差或利用率周期影响，因此使用周期调整EPS/PE区间，而不是把某一个季度的PE当作全年估值。H2转换比例用于检验当前周期能否延续。
为什么选这个分母：使用H1扣非利润以减少非经常性收益污染；股本由总市值/当前价反推；H2/H1转换比例是明确的季节性和交付假设，不是机械年化。
为什么用这组倍数：熊/基准/牛PE区间是针对该周期业务类型的House假设，不是外部券商目标价；熊值较低反映周期反转，牛值较高必须由价格、销量和现金流共同确认。
为什么用这组概率：30\%/50\%/20\%是统一筛选先验：基准情景权重最高，同时保留熊市和牛市的实质性下行/上行贡献；它是House情景权重，不是经过统计估计的客观事件概率。
证据如何参与：官方半年报预告、Q1财务包和行情/K线提供事实输入；研报/来源用于分母或可比背景。没有正权重外部目标锚，House价值不伪装为Street共识。
\begin{align*}
\mathrm{Shares} &= \frac{599.01}{17.50}=34.2291\ \mathrm{100m\ shares}\\
\mathrm{EPS}_{H1}^d &= \frac{31.90}{34.2291}=0.9320\\
\mathrm{EPS}_s &= \mathrm{EPS}_{H1}^d(1+r_s)=1.44/1.83/2.00\\
V_s &= \mathrm{EPS}_s\times \mathrm{PE}_s=11.56/21.96/32.06\\
V_{prob} &= 0.30\times 11.56+0.50\times 21.96+0.20\times 32.06=20.86\\
P_{target} &= V_{prob}=20.86\quad(\text{外部锚权重}=0)\\
\mathrm{Upside} &= \frac{20.86}{17.50}-1=19.2\%\\
\end{align*}

\paragraph{600392 盛和资源：周期调整扣非 EPS / PE}
分母：H1 deducted EPS bridged with H2/H1 conversion；证据质量：中；审计状态：PASS。
为什么选这个方法：有色金属盈利受商品价格、运价、价差或利用率周期影响，因此使用周期调整EPS/PE区间，而不是把某一个季度的PE当作全年估值。H2转换比例用于检验当前周期能否延续。
为什么选这个分母：使用H1扣非利润以减少非经常性收益污染；股本由总市值/当前价反推；H2/H1转换比例是明确的季节性和交付假设，不是机械年化。
为什么用这组倍数：熊/基准/牛PE区间是针对该周期业务类型的House假设，不是外部券商目标价；熊值较低反映周期反转，牛值较高必须由价格、销量和现金流共同确认。
为什么用这组概率：30\%/50\%/20\%是统一筛选先验：基准情景权重最高，同时保留熊市和牛市的实质性下行/上行贡献；它是House情景权重，不是经过统计估计的客观事件概率。
证据如何参与：官方半年报预告、Q1财务包和行情/K线提供事实输入；研报/来源用于分母或可比背景。没有正权重外部目标锚，House价值不伪装为Street共识。
\begin{align*}
\mathrm{Shares} &= \frac{418.22}{23.86}=17.5281\ \mathrm{100m\ shares}\\
\mathrm{EPS}_{H1}^d &= \frac{8.55}{17.5281}=0.4878\\
\mathrm{EPS}_s &= \mathrm{EPS}_{H1}^d(1+r_s)=0.76/0.90/1.05\\
V_s &= \mathrm{EPS}_s\times \mathrm{PE}_s=6.05/10.83/16.78\\
V_{prob} &= 0.30\times 6.05+0.50\times 10.83+0.20\times 16.78=10.59\\
P_{target} &= V_{prob}=10.59\quad(\text{外部锚权重}=0)\\
\mathrm{Upside} &= \frac{10.59}{23.86}-1=-55.6\%\\
\end{align*}

\paragraph{603225 新凤鸣：周期调整扣非 EPS / PE}
分母：H1 deducted EPS bridged with H2/H1 conversion；证据质量：中高；审计状态：PASS。
为什么选这个方法：基础化工盈利受商品价格、运价、价差或利用率周期影响，因此使用周期调整EPS/PE区间，而不是把某一个季度的PE当作全年估值。H2转换比例用于检验当前周期能否延续。
为什么选这个分母：使用H1扣非利润以减少非经常性收益污染；股本由总市值/当前价反推；H2/H1转换比例是明确的季节性和交付假设，不是机械年化。
为什么用这组倍数：熊/基准/牛PE区间是针对该周期业务类型的House假设，不是外部券商目标价；熊值较低反映周期反转，牛值较高必须由价格、销量和现金流共同确认。
为什么用这组概率：30\%/50\%/20\%是统一筛选先验：基准情景权重最高，同时保留熊市和牛市的实质性下行/上行贡献；它是House情景权重，不是经过统计估计的客观事件概率。
证据如何参与：官方半年报预告、Q1财务包和行情/K线提供事实输入；研报/来源用于分母或可比背景。没有正权重外部目标锚，House价值不伪装为Street共识。
\begin{align*}
\mathrm{Shares} &= \frac{301.45}{17.97}=16.7752\ \mathrm{100m\ shares}\\
\mathrm{EPS}_{H1}^d &= \frac{15.15}{16.7752}=0.9031\\
\mathrm{EPS}_s &= \mathrm{EPS}_{H1}^d(1+r_s)=1.40/1.67/1.94\\
V_s &= \mathrm{EPS}_s\times \mathrm{PE}_s=13.48/23.39/34.95\\
V_{prob} &= 0.30\times 13.48+0.50\times 23.39+0.20\times 34.95=22.73\\
P_{target} &= V_{prob}=22.73\quad(\text{外部锚权重}=0)\\
\mathrm{Upside} &= \frac{22.73}{17.97}-1=26.5\%\\
\end{align*}

\paragraph{002008 大族激光：常规扣非 EPS / PE}
分母：H1 deducted EPS bridged with H2/H1 conversion；证据质量：中高；审计状态：PASS。
为什么选这个方法：该标的H1扣非利润为正且未触发恢复估值阻断，因此扣非EPS是最透明的筛选分母；情景桥接避免把一个H1数字直接当作全年点预测。
为什么选这个分母：使用H1扣非利润以减少非经常性收益污染；股本由总市值/当前价反推；H2/H1转换比例是明确的季节性和交付假设，不是机械年化。
为什么用这组倍数：PE区间按公司所属行业匹配，表达估值压缩/正常/重估三种情景；这是House筛选区间，正式升级前仍需当前可审计的Street区间。
为什么用这组概率：30\%/50\%/20\%是统一筛选先验：基准情景权重最高，同时保留熊市和牛市的实质性下行/上行贡献；它是House情景权重，不是经过统计估计的客观事件概率。
证据如何参与：官方半年报预告、Q1财务包和行情/K线提供事实输入；研报/来源用于分母或可比背景。没有正权重外部目标锚，House价值不伪装为Street共识。
\begin{align*}
\mathrm{Shares} &= \frac{1177.45}{114.36}=10.2960\ \mathrm{100m\ shares}\\
\mathrm{EPS}_{H1}^d &= \frac{14.00}{10.2960}=1.3598\\
\mathrm{EPS}_s &= \mathrm{EPS}_{H1}^d(1+r_s)=2.24/2.62/2.99\\
V_s &= \mathrm{EPS}_s\times \mathrm{PE}_s=31.41/52.40/83.76\\
V_{prob} &= 0.30\times 31.41+0.50\times 52.40+0.20\times 83.76=52.38\\
P_{target} &= V_{prob}=52.38\quad(\text{外部锚权重}=0)\\
\mathrm{Upside} &= \frac{52.38}{114.36}-1=-54.2\%\\
\end{align*}

\paragraph{000685 中山公用：常规扣非 EPS / PE}
分母：H1 deducted EPS bridged with H2/H1 conversion；证据质量：中高；审计状态：PASS。
为什么选这个方法：该标的H1扣非利润为正且未触发恢复估值阻断，因此扣非EPS是最透明的筛选分母；情景桥接避免把一个H1数字直接当作全年点预测。
为什么选这个分母：使用H1扣非利润以减少非经常性收益污染；股本由总市值/当前价反推；H2/H1转换比例是明确的季节性和交付假设，不是机械年化。
为什么用这组倍数：PE区间按公司所属行业匹配，表达估值压缩/正常/重估三种情景；这是House筛选区间，正式升级前仍需当前可审计的Street区间。
为什么用这组概率：30\%/50\%/20\%是统一筛选先验：基准情景权重最高，同时保留熊市和牛市的实质性下行/上行贡献；它是House情景权重，不是经过统计估计的客观事件概率。
证据如何参与：官方半年报预告、Q1财务包和行情/K线提供事实输入；研报/来源用于分母或可比背景。没有正权重外部目标锚，House价值不伪装为Street共识。
\begin{align*}
\mathrm{Shares} &= \frac{164.24}{11.19}=14.6774\ \mathrm{100m\ shares}\\
\mathrm{EPS}_{H1}^d &= \frac{14.27}{14.6774}=0.9724\\
\mathrm{EPS}_s &= \mathrm{EPS}_{H1}^d(1+r_s)=1.60/1.85/2.14\\
V_s &= \mathrm{EPS}_s\times \mathrm{PE}_s=8.39/33.26/51.34\\
V_{prob} &= 0.30\times 8.39+0.50\times 33.26+0.20\times 51.34=29.41\\
P_{target} &= V_{prob}=29.41\quad(\text{外部锚权重}=0)\\
\mathrm{Upside} &= \frac{29.41}{11.19}-1=162.8\%\\
\end{align*}

\paragraph{600909 华安证券：金融 PB/ROE + 盈利混合}
分母：Q1 BPS plus H1 deducted EPS bridged to H2；证据质量：中高；审计状态：PASS。
为什么选这个方法：金融公司主要通过净资产、ROE和盈利质量估值；PB保留资产负债表/ROE锚，PE提供独立盈利交叉检查，混合使用可以避免只依赖投资收益驱动的单季EPS。
为什么选这个分母：Q1 BPS作为资产分母，H1扣非EPS通过H2情景转换为熊/基准/牛盈利；这样将资产价值和盈利能力分开处理。
为什么用这组倍数：PB 0.75/1.05/1.35倍分别代表资产质量压力/正常/重估；行业PE区间提供盈利交叉检查；这些是House区间，不是外部目标价。
为什么用这组概率：30\%/50\%/20\%是统一筛选先验：基准情景权重最高，同时保留熊市和牛市的实质性下行/上行贡献；它是House情景权重，不是经过统计估计的客观事件概率。
证据如何参与：官方半年报预告、Q1财务包和行情/K线提供事实输入；研报/来源用于分母或可比背景。没有正权重外部目标锚，House价值不伪装为Street共识。
\begin{align*}
\mathrm{EPS}_{H1}^d &= \frac{21.60}{50.2421} = 0.4299\ \mathrm{CNY/share}\\
V_s &= 0.65\left(\mathrm{BPS}_1\times \mathrm{PB}_s\right)+0.35\left(\mathrm{EPS}_s\times \mathrm{PE}_s\right)\\
V_{b/m/b} &= 4.61/6.89/9.35\ \mathrm{CNY/share}\quad(\mathrm{PB}=3.94/5.51/7.09,\ \mathrm{PE}=5.85/9.46/13.54)\\
V_{prob} &= 0.30\times 4.61+0.50\times 6.89+0.20\times 9.35=6.70\\
P_{target} &= V_{prob}=6.70\quad(\text{外部锚权重}=0)\\
\mathrm{Upside} &= \frac{6.70}{9.50}-1=-29.5\%\\
\end{align*}

\paragraph{000833 粤桂股份：常规扣非 EPS / PE}
分母：H1 deducted EPS bridged with H2/H1 conversion；证据质量：中；审计状态：PASS。
为什么选这个方法：该标的H1扣非利润为正且未触发恢复估值阻断，因此扣非EPS是最透明的筛选分母；情景桥接避免把一个H1数字直接当作全年点预测。
为什么选这个分母：使用H1扣非利润以减少非经常性收益污染；股本由总市值/当前价反推；H2/H1转换比例是明确的季节性和交付假设，不是机械年化。
为什么用这组倍数：PE区间按公司所属行业匹配，表达估值压缩/正常/重估三种情景；这是House筛选区间，正式升级前仍需当前可审计的Street区间。
为什么用这组概率：30\%/50\%/20\%是统一筛选先验：基准情景权重最高，同时保留熊市和牛市的实质性下行/上行贡献；它是House情景权重，不是经过统计估计的客观事件概率。
证据如何参与：官方半年报预告、Q1财务包和行情/K线提供事实输入；研报/来源用于分母或可比背景。没有正权重外部目标锚，House价值不伪装为Street共识。
\begin{align*}
\mathrm{Shares} &= \frac{160.98}{20.07}=8.0209\ \mathrm{100m\ shares}\\
\mathrm{EPS}_{H1}^d &= \frac{6.53}{8.0209}=0.8141\\
\mathrm{EPS}_s &= \mathrm{EPS}_{H1}^d(1+r_s)=1.34/1.55/1.79\\
V_s &= \mathrm{EPS}_s\times \mathrm{PE}_s=15.05/27.84/42.99\\
V_{prob} &= 0.30\times 15.05+0.50\times 27.84+0.20\times 42.99=27.03\\
P_{target} &= V_{prob}=27.03\quad(\text{外部锚权重}=0)\\
\mathrm{Upside} &= \frac{27.03}{20.07}-1=34.7\%\\
\end{align*}

\paragraph{300037 新宙邦：成长盈利情景 PE}
分母：H1 deducted EPS bridged with H2/H1 conversion；证据质量：高；审计状态：PASS。
为什么选这个方法：电力设备具有成长或盈利持续期特征，因此模型给盈利桥接配置较高PE区间，而不是直接给主题收入估值；价格和现金流检查防止把泛化成长叙事直接变成EPS。
为什么选这个分母：使用H1扣非利润以减少非经常性收益污染；股本由总市值/当前价反推；H2/H1转换比例是明确的季节性和交付假设，不是机械年化。
为什么用这组倍数：PE区间是House对成长持续期的情景判断；牛市情景要求业绩桥接和经营证据持续，不是市场一致目标价。
为什么用这组概率：30\%/50\%/20\%是统一筛选先验：基准情景权重最高，同时保留熊市和牛市的实质性下行/上行贡献；它是House情景权重，不是经过统计估计的客观事件概率。
证据如何参与：官方半年报预告、Q1财务包和行情/K线提供事实输入；研报/来源用于分母或可比背景。外部目标价123.30元，赋予正权重10\%。
\begin{align*}
\mathrm{Shares} &= \frac{506.61}{67.20}=7.5388\ \mathrm{100m\ shares}\\
\mathrm{EPS}_{H1}^d &= \frac{9.85}{7.5388}=1.3066\\
\mathrm{EPS}_s &= \mathrm{EPS}_{H1}^d(1+r_s)=2.22/3.21/3.01\\
V_s &= \mathrm{EPS}_s\times \mathrm{PE}_s=35.54/70.62/84.14\\
V_{prob} &= 0.30\times 35.54+0.50\times 70.62+0.20\times 84.14=62.80\\
P_{target} &= 62.80\times 0.90+123.30\times 0.10=68.85\\
\mathrm{Upside} &= \frac{68.85}{67.20}-1=2.5\%\\
\end{align*}

\paragraph{000060 中金岭南：周期调整扣非 EPS / PE}
分母：H1 deducted EPS bridged with H2/H1 conversion；证据质量：中高；审计状态：PASS。
为什么选这个方法：有色金属盈利受商品价格、运价、价差或利用率周期影响，因此使用周期调整EPS/PE区间，而不是把某一个季度的PE当作全年估值。H2转换比例用于检验当前周期能否延续。
为什么选这个分母：使用H1扣非利润以减少非经常性收益污染；股本由总市值/当前价反推；H2/H1转换比例是明确的季节性和交付假设，不是机械年化。
为什么用这组倍数：熊/基准/牛PE区间是针对该周期业务类型的House假设，不是外部券商目标价；熊值较低反映周期反转，牛值较高必须由价格、销量和现金流共同确认。
为什么用这组概率：30\%/50\%/20\%是统一筛选先验：基准情景权重最高，同时保留熊市和牛市的实质性下行/上行贡献；它是House情景权重，不是经过统计估计的客观事件概率。
证据如何参与：官方半年报预告、Q1财务包和行情/K线提供事实输入；研报/来源用于分母或可比背景。没有正权重外部目标锚，House价值不伪装为Street共识。
\begin{align*}
\mathrm{Shares} &= \frac{278.02}{6.26}=44.4121\ \mathrm{100m\ shares}\\
\mathrm{EPS}_{H1}^d &= \frac{11.05}{44.4121}=0.2488\\
\mathrm{EPS}_s &= \mathrm{EPS}_{H1}^d(1+r_s)=0.39/0.51/0.53\\
V_s &= \mathrm{EPS}_s\times \mathrm{PE}_s=3.09/6.12/8.56\\
V_{prob} &= 0.30\times 3.09+0.50\times 6.12+0.20\times 8.56=5.70\\
P_{target} &= V_{prob}=5.70\quad(\text{外部锚权重}=0)\\
\mathrm{Upside} &= \frac{5.70}{6.26}-1=-8.9\%\\
\end{align*}

\paragraph{600522 中天科技：成长盈利情景 PE}
分母：H1 deducted EPS bridged with H2/H1 conversion；证据质量：中；审计状态：PASS。
为什么选这个方法：通信具有成长或盈利持续期特征，因此模型给盈利桥接配置较高PE区间，而不是直接给主题收入估值；价格和现金流检查防止把泛化成长叙事直接变成EPS。
为什么选这个分母：使用H1扣非利润以减少非经常性收益污染；股本由总市值/当前价反推；H2/H1转换比例是明确的季节性和交付假设，不是机械年化。
为什么用这组倍数：PE区间是House对成长持续期的情景判断；牛市情景要求业绩桥接和经营证据持续，不是市场一致目标价。
为什么用这组概率：30\%/50\%/20\%是统一筛选先验：基准情景权重最高，同时保留熊市和牛市的实质性下行/上行贡献；它是House情景权重，不是经过统计估计的客观事件概率。
证据如何参与：官方半年报预告、Q1财务包和行情/K线提供事实输入；研报/来源用于分母或可比背景。没有正权重外部目标锚，House价值不伪装为Street共识。
\begin{align*}
\mathrm{Shares} &= \frac{1389.07}{40.70}=34.1295\ \mathrm{100m\ shares}\\
\mathrm{EPS}_{H1}^d &= \frac{22.33}{34.1295}=0.6543\\
\mathrm{EPS}_s &= \mathrm{EPS}_{H1}^d(1+r_s)=1.11/1.31/1.50\\
V_s &= \mathrm{EPS}_s\times \mathrm{PE}_s=20.02/34.02/51.16\\
V_{prob} &= 0.30\times 20.02+0.50\times 34.02+0.20\times 51.16=33.25\\
P_{target} &= V_{prob}=33.25\quad(\text{外部锚权重}=0)\\
\mathrm{Upside} &= \frac{33.25}{40.70}-1=-18.3\%\\
\end{align*}

\paragraph{002468 申通快递：周期调整扣非 EPS / PE}
分母：H1 deducted EPS bridged with H2/H1 conversion；证据质量：中高；审计状态：PASS。
为什么选这个方法：交通运输盈利受商品价格、运价、价差或利用率周期影响，因此使用周期调整EPS/PE区间，而不是把某一个季度的PE当作全年估值。H2转换比例用于检验当前周期能否延续。
为什么选这个分母：使用H1扣非利润以减少非经常性收益污染；股本由总市值/当前价反推；H2/H1转换比例是明确的季节性和交付假设，不是机械年化。
为什么用这组倍数：熊/基准/牛PE区间是针对该周期业务类型的House假设，不是外部券商目标价；熊值较低反映周期反转，牛值较高必须由价格、销量和现金流共同确认。
为什么用这组概率：30\%/50\%/20\%是统一筛选先验：基准情景权重最高，同时保留熊市和牛市的实质性下行/上行贡献；它是House情景权重，不是经过统计估计的客观事件概率。
证据如何参与：官方半年报预告、Q1财务包和行情/K线提供事实输入；研报/来源用于分母或可比背景。没有正权重外部目标锚，House价值不伪装为Street共识。
\begin{align*}
\mathrm{Shares} &= \frac{232.99}{15.22}=15.3081\ \mathrm{100m\ shares}\\
\mathrm{EPS}_{H1}^d &= \frac{10.05}{15.3081}=0.6565\\
\mathrm{EPS}_s &= \mathrm{EPS}_{H1}^d(1+r_s)=1.02/1.24/1.41\\
V_s &= \mathrm{EPS}_s\times \mathrm{PE}_s=8.14/14.88/22.58\\
V_{prob} &= 0.30\times 8.14+0.50\times 14.88+0.20\times 22.58=14.40\\
P_{target} &= V_{prob}=14.40\quad(\text{外部锚权重}=0)\\
\mathrm{Upside} &= \frac{14.40}{15.22}-1=-5.4\%\\
\end{align*}

\paragraph{600267 海正药业：成长盈利情景 PE}
分母：H1 deducted EPS bridged with H2/H1 conversion；证据质量：中；审计状态：PASS。
为什么选这个方法：医药生物具有成长或盈利持续期特征，因此模型给盈利桥接配置较高PE区间，而不是直接给主题收入估值；价格和现金流检查防止把泛化成长叙事直接变成EPS。
为什么选这个分母：使用H1扣非利润以减少非经常性收益污染；股本由总市值/当前价反推；H2/H1转换比例是明确的季节性和交付假设，不是机械年化。
为什么用这组倍数：PE区间是House对成长持续期的情景判断；牛市情景要求业绩桥接和经营证据持续，不是市场一致目标价。
为什么用这组概率：30\%/50\%/20\%是统一筛选先验：基准情景权重最高，同时保留熊市和牛市的实质性下行/上行贡献；它是House情景权重，不是经过统计估计的客观事件概率。
证据如何参与：官方半年报预告、Q1财务包和行情/K线提供事实输入；研报/来源用于分母或可比背景。没有正权重外部目标锚，House价值不伪装为Street共识。
\begin{align*}
\mathrm{Shares} &= \frac{141.40}{11.98}=11.8030\ \mathrm{100m\ shares}\\
\mathrm{EPS}_{H1}^d &= \frac{5.10}{11.8030}=0.4321\\
\mathrm{EPS}_s &= \mathrm{EPS}_{H1}^d(1+r_s)=0.73/0.86/0.99\\
V_s &= \mathrm{EPS}_s\times \mathrm{PE}_s=8.98/21.60/31.80\\
V_{prob} &= 0.30\times 8.98+0.50\times 21.60+0.20\times 31.80=19.85\\
P_{target} &= V_{prob}=19.85\quad(\text{外部锚权重}=0)\\
\mathrm{Upside} &= \frac{19.85}{11.98}-1=65.7\%\\
\end{align*}

\paragraph{600906 财达证券：金融 PB/ROE + 盈利混合}
分母：Q1 BPS plus H1 deducted EPS bridged to H2；证据质量：中；审计状态：PASS。
为什么选这个方法：金融公司主要通过净资产、ROE和盈利质量估值；PB保留资产负债表/ROE锚，PE提供独立盈利交叉检查，混合使用可以避免只依赖投资收益驱动的单季EPS。
为什么选这个分母：Q1 BPS作为资产分母，H1扣非EPS通过H2情景转换为熊/基准/牛盈利；这样将资产价值和盈利能力分开处理。
为什么用这组倍数：PB 0.75/1.05/1.35倍分别代表资产质量压力/正常/重估；行业PE区间提供盈利交叉检查；这些是House区间，不是外部目标价。
为什么用这组概率：30\%/50\%/20\%是统一筛选先验：基准情景权重最高，同时保留熊市和牛市的实质性下行/上行贡献；它是House情景权重，不是经过统计估计的客观事件概率。
证据如何参与：官方半年报预告、Q1财务包和行情/K线提供事实输入；研报/来源用于分母或可比背景。没有正权重外部目标锚，House价值不伪装为Street共识。
\begin{align*}
\mathrm{EPS}_{H1}^d &= \frac{7.60}{32.4497} = 0.2344\ \mathrm{CNY/share}\\
V_s &= 0.65\left(\mathrm{BPS}_1\times \mathrm{PB}_s\right)+0.35\left(\mathrm{EPS}_s\times \mathrm{PE}_s\right)\\
V_{b/m/b} &= 3.00/4.44/5.97\ \mathrm{CNY/share}\quad(\mathrm{PB}=2.89/4.05/5.20,\ \mathrm{PE}=3.19/5.16/7.38)\\
V_{prob} &= 0.30\times 3.00+0.50\times 4.44+0.20\times 5.97=4.31\\
P_{target} &= V_{prob}=4.31\quad(\text{外部锚权重}=0)\\
\mathrm{Upside} &= \frac{4.31}{7.25}-1=-40.6\%\\
\end{align*}

\paragraph{603268 松发股份：成长盈利情景 PE}
分母：H1 deducted EPS bridged with H2/H1 conversion；证据质量：中高；审计状态：PASS。
为什么选这个方法：国防军工具有成长或盈利持续期特征，因此模型给盈利桥接配置较高PE区间，而不是直接给主题收入估值；价格和现金流检查防止把泛化成长叙事直接变成EPS。
为什么选这个分母：使用H1扣非利润以减少非经常性收益污染；股本由总市值/当前价反推；H2/H1转换比例是明确的季节性和交付假设，不是机械年化。
为什么用这组倍数：PE区间是House对成长持续期的情景判断；牛市情景要求业绩桥接和经营证据持续，不是市场一致目标价。
为什么用这组概率：30\%/50\%/20\%是统一筛选先验：基准情景权重最高，同时保留熊市和牛市的实质性下行/上行贡献；它是House情景权重，不是经过统计估计的客观事件概率。
证据如何参与：官方半年报预告、Q1财务包和行情/K线提供事实输入；研报/来源用于分母或可比背景。没有正权重外部目标锚，House价值不伪装为Street共识。
\begin{align*}
\mathrm{Shares} &= \frac{1557.71}{160.46}=9.7078\ \mathrm{100m\ shares}\\
\mathrm{EPS}_{H1}^d &= \frac{35.00}{9.7078}=3.6054\\
\mathrm{EPS}_s &= \mathrm{EPS}_{H1}^d(1+r_s)=6.13/7.40/8.65\\
V_s &= \mathrm{EPS}_s\times \mathrm{PE}_s=120.34/259.00/389.38\\
V_{prob} &= 0.30\times 120.34+0.50\times 259.00+0.20\times 389.38=243.48\\
P_{target} &= V_{prob}=243.48\quad(\text{外部锚权重}=0)\\
\mathrm{Upside} &= \frac{243.48}{160.46}-1=51.7\%\\
\end{align*}

\paragraph{002458 益生股份：常规扣非 EPS / PE}
分母：H1 deducted EPS bridged with H2/H1 conversion；证据质量：中；审计状态：PASS。
为什么选这个方法：该标的H1扣非利润为正且未触发恢复估值阻断，因此扣非EPS是最透明的筛选分母；情景桥接避免把一个H1数字直接当作全年点预测。
为什么选这个分母：使用H1扣非利润以减少非经常性收益污染；股本由总市值/当前价反推；H2/H1转换比例是明确的季节性和交付假设，不是机械年化。
为什么用这组倍数：PE区间按公司所属行业匹配，表达估值压缩/正常/重估三种情景；这是House筛选区间，正式升级前仍需当前可审计的Street区间。
为什么用这组概率：30\%/50\%/20\%是统一筛选先验：基准情景权重最高，同时保留熊市和牛市的实质性下行/上行贡献；它是House情景权重，不是经过统计估计的客观事件概率。
证据如何参与：官方半年报预告、Q1财务包和行情/K线提供事实输入；研报/来源用于分母或可比背景。没有正权重外部目标锚，House价值不伪装为Street共识。
\begin{align*}
\mathrm{Shares} &= \frac{129.09}{9.02}=14.3115\ \mathrm{100m\ shares}\\
\mathrm{EPS}_{H1}^d &= \frac{2.83}{14.3115}=0.1977\\
\mathrm{EPS}_s &= \mathrm{EPS}_{H1}^d(1+r_s)=0.33/0.49/0.43\\
V_s &= \mathrm{EPS}_s\times \mathrm{PE}_s=3.92/8.82/10.44\\
V_{prob} &= 0.30\times 3.92+0.50\times 8.82+0.20\times 10.44=7.67\\
P_{target} &= V_{prob}=7.67\quad(\text{外部锚权重}=0)\\
\mathrm{Upside} &= \frac{7.67}{9.02}-1=-15.0\%\\
\end{align*}

\paragraph{002384 东山精密：成长盈利情景 PE}
分母：H1 deducted EPS bridged with H2/H1 conversion；证据质量：中高；审计状态：PASS。
为什么选这个方法：电子具有成长或盈利持续期特征，因此模型给盈利桥接配置较高PE区间，而不是直接给主题收入估值；价格和现金流检查防止把泛化成长叙事直接变成EPS。
为什么选这个分母：使用H1扣非利润以减少非经常性收益污染；股本由总市值/当前价反推；H2/H1转换比例是明确的季节性和交付假设，不是机械年化。
为什么用这组倍数：PE区间是House对成长持续期的情景判断；牛市情景要求业绩桥接和经营证据持续，不是市场一致目标价。
为什么用这组概率：30\%/50\%/20\%是统一筛选先验：基准情景权重最高，同时保留熊市和牛市的实质性下行/上行贡献；它是House情景权重，不是经过统计估计的客观事件概率。
证据如何参与：官方半年报预告、Q1财务包和行情/K线提供事实输入；研报/来源用于分母或可比背景。没有正权重外部目标锚，House价值不伪装为Street共识。
\begin{align*}
\mathrm{Shares} &= \frac{4807.79}{262.49}=18.3161\ \mathrm{100m\ shares}\\
\mathrm{EPS}_{H1}^d &= \frac{24.50}{18.3161}=1.3376\\
\mathrm{EPS}_s &= \mathrm{EPS}_{H1}^d(1+r_s)=2.27/3.76/3.08\\
V_s &= \mathrm{EPS}_s\times \mathrm{PE}_s=40.93/97.76/104.60\\
V_{prob} &= 0.30\times 40.93+0.50\times 97.76+0.20\times 104.60=82.08\\
P_{target} &= V_{prob}=82.08\quad(\text{外部锚权重}=0)\\
\mathrm{Upside} &= \frac{82.08}{262.49}-1=-68.7\%\\
\end{align*}

\paragraph{000977 浪潮信息：成长盈利情景 PE}
分母：H1 deducted EPS bridged with H2/H1 conversion；证据质量：中高；审计状态：PASS。
为什么选这个方法：计算机具有成长或盈利持续期特征，因此模型给盈利桥接配置较高PE区间，而不是直接给主题收入估值；价格和现金流检查防止把泛化成长叙事直接变成EPS。
为什么选这个分母：使用H1扣非利润以减少非经常性收益污染；股本由总市值/当前价反推；H2/H1转换比例是明确的季节性和交付假设，不是机械年化。
为什么用这组倍数：PE区间是House对成长持续期的情景判断；牛市情景要求业绩桥接和经营证据持续，不是市场一致目标价。
为什么用这组概率：30\%/50\%/20\%是统一筛选先验：基准情景权重最高，同时保留熊市和牛市的实质性下行/上行贡献；它是House情景权重，不是经过统计估计的客观事件概率。
证据如何参与：官方半年报预告、Q1财务包和行情/K线提供事实输入；研报/来源用于分母或可比背景。没有正权重外部目标锚，House价值不伪装为Street共识。
\begin{align*}
\mathrm{Shares} &= \frac{1243.80}{84.70}=14.6848\ \mathrm{100m\ shares}\\
\mathrm{EPS}_{H1}^d &= \frac{23.05}{14.6848}=1.5697\\
\mathrm{EPS}_s &= \mathrm{EPS}_{H1}^d(1+r_s)=2.67/3.22/3.77\\
V_s &= \mathrm{EPS}_s\times \mathrm{PE}_s=48.03/90.10/143.15\\
V_{prob} &= 0.30\times 48.03+0.50\times 90.10+0.20\times 143.15=88.09\\
P_{target} &= V_{prob}=88.09\quad(\text{外部锚权重}=0)\\
\mathrm{Upside} &= \frac{88.09}{84.70}-1=4.0\%\\
\end{align*}

\paragraph{603127 昭衍新药：成长盈利情景 PE}
分母：H1 deducted EPS bridged with H2/H1 conversion；证据质量：中高；审计状态：PASS。
为什么选这个方法：医药生物具有成长或盈利持续期特征，因此模型给盈利桥接配置较高PE区间，而不是直接给主题收入估值；价格和现金流检查防止把泛化成长叙事直接变成EPS。
为什么选这个分母：使用H1扣非利润以减少非经常性收益污染；股本由总市值/当前价反推；H2/H1转换比例是明确的季节性和交付假设，不是机械年化。
为什么用这组倍数：PE区间是House对成长持续期的情景判断；牛市情景要求业绩桥接和经营证据持续，不是市场一致目标价。
为什么用这组概率：30\%/50\%/20\%是统一筛选先验：基准情景权重最高，同时保留熊市和牛市的实质性下行/上行贡献；它是House情景权重，不是经过统计估计的客观事件概率。
证据如何参与：官方半年报预告、Q1财务包和行情/K线提供事实输入；研报/来源用于分母或可比背景。没有正权重外部目标锚，House价值不伪装为Street共识。
\begin{align*}
\mathrm{Shares} &= \frac{379.32}{50.62}=7.4935\ \mathrm{100m\ shares}\\
\mathrm{EPS}_{H1}^d &= \frac{7.01}{7.4935}=0.9361\\
\mathrm{EPS}_s &= \mathrm{EPS}_{H1}^d(1+r_s)=1.59/1.87/2.15\\
V_s &= \mathrm{EPS}_s\times \mathrm{PE}_s=28.64/46.80/68.89\\
V_{prob} &= 0.30\times 28.64+0.50\times 46.80+0.20\times 68.89=45.77\\
P_{target} &= V_{prob}=45.77\quad(\text{外部锚权重}=0)\\
\mathrm{Upside} &= \frac{45.77}{50.62}-1=-9.6\%\\
\end{align*}

\paragraph{301200 大族数控：常规扣非 EPS / PE}
分母：H1 deducted EPS bridged with H2/H1 conversion；证据质量：中高；审计状态：PASS。
为什么选这个方法：该标的H1扣非利润为正且未触发恢复估值阻断，因此扣非EPS是最透明的筛选分母；情景桥接避免把一个H1数字直接当作全年点预测。
为什么选这个分母：使用H1扣非利润以减少非经常性收益污染；股本由总市值/当前价反推；H2/H1转换比例是明确的季节性和交付假设，不是机械年化。
为什么用这组倍数：PE区间按公司所属行业匹配，表达估值压缩/正常/重估三种情景；这是House筛选区间，正式升级前仍需当前可审计的Street区间。
为什么用这组概率：30\%/50\%/20\%是统一筛选先验：基准情景权重最高，同时保留熊市和牛市的实质性下行/上行贡献；它是House情景权重，不是经过统计估计的客观事件概率。
证据如何参与：官方半年报预告、Q1财务包和行情/K线提供事实输入；研报/来源用于分母或可比背景。没有正权重外部目标锚，House价值不伪装为Street共识。
\begin{align*}
\mathrm{Shares} &= \frac{1525.70}{312.00}=4.8901\ \mathrm{100m\ shares}\\
\mathrm{EPS}_{H1}^d &= \frac{9.50}{4.8901}=1.9427\\
\mathrm{EPS}_s &= \mathrm{EPS}_{H1}^d(1+r_s)=3.21/4.25/4.27\\
V_s &= \mathrm{EPS}_s\times \mathrm{PE}_s=44.88/85.00/119.67\\
V_{prob} &= 0.30\times 44.88+0.50\times 85.00+0.20\times 119.67=79.90\\
P_{target} &= V_{prob}=79.90\quad(\text{外部锚权重}=0)\\
\mathrm{Upside} &= \frac{79.90}{312.00}-1=-74.4\%\\
\end{align*}

\paragraph{000938 紫光股份：成长盈利情景 PE}
分母：H1 deducted EPS bridged with H2/H1 conversion；证据质量：中高；审计状态：PASS。
为什么选这个方法：计算机具有成长或盈利持续期特征，因此模型给盈利桥接配置较高PE区间，而不是直接给主题收入估值；价格和现金流检查防止把泛化成长叙事直接变成EPS。
为什么选这个分母：使用H1扣非利润以减少非经常性收益污染；股本由总市值/当前价反推；H2/H1转换比例是明确的季节性和交付假设，不是机械年化。
为什么用这组倍数：PE区间是House对成长持续期的情景判断；牛市情景要求业绩桥接和经营证据持续，不是市场一致目标价。
为什么用这组概率：30\%/50\%/20\%是统一筛选先验：基准情景权重最高，同时保留熊市和牛市的实质性下行/上行贡献；它是House情景权重，不是经过统计估计的客观事件概率。
证据如何参与：官方半年报预告、Q1财务包和行情/K线提供事实输入；研报/来源用于分母或可比背景。没有正权重外部目标锚，House价值不伪装为Street共识。
\begin{align*}
\mathrm{Shares} &= \frac{1016.19}{35.53}=28.6009\ \mathrm{100m\ shares}\\
\mathrm{EPS}_{H1}^d &= \frac{18.15}{28.6009}=0.6346\\
\mathrm{EPS}_s &= \mathrm{EPS}_{H1}^d(1+r_s)=1.08/1.51/1.52\\
V_s &= \mathrm{EPS}_s\times \mathrm{PE}_s=19.42/42.28/57.88\\
V_{prob} &= 0.30\times 19.42+0.50\times 42.28+0.20\times 57.88=38.54\\
P_{target} &= V_{prob}=38.54\quad(\text{外部锚权重}=0)\\
\mathrm{Upside} &= \frac{38.54}{35.53}-1=8.5\%\\
\end{align*}

\paragraph{002463 沪电股份：成长盈利情景 PE}
分母：H1 deducted EPS bridged with H2/H1 conversion；证据质量：中高；审计状态：PASS。
为什么选这个方法：电子具有成长或盈利持续期特征，因此模型给盈利桥接配置较高PE区间，而不是直接给主题收入估值；价格和现金流检查防止把泛化成长叙事直接变成EPS。
为什么选这个分母：使用H1扣非利润以减少非经常性收益污染；股本由总市值/当前价反推；H2/H1转换比例是明确的季节性和交付假设，不是机械年化。
为什么用这组倍数：PE区间是House对成长持续期的情景判断；牛市情景要求业绩桥接和经营证据持续，不是市场一致目标价。
为什么用这组概率：30\%/50\%/20\%是统一筛选先验：基准情景权重最高，同时保留熊市和牛市的实质性下行/上行贡献；它是House情景权重，不是经过统计估计的客观事件概率。
证据如何参与：官方半年报预告、Q1财务包和行情/K线提供事实输入；研报/来源用于分母或可比背景。没有正权重外部目标锚，House价值不伪装为Street共识。
\begin{align*}
\mathrm{Shares} &= \frac{2640.80}{137.23}=19.2436\ \mathrm{100m\ shares}\\
\mathrm{EPS}_{H1}^d &= \frac{28.05}{19.2436}=1.4576\\
\mathrm{EPS}_s &= \mathrm{EPS}_{H1}^d(1+r_s)=2.48/2.99/3.35\\
V_s &= \mathrm{EPS}_s\times \mathrm{PE}_s=44.60/77.74/113.99\\
V_{prob} &= 0.30\times 44.60+0.50\times 77.74+0.20\times 113.99=75.05\\
P_{target} &= V_{prob}=75.05\quad(\text{外部锚权重}=0)\\
\mathrm{Upside} &= \frac{75.05}{137.23}-1=-45.3\%\\
\end{align*}

\paragraph{688183 生益电子：成长盈利情景 PE}
分母：H1 deducted EPS bridged with H2/H1 conversion；证据质量：中；审计状态：PASS。
为什么选这个方法：电子具有成长或盈利持续期特征，因此模型给盈利桥接配置较高PE区间，而不是直接给主题收入估值；价格和现金流检查防止把泛化成长叙事直接变成EPS。
为什么选这个分母：使用H1扣非利润以减少非经常性收益污染；股本由总市值/当前价反推；H2/H1转换比例是明确的季节性和交付假设，不是机械年化。
为什么用这组倍数：PE区间是House对成长持续期的情景判断；牛市情景要求业绩桥接和经营证据持续，不是市场一致目标价。
为什么用这组概率：30\%/50\%/20\%是统一筛选先验：基准情景权重最高，同时保留熊市和牛市的实质性下行/上行贡献；它是House情景权重，不是经过统计估计的客观事件概率。
证据如何参与：官方半年报预告、Q1财务包和行情/K线提供事实输入；研报/来源用于分母或可比背景。没有正权重外部目标锚，House价值不伪装为Street共识。
\begin{align*}
\mathrm{Shares} &= \frac{1079.40}{128.87}=8.3759\ \mathrm{100m\ shares}\\
\mathrm{EPS}_{H1}^d &= \frac{11.09}{8.3759}=1.3243\\
\mathrm{EPS}_s &= \mathrm{EPS}_{H1}^d(1+r_s)=2.25/2.65/3.05\\
V_s &= \mathrm{EPS}_s\times \mathrm{PE}_s=40.52/68.86/103.56\\
V_{prob} &= 0.30\times 40.52+0.50\times 68.86+0.20\times 103.56=67.30\\
P_{target} &= V_{prob}=67.30\quad(\text{外部锚权重}=0)\\
\mathrm{Upside} &= \frac{67.30}{128.87}-1=-47.8\%\\
\end{align*}

\paragraph{300604 长川科技：成长盈利情景 PE}
分母：H1 deducted EPS bridged with H2/H1 conversion；证据质量：高；审计状态：PASS。
为什么选这个方法：电子具有成长或盈利持续期特征，因此模型给盈利桥接配置较高PE区间，而不是直接给主题收入估值；价格和现金流检查防止把泛化成长叙事直接变成EPS。
为什么选这个分母：使用H1扣非利润以减少非经常性收益污染；股本由总市值/当前价反推；H2/H1转换比例是明确的季节性和交付假设，不是机械年化。
为什么用这组倍数：PE区间是House对成长持续期的情景判断；牛市情景要求业绩桥接和经营证据持续，不是市场一致目标价。
为什么用这组概率：30\%/50\%/20\%是统一筛选先验：基准情景权重最高，同时保留熊市和牛市的实质性下行/上行贡献；它是House情景权重，不是经过统计估计的客观事件概率。
证据如何参与：官方半年报预告、Q1财务包和行情/K线提供事实输入；研报/来源用于分母或可比背景。外部目标价380.00元，赋予正权重10\%。
\begin{align*}
\mathrm{Shares} &= \frac{1931.42}{304.44}=6.3442\ \mathrm{100m\ shares}\\
\mathrm{EPS}_{H1}^d &= \frac{9.05}{6.3442}=1.4265\\
\mathrm{EPS}_s &= \mathrm{EPS}_{H1}^d(1+r_s)=2.43/3.31/3.28\\
V_s &= \mathrm{EPS}_s\times \mathrm{PE}_s=43.65/86.06/111.55\\
V_{prob} &= 0.30\times 43.65+0.50\times 86.06+0.20\times 111.55=78.44\\
P_{target} &= 78.44\times 0.90+380.00\times 0.10=108.60\\
\mathrm{Upside} &= \frac{108.60}{304.44}-1=-64.3\%\\
\end{align*}

\paragraph{002916 深南电路：成长盈利情景 PE}
分母：H1 deducted EPS bridged with H2/H1 conversion；证据质量：中；审计状态：PASS。
为什么选这个方法：电子具有成长或盈利持续期特征，因此模型给盈利桥接配置较高PE区间，而不是直接给主题收入估值；价格和现金流检查防止把泛化成长叙事直接变成EPS。
为什么选这个分母：使用H1扣非利润以减少非经常性收益污染；股本由总市值/当前价反推；H2/H1转换比例是明确的季节性和交付假设，不是机械年化。
为什么用这组倍数：PE区间是House对成长持续期的情景判断；牛市情景要求业绩桥接和经营证据持续，不是市场一致目标价。
为什么用这组概率：30\%/50\%/20\%是统一筛选先验：基准情景权重最高，同时保留熊市和牛市的实质性下行/上行贡献；它是House情景权重，不是经过统计估计的客观事件概率。
证据如何参与：官方半年报预告、Q1财务包和行情/K线提供事实输入；研报/来源用于分母或可比背景。没有正权重外部目标锚，House价值不伪装为Street共识。
\begin{align*}
\mathrm{Shares} &= \frac{2569.56}{377.23}=6.8117\ \mathrm{100m\ shares}\\
\mathrm{EPS}_{H1}^d &= \frac{21.80}{6.8117}=3.2004\\
\mathrm{EPS}_s &= \mathrm{EPS}_{H1}^d(1+r_s)=5.44/8.14/7.36\\
V_s &= \mathrm{EPS}_s\times \mathrm{PE}_s=97.93/211.64/250.27\\
V_{prob} &= 0.30\times 97.93+0.50\times 211.64+0.20\times 250.27=185.25\\
P_{target} &= V_{prob}=185.25\quad(\text{外部锚权重}=0)\\
\mathrm{Upside} &= \frac{185.25}{377.23}-1=-50.9\%\\
\end{align*}

\paragraph{688825 长鑫科技：上市边界，不作二级市场估值}
分母：无二级市场当前价；证据质量：低；审计状态：PASS。
为什么选这个方法：截至数据截点尚未开始二级市场交易，直接估算当前价会制造伪精确结果。
为什么选这个分母：不存在可观察的二级市场当前价或交易市值。
为什么用这组倍数：上市前不适用；形成有效市场价格后再建立可比公司/盈利模型。
为什么用这组概率：尚未建立当前价模型，因此不适用。
证据如何参与：仅使用IPO发行价和上市状态证据，不推断Street目标价。
\begin{align*}
P_{\mathrm{current}} &= \mathrm{N/A}\quad(\text{未上市})\\
P_{\mathrm{target}} &= \mathrm{N/A}\quad(\text{上市交易后重建})\\
P_{target} &= V_{prob}=\mathrm{N/A}\quad(\text{外部锚权重}=0)\\
\end{align*}

\paragraph{001248 华润新能源：常规扣非 EPS / PE}
分母：H1 deducted EPS bridged with H2/H1 conversion；证据质量：中高；审计状态：PASS。
为什么选这个方法：该标的H1扣非利润为正且未触发恢复估值阻断，因此扣非EPS是最透明的筛选分母；情景桥接避免把一个H1数字直接当作全年点预测。
为什么选这个分母：使用H1扣非利润以减少非经常性收益污染；股本由总市值/当前价反推；H2/H1转换比例是明确的季节性和交付假设，不是机械年化。
为什么用这组倍数：PE区间按公司所属行业匹配，表达估值压缩/正常/重估三种情景；这是House筛选区间，正式升级前仍需当前可审计的Street区间。
为什么用这组概率：30\%/50\%/20\%是统一筛选先验：基准情景权重最高，同时保留熊市和牛市的实质性下行/上行贡献；它是House情景权重，不是经过统计估计的客观事件概率。
证据如何参与：官方半年报预告、Q1财务包和行情/K线提供事实输入；研报/来源用于分母或可比背景。没有正权重外部目标锚，House价值不伪装为Street共识。
\begin{align*}
\mathrm{Shares} &= \frac{1744.04}{13.41}=130.0552\ \mathrm{100m\ shares}\\
\mathrm{EPS}_{H1}^d &= \frac{35.40}{130.0552}=0.2722\\
\mathrm{EPS}_s &= \mathrm{EPS}_{H1}^d(1+r_s)=0.45/0.52/0.60\\
V_s &= \mathrm{EPS}_s\times \mathrm{PE}_s=4.04/6.72/10.18\\
V_{prob} &= 0.30\times 4.04+0.50\times 6.72+0.20\times 10.18=6.61\\
P_{target} &= V_{prob}=6.61\quad(\text{外部锚权重}=0)\\
\mathrm{Upside} &= \frac{6.61}{13.41}-1=-50.7\%\\
\end{align*}

\paragraph{601633 长城汽车：常规扣非 EPS / PE}
分母：H1 deducted EPS bridged with H2/H1 conversion；证据质量：中高；审计状态：PASS。
为什么选这个方法：该标的H1扣非利润为正且未触发恢复估值阻断，因此扣非EPS是最透明的筛选分母；情景桥接避免把一个H1数字直接当作全年点预测。
为什么选这个分母：使用H1扣非利润以减少非经常性收益污染；股本由总市值/当前价反推；H2/H1转换比例是明确的季节性和交付假设，不是机械年化。
为什么用这组倍数：PE区间按公司所属行业匹配，表达估值压缩/正常/重估三种情景；这是House筛选区间，正式升级前仍需当前可审计的Street区间。
为什么用这组概率：30\%/50\%/20\%是统一筛选先验：基准情景权重最高，同时保留熊市和牛市的实质性下行/上行贡献；它是House情景权重，不是经过统计估计的客观事件概率。
证据如何参与：官方半年报预告、Q1财务包和行情/K线提供事实输入；研报/来源用于分母或可比背景。没有正权重外部目标锚，House价值不伪装为Street共识。
\begin{align*}
\mathrm{Shares} &= \frac{1332.62}{15.58}=85.5340\ \mathrm{100m\ shares}\\
\mathrm{EPS}_{H1}^d &= \frac{16.25}{85.5340}=0.1900\\
\mathrm{EPS}_s &= \mathrm{EPS}_{H1}^d(1+r_s)=0.31/1.06/0.42\\
V_s &= \mathrm{EPS}_s\times \mathrm{PE}_s=3.76/8.36/16.99\\
V_{prob} &= 0.30\times 3.76+0.50\times 8.36+0.20\times 16.99=8.71\\
P_{target} &= V_{prob}=8.71\quad(\text{外部锚权重}=0)\\
\mathrm{Upside} &= \frac{8.71}{15.58}-1=-44.1\%\\
\end{align*}

\paragraph{600736 苏州高新：normalized EPS / PE with PB cross-check}
分母：2026E normalized EPS after property settlement and investment-income normalization；证据质量：中高；审计状态：PASS。
为什么选这个方法：该业务资产较重或受项目周期影响，因此资产价值/BPS和周期调整盈利比直接年化结算利润更稳定；业务匹配可比包括张江高科、陆家嘴、外高桥。
为什么选这个分母：正常化分母为：2026E normalized EPS after property settlement and investment-income normalization；H1机械年化指标只作为市场隐含压力测试，不作为最终目标价分母。
为什么用这组倍数：熊/基准/牛倍数分别表达下行、正常兑现和上行重估；区间通过可比方法和当前价格隐含指标交叉检查（产业园和投资收益提供支撑，但地产结算波动和现金流承压），因此只能作为条件性House价值。
为什么用这组概率：30\%/50\%/20\%是统一筛选先验：基准情景权重最高，同时保留熊市和牛市的实质性下行/上行贡献；它是House情景权重，不是经过统计估计的客观事件概率。
证据如何参与：官方半年报预告、Q1财务包和行情/K线提供事实输入；研报/来源用于分母或可比背景。没有正权重外部目标锚，House价值不伪装为Street共识。
\begin{align*}
V_s=\mathrm{BPS}_{normalized}\times \mathrm{PB}_s\\
V_{b/m/b} &= 4.20/7.70/10.40\ \mathrm{CNY/share}\\
V_{prob} &= 0.30\times 4.20+0.50\times 7.70+0.20\times 10.40=7.19\\
P_{target} &= V_{prob}=7.19\quad(\text{外部锚权重}=0)\\
\mathrm{Upside} &= \frac{7.19}{6.87}-1=4.7\%\\
\end{align*}

\paragraph{000415 渤海租赁：normalized EPS / PE with PB cross-check}
分母：2026E normalized EPS after lease income, asset quality and FX normalization；证据质量：中高；审计状态：PASS。
为什么选这个方法：该业务资产较重或受项目周期影响，因此资产价值/BPS和周期调整盈利比直接年化结算利润更稳定；业务匹配可比包括中银航空租赁、中国飞机租赁、国银租赁。
为什么选这个分母：正常化分母为：2026E normalized EPS after lease income, asset quality and FX normalization；H1机械年化指标只作为市场隐含压力测试，不作为最终目标价分母。
为什么用这组倍数：熊/基准/牛倍数分别表达下行、正常兑现和上行重估；区间通过可比方法和当前价格隐含指标交叉检查（航空租赁景气改善，但资产质量、汇率和杠杆仍压制估值），因此只能作为条件性House价值。
为什么用这组概率：30\%/50\%/20\%是统一筛选先验：基准情景权重最高，同时保留熊市和牛市的实质性下行/上行贡献；它是House情景权重，不是经过统计估计的客观事件概率。
证据如何参与：官方半年报预告、Q1财务包和行情/K线提供事实输入；研报/来源用于分母或可比背景。没有正权重外部目标锚，House价值不伪装为Street共识。
\begin{align*}
V_s=\mathrm{BPS}_{normalized}\times \mathrm{PB}_s\\
V_{b/m/b} &= 3.00/4.62/5.95\ \mathrm{CNY/share}\\
V_{prob} &= 0.30\times 3.00+0.50\times 4.62+0.20\times 5.95=4.40\\
P_{target} &= V_{prob}=4.40\quad(\text{外部锚权重}=0)\\
\mathrm{Upside} &= \frac{4.40}{4.45}-1=-1.1\%\\
\end{align*}

\paragraph{002221 东华能源：周期调整扣非 EPS / PE}
分母：H1 deducted EPS bridged with H2/H1 conversion；证据质量：中；审计状态：PASS。
为什么选这个方法：石油石化盈利受商品价格、运价、价差或利用率周期影响，因此使用周期调整EPS/PE区间，而不是把某一个季度的PE当作全年估值。H2转换比例用于检验当前周期能否延续。
为什么选这个分母：使用H1扣非利润以减少非经常性收益污染；股本由总市值/当前价反推；H2/H1转换比例是明确的季节性和交付假设，不是机械年化。
为什么用这组倍数：熊/基准/牛PE区间是针对该周期业务类型的House假设，不是外部券商目标价；熊值较低反映周期反转，牛值较高必须由价格、销量和现金流共同确认。
为什么用这组概率：30\%/50\%/20\%是统一筛选先验：基准情景权重最高，同时保留熊市和牛市的实质性下行/上行贡献；它是House情景权重，不是经过统计估计的客观事件概率。
证据如何参与：官方半年报预告、Q1财务包和行情/K线提供事实输入；研报/来源用于分母或可比背景。没有正权重外部目标锚，House价值不伪装为Street共识。
\begin{align*}
\mathrm{Shares} &= \frac{85.74}{5.44}=15.7610\ \mathrm{100m\ shares}\\
\mathrm{EPS}_{H1}^d &= \frac{0.95}{15.7610}=0.0603\\
\mathrm{EPS}_s &= \mathrm{EPS}_{H1}^d(1+r_s)=0.09/0.11/0.14\\
V_s &= \mathrm{EPS}_s\times \mathrm{PE}_s=0.56/1.03/1.63\\
V_{prob} &= 0.30\times 0.56+0.50\times 1.03+0.20\times 1.63=1.01\\
P_{target} &= V_{prob}=1.01\quad(\text{外部锚权重}=0)\\
\mathrm{Upside} &= \frac{1.01}{5.44}-1=-81.4\%\\
\end{align*}

\paragraph{002611 东方精工：常规扣非 EPS / PE}
分母：H1 deducted EPS bridged with H2/H1 conversion；证据质量：中；审计状态：PASS。
为什么选这个方法：该标的H1扣非利润为正且未触发恢复估值阻断，因此扣非EPS是最透明的筛选分母；情景桥接避免把一个H1数字直接当作全年点预测。
为什么选这个分母：使用H1扣非利润以减少非经常性收益污染；股本由总市值/当前价反推；H2/H1转换比例是明确的季节性和交付假设，不是机械年化。
为什么用这组倍数：PE区间按公司所属行业匹配，表达估值压缩/正常/重估三种情景；这是House筛选区间，正式升级前仍需当前可审计的Street区间。
为什么用这组概率：30\%/50\%/20\%是统一筛选先验：基准情景权重最高，同时保留熊市和牛市的实质性下行/上行贡献；它是House情景权重，不是经过统计估计的客观事件概率。
证据如何参与：官方半年报预告、Q1财务包和行情/K线提供事实输入；研报/来源用于分母或可比背景。没有正权重外部目标锚，House价值不伪装为Street共识。
\begin{align*}
\mathrm{Shares} &= \frac{182.84}{15.02}=12.1731\ \mathrm{100m\ shares}\\
\mathrm{EPS}_{H1}^d &= \frac{0.90}{12.1731}=0.0739\\
\mathrm{EPS}_s &= \mathrm{EPS}_{H1}^d(1+r_s)=0.12/0.14/0.16\\
V_s &= \mathrm{EPS}_s\times \mathrm{PE}_s=1.71/2.81/4.55\\
V_{prob} &= 0.30\times 1.71+0.50\times 2.81+0.20\times 4.55=2.83\\
P_{target} &= V_{prob}=2.83\quad(\text{外部锚权重}=0)\\
\mathrm{Upside} &= \frac{2.83}{15.02}-1=-81.2\%\\
\end{align*}

\paragraph{002174 游族网络：成长盈利情景 PE}
分母：H1 deducted EPS bridged with H2/H1 conversion；证据质量：中；审计状态：PASS。
为什么选这个方法：传媒具有成长或盈利持续期特征，因此模型给盈利桥接配置较高PE区间，而不是直接给主题收入估值；价格和现金流检查防止把泛化成长叙事直接变成EPS。
为什么选这个分母：使用H1扣非利润以减少非经常性收益污染；股本由总市值/当前价反推；H2/H1转换比例是明确的季节性和交付假设，不是机械年化。
为什么用这组倍数：PE区间是House对成长持续期的情景判断；牛市情景要求业绩桥接和经营证据持续，不是市场一致目标价。
为什么用这组概率：30\%/50\%/20\%是统一筛选先验：基准情景权重最高，同时保留熊市和牛市的实质性下行/上行贡献；它是House情景权重，不是经过统计估计的客观事件概率。
证据如何参与：官方半年报预告、Q1财务包和行情/K线提供事实输入；研报/来源用于分母或可比背景。没有正权重外部目标锚，House价值不伪装为Street共识。
\begin{align*}
\mathrm{Shares} &= \frac{122.87}{12.54}=9.7982\ \mathrm{100m\ shares}\\
\mathrm{EPS}_{H1}^d &= \frac{0.20}{9.7982}=0.0204\\
\mathrm{EPS}_s &= \mathrm{EPS}_{H1}^d(1+r_s)=0.03/0.04/0.05\\
V_s &= \mathrm{EPS}_s\times \mathrm{PE}_s=0.35/0.59/0.88\\
V_{prob} &= 0.30\times 0.35+0.50\times 0.59+0.20\times 0.88=0.58\\
P_{target} &= V_{prob}=0.58\quad(\text{外部锚权重}=0)\\
\mathrm{Upside} &= \frac{0.58}{12.54}-1=-95.4\%\\
\end{align*}

\paragraph{603296 华勤技术：成长盈利情景 PE}
分母：H1 deducted EPS bridged with H2/H1 conversion；证据质量：中高；审计状态：PASS。
为什么选这个方法：电子具有成长或盈利持续期特征，因此模型给盈利桥接配置较高PE区间，而不是直接给主题收入估值；价格和现金流检查防止把泛化成长叙事直接变成EPS。
为什么选这个分母：使用H1扣非利润以减少非经常性收益污染；股本由总市值/当前价反推；H2/H1转换比例是明确的季节性和交付假设，不是机械年化。
为什么用这组倍数：PE区间是House对成长持续期的情景判断；牛市情景要求业绩桥接和经营证据持续，不是市场一致目标价。
为什么用这组概率：30\%/50\%/20\%是统一筛选先验：基准情景权重最高，同时保留熊市和牛市的实质性下行/上行贡献；它是House情景权重，不是经过统计估计的客观事件概率。
证据如何参与：官方半年报预告、Q1财务包和行情/K线提供事实输入；研报/来源用于分母或可比背景。没有正权重外部目标锚，House价值不伪装为Street共识。
\begin{align*}
\mathrm{Shares} &= \frac{1185.67}{78.20}=15.1620\ \mathrm{100m\ shares}\\
\mathrm{EPS}_{H1}^d &= \frac{16.82}{15.1620}=1.1097\\
\mathrm{EPS}_s &= \mathrm{EPS}_{H1}^d(1+r_s)=1.89/4.72/2.55\\
V_s &= \mathrm{EPS}_s\times \mathrm{PE}_s=33.96/86.78/122.72\\
V_{prob} &= 0.30\times 33.96+0.50\times 86.78+0.20\times 122.72=78.12\\
P_{target} &= V_{prob}=78.12\quad(\text{外部锚权重}=0)\\
\mathrm{Upside} &= \frac{78.12}{78.20}-1=-0.1\%\\
\end{align*}

\normalsize
